% Biotechnology — Biology Upper Sixth — Cours signé OnBuch+
% Official MINESEC syllabus. Compiler avec : tectonic BIO07-biotechnology.tex
\newcommand{\DOCMATIERE}{Biology}
\newcommand{\DOCNIVEAU}{Upper Sixth}
\newcommand{\DOCMODULE}{Upper Sixth — Biology}
\newcommand{\DOCLECON}{7}
\newcommand{\DOCTITRE}{Biotechnology}
% OnBuch+ — préambule « cours signés » ENRICHI (compilé avec tectonic / XeLaTeX)
% Système visuel « Pop » d'OnBuch+ : fond crème (jamais blanc), bordures encre
% épaisses, ombres « dures » décalées, titres Archivo Black en capitales.
% Version MATHÉMATIQUES : graphes (pgfplots/tikz), boîtes pédagogiques
% (définition, propriété, méthode, attention, le savais-tu, exercice résolu,
% exercice, TP, bilan), page de garde et signature OnBuch+ ancrée en bas.
%
% Macros attendues dans main.tex AVANT \input{preamble} :
%   \DOCMATIERE  \DOCNIVEAU  \DOCMODULE  \DOCLECON (numéro)  \DOCTITRE
\documentclass[11pt]{article}

\usepackage{fontspec}
\usepackage[english]{babel}
\usepackage[a4paper,margin=2cm,top=3.4cm,bottom=2.8cm,headheight=1.4cm,footskip=1.3cm]{geometry}
\usepackage{amsmath,amssymb,amsfonts}
\usepackage{mathtools} % \xmapsto, \coloneqq... souvent utilisés par les modèles
\usepackage{cancel} % simplifications barrées (bilans de forces)
% \abs{z} (module, valeur absolue) et \norm{u} (norme), utilisés par les modèles
\providecommand{\abs}{}\let\abs\relax\DeclarePairedDelimiter{\abs}{\lvert}{\rvert}
\providecommand{\norm}{}\let\norm\relax\DeclarePairedDelimiter{\norm}{\lVert}{\rVert}
\usepackage[table]{xcolor} % [table] : fournit aussi \rowcolors (alternance de lignes)
\usepackage{pagecolor}
\usepackage{tikz}
\usetikzlibrary{arrows.meta,positioning,calc,shapes.geometric,shapes.misc,shapes.arrows,decorations.pathmorphing,decorations.pathreplacing,decorations.markings,patterns,shadows,mindmap,trees,fit,backgrounds,matrix,intersections,angles,quotes}
\usepackage{pgfplots}
\usepackage[version=4]{mhchem} % équations nucléaires \ce{^{235}_{92}U}
\usepackage[european,siunitx]{circuitikz} % schémas électriques
\pgfplotsset{compat=1.18}
\usepgfplotslibrary{fillbetween,groupplots} % aires entre courbes (calcul intégral)
\usepackage{enumitem}
\usepackage{fancyhdr}
% [most] charge toutes les bibliothèques tcolorbox (listings, minted, raster,
% documentation...) et gonfle inutilement la mémoire TeX sur un document long
% (~300 boîtes) : on ne charge que ce qui est réellement utilisé.
\usepackage[skins,breakable]{tcolorbox}
\usepackage{graphicx}
\usepackage{colortbl}
\usepackage{booktabs}
\usepackage{tabularx}
\usepackage{diagbox} % case d'en-tête diagonale des tableaux à double entrée
% Colonnes extensibles centrée / gauche / droite (C, L, R), souvent utilisées par les modèles
\newcolumntype{C}{>{\centering\arraybackslash}X}
\newcolumntype{L}{>{\raggedright\arraybackslash}X}
\newcolumntype{R}{>{\raggedleft\arraybackslash}X}
\usepackage{array}
\usepackage{multicol}
\usepackage{siunitx}
% parse-numbers=false : accepte aussi \SI{2,5 \times 10^{-3}}{...}
\sisetup{locale=UK,output-decimal-marker={.},group-minimum-digits=4,parse-numbers=false}
\DeclareSIUnit\ohm{\ensuremath{\symup{\Omega}}} % U+2126 absent de la police
\DeclareSIUnit\division{div} % réglages d'oscilloscope (ms/div, V/div)
\DeclareSIUnit\tr{tr} % tours (vitesse de rotation, ex. \SI{1500}{\tr\per\minute})
\usepackage{qrcode}
% \degree : commande fréquemment utilisée par les modèles pour un angle ou une
% température en dehors de \SI{}{} (non fournie par les paquets chargés).
\providecommand{\degree}{^{\circ}}
% \qed (fin de démonstration), utilisé par les modèles sans amsthm
\providecommand{\qed}{\hfill$\square$}
% \barre : double barre des valeurs interdites dans un tableau de variations (tkz-tab)
\providecommand{\barre}{\ensuremath{\|}}
% environnement proof (amsthm non chargé) utilisé par les modèles
\ifdefined\proof\else\newenvironment{proof}[1][Proof]{\par\noindent\textit{#1.}\ }{\qed\par}\fi
\usepackage{lastpage}
\usepackage{hyperref}
\hypersetup{hidelinks}

% ============================================================
% Palette « Pop » OnBuch+ — mêmes hex que la classe P (plus_ui.dart)
% ============================================================
\definecolor{poporange}{HTML}{F59321}
\definecolor{popdark}{HTML}{E07A0C}
\definecolor{popcream}{HTML}{FFF6E8}
\definecolor{popink}{HTML}{1C1714}
\definecolor{poppurple}{HTML}{7A5AE0}
\definecolor{popgreen}{HTML}{2E9E6A}
\definecolor{poppink}{HTML}{E0457B}
\definecolor{popblue}{HTML}{2F7DE1}
\definecolor{popgold}{HTML}{C98A12}
\definecolor{popmuted}{HTML}{8A7F76}
\definecolor{popline}{HTML}{EADFD2}
\definecolor{popgray}{HTML}{8A7F76}
\colorlet{popgrey}{popgray}
% Alias tolérants (le modèle nomme parfois les couleurs différemment)
\colorlet{popwhite}{white}
\colorlet{popblack}{popink}
\colorlet{popred}{poppink}
\colorlet{popyellow}{popgold}
% Teintes claires pour fonds de tableaux / zones
\colorlet{popblueL}{popblue!10!white}
\colorlet{poporangeL}{poporange!14!white}
\colorlet{popgreenL}{popgreen!12!white}
\colorlet{poppurpleL}{poppurple!12!white}
\colorlet{poppinkL}{poppink!12!white}
\colorlet{popgoldL}{popgold!14!white}

\pagecolor{popcream}

% ============================================================
% Typographie
% ============================================================
\defaultfontfeatures{Ligatures=TeX}
\setmainfont{PlusJakartaSans-Medium.ttf}[Path=../../../fonts/,BoldFont=PlusJakartaSans-Bold.ttf]
% Maths en sans-serif (Fira Math) avec chiffres et lettres droites de Plus
% Jakarta, pour que les formules se fondent dans le texte.
\usepackage[math-style=ISO,bold-style=ISO]{unicode-math}
\setmathfont{FiraMath-Regular.otf}[Path=../../../fonts/]
\setmathfont{PlusJakartaSans-Medium.ttf}[Path=../../../fonts/,range={up/num,up/{Latin,latin},\mathup}]
% (icomma retiré : décimales avec point en anglais)
% Caractères mathématiques Unicode absents des polices de texte (Plus Jakarta,
% Archivo Black) : rendus par leur équivalent mathématique, en texte comme en
% formule — sinon ils s'affichent en carré vide (ex. « Arithmétique dans ℤ »).
\usepackage{newunicodechar}
\newunicodechar{ℝ}{\ensuremath{\mathbb{R}}}
\newunicodechar{ℤ}{\ensuremath{\mathbb{Z}}}
\newunicodechar{ℂ}{\ensuremath{\mathbb{C}}}
\newunicodechar{ℕ}{\ensuremath{\mathbb{N}}}
\newunicodechar{ℚ}{\ensuremath{\mathbb{Q}}}
\newunicodechar{𝔻}{\ensuremath{\mathbb{D}}}
\newunicodechar{α}{\ensuremath{\alpha}}
\newunicodechar{β}{\ensuremath{\beta}}
\newunicodechar{γ}{\ensuremath{\gamma}}
\newunicodechar{δ}{\ensuremath{\delta}}
\newunicodechar{ε}{\ensuremath{\varepsilon}}
\newunicodechar{θ}{\ensuremath{\theta}}
\newunicodechar{λ}{\ensuremath{\lambda}}
\newunicodechar{μ}{\ensuremath{\mu}}
\newunicodechar{σ}{\ensuremath{\sigma}}
\newunicodechar{τ}{\ensuremath{\tau}}
\newunicodechar{φ}{\ensuremath{\varphi}}
\newunicodechar{ψ}{\ensuremath{\psi}}
\newunicodechar{ω}{\ensuremath{\omega}}
\newunicodechar{Σ}{\ensuremath{\Sigma}}
% majuscules produites par \MakeUppercase dans les titres (α-aminés -> Α)
\newunicodechar{Α}{\ensuremath{\alpha}}
\newunicodechar{Β}{\ensuremath{\beta}}
\newunicodechar{Γ}{\ensuremath{\Gamma}}
\newunicodechar{Δ}{\ensuremath{\Delta}}
\newunicodechar{Ε}{\ensuremath{\varepsilon}}
\newunicodechar{Θ}{\ensuremath{\Theta}}
\newunicodechar{Λ}{\ensuremath{\Lambda}}
\newunicodechar{Μ}{\ensuremath{\mu}}
\newunicodechar{Τ}{\ensuremath{\tau}}
\newunicodechar{Φ}{\ensuremath{\Phi}}
\newunicodechar{Ψ}{\ensuremath{\Psi}}
\newunicodechar{Ω}{\ensuremath{\Omega}}
\newunicodechar{↦}{\ensuremath{\mapsto}}
\newunicodechar{∈}{\ensuremath{\in}}
\newunicodechar{∉}{\ensuremath{\notin}}
\newunicodechar{∧}{\ensuremath{\wedge}}
\newunicodechar{⇔}{\ensuremath{\Leftrightarrow}}
\newunicodechar{⟺}{\ensuremath{\Longleftrightarrow}}
\newunicodechar{⇒}{\ensuremath{\Rightarrow}}
\newunicodechar{⟹}{\ensuremath{\Longrightarrow}}
\newunicodechar{∘}{\ensuremath{\circ}}
\newunicodechar{∪}{\ensuremath{\cup}}
\newunicodechar{∩}{\ensuremath{\cap}}
\newunicodechar{≡}{\ensuremath{\equiv}}
\newunicodechar{⊂}{\ensuremath{\subset}}
\newunicodechar{⊥}{\ensuremath{\perp}}
\newunicodechar{∃}{\ensuremath{\exists}}
\newunicodechar{∀}{\ensuremath{\forall}}
\newunicodechar{∛}{\ensuremath{\sqrt[3]{\,}}}
\newunicodechar{∜}{\ensuremath{\sqrt[4]{\,}}}
\newunicodechar{⃗}{\ensuremath{{}^{\to}}}
\newunicodechar{ⁿ}{\ifmmode^{n}\else\textsuperscript{\ensuremath{n}}\fi}
\newunicodechar{ˣ}{\ifmmode^{x}\else\textsuperscript{\ensuremath{x}}\fi}
\newunicodechar{ᵇ}{\ifmmode^{b}\else\textsuperscript{\ensuremath{b}}\fi}
\newunicodechar{ʳ}{\ifmmode^{r}\else\textsuperscript{\ensuremath{r}}\fi}
\newunicodechar{ᵘ}{\ifmmode^{u}\else\textsuperscript{\ensuremath{u}}\fi}
\newunicodechar{ʸ}{\ifmmode^{y}\else\textsuperscript{\ensuremath{y}}\fi}
\newunicodechar{ᵃ}{\ifmmode^{a}\else\textsuperscript{\ensuremath{a}}\fi}
\newunicodechar{ᵉ}{\ifmmode^{e}\else\textsuperscript{\ensuremath{e}}\fi}
\newunicodechar{ᵏ}{\ifmmode^{k}\else\textsuperscript{\ensuremath{k}}\fi}
\newunicodechar{ᵖ}{\ifmmode^{p}\else\textsuperscript{\ensuremath{p}}\fi}
\newunicodechar{ᴺ}{\ifmmode^{N}\else\textsuperscript{\ensuremath{N}}\fi}
\newunicodechar{ᶜ}{\ifmmode^{c}\else\textsuperscript{\ensuremath{c}}\fi}
\newunicodechar{ʰ}{\ifmmode^{h}\else\textsuperscript{\ensuremath{h}}\fi}
\newunicodechar{ᵠ}{\ifmmode^{q}\else\textsuperscript{\ensuremath{q}}\fi}
\newunicodechar{ₙ}{\ifmmode_{n}\else\textsubscript{\ensuremath{n}}\fi}
\newunicodechar{ₐ}{\ifmmode_{a}\else\textsubscript{\ensuremath{a}}\fi}
\newunicodechar{ᵢ}{\ifmmode_{i}\else\textsubscript{\ensuremath{i}}\fi}
\newunicodechar{ᵣ}{\ifmmode_{r}\else\textsubscript{\ensuremath{r}}\fi}
\newunicodechar{ₖ}{\ifmmode_{k}\else\textsubscript{\ensuremath{k}}\fi}
\newunicodechar{ᵦ}{\ifmmode_{\beta}\else\textsubscript{\ensuremath{\beta}}\fi}
\newunicodechar{ₓ}{\ifmmode_{x}\else\textsubscript{\ensuremath{x}}\fi}
\newfontfamily\popdisplay{ArchivoBlack-Regular.ttf}[Path=../../../fonts/]
\newfontfamily\poplogo{SpaceGrotesk-Bold.ttf}[Path=../../../fonts/]
\newfontfamily\popsemi{PlusJakartaSans-SemiBold.ttf}[Path=../../../fonts/]
\newfontfamily\popxbold{PlusJakartaSans-ExtraBold.ttf}[Path=../../../fonts/]
\newfontfamily\popxboldls{PlusJakartaSans-ExtraBold.ttf}[Path=../../../fonts/,LetterSpace=3.0]

\setlist[enumerate]{leftmargin=1.7em,itemsep=5pt,topsep=4pt}
\setlist[itemize]{leftmargin=1.7em,itemsep=4pt,topsep=4pt,label={\color{poporange}\textbullet}}
\setlength{\parindent}{0pt}
\setlength{\parskip}{5pt}
\renewcommand{\arraystretch}{1.25}

% pgfplots : style Pop par défaut
\pgfplotsset{
  every axis/.append style={axis line style={popink,line width=1pt},tick style={popink},
    grid style={popline,line width=0.5pt},label style={font=\small},tick label style={font=\footnotesize}},
  popcurve/.style={poporange,line width=1.8pt,no markers,smooth},
}
% Même style disponible dans un \draw TikZ simple (hors pgfplots)
\tikzset{popcurve/.style={poporange,line width=1.8pt}}
\tikzset{popgrid/.style={popline,very thin}}
\colorlet{popcurve}{poporange} % le nom du style est parfois utilisé comme couleur

% ============================================================
% Pill / squiggle
% ============================================================
\newcommand{\poppill}[3]{%
  \tikz[baseline=-0.55ex]\node[draw=popink,line width=1.2pt,rounded corners=2.6pt,%
    fill=#2,inner xsep=6.5pt,inner ysep=3pt]{%
    \popxboldls\fontsize{7}{7}\selectfont\color{#3}\MakeUppercase{#1}};%
}
\newcommand{\popsquiggle}[1]{%
  \tikz[baseline=-0.4ex]{
    \draw[#1,line width=2.6pt,line cap=round]
      plot [smooth] coordinates {(0,0.09) (0.34,0.01) (0.68,0.21) (1.02,0.01) (1.36,0.10)};
  }%
}
% Mot-clé surligné (marqueur orange sous le texte)
\newcommand{\cle}[1]{{\popxbold\color{popdark}#1}}

% ============================================================
% En-tête / pied
% ============================================================
\pagestyle{fancy}
\fancyhf{}
\renewcommand{\headrulewidth}{0pt}
\renewcommand{\footrulewidth}{0pt}
\lhead{\raisebox{-0.15ex}{\poplogo\fontsize{13}{13}\selectfont\color{popink}OnBuch\color{poporange}+}}
\rhead{\poppill{\DOCMATIERE\ \textbullet\ \DOCNIVEAU\ \textbullet\ Chapter \DOCLECON}{white}{popink}}
\lfoot{\popxbold\fontsize{7.6}{7.6}\selectfont\color{popmuted}\MakeUppercase{Signed course OnBuch+}\ \textbullet\ {\popsemi\DOCTITRE}}
\rfoot{\tikz[baseline=-0.6ex]\node[rounded corners=8.5pt,fill=popink,minimum height=17pt,minimum width=17pt,inner xsep=5pt,inner sep=0]{\popxbold\fontsize{8}{8}\selectfont\color{popcream}\thepage\,/\,\pageref*{LastPage}};}

% ============================================================
% Titres
% ============================================================
\newcommand{\chaptitre}[2]{%
  \begin{tcolorbox}[enhanced,colback=poporange,colframe=popink,boxrule=2.6pt,arc=11pt,
      drop shadow=popink,left=15pt,right=15pt,top=12pt,bottom=11pt,before skip=3pt,after skip=18pt]
    \poppill{#2}{popink}{popcream}\par\vspace{9pt}
    {\raggedright\hyphenpenalty=10000\exhyphenpenalty=10000\popdisplay\fontsize{22}{24}\selectfont\color{popcream}\MakeUppercase{#1}\par}
  \end{tcolorbox}
}
% Section numérotée : pastille orange avec le numéro + titre Archivo
\newcounter{popsec}
\newcommand{\coursec}[1]{%
  \stepcounter{popsec}\setcounter{popsub}{0}%
  \par\vspace{16pt}\noindent
  \begin{minipage}{\linewidth}
  \tikz[baseline=-0.7ex]\node[draw=popink,line width=1.6pt,fill=poporange,rounded corners=4pt,minimum size=22pt,inner sep=0]{\popdisplay\fontsize{12}{12}\selectfont\color{popcream}\thepopsec};%
  \hspace{8pt}{\popdisplay\fontsize{15}{17}\selectfont\color{popink}\MakeUppercase{#1}}\par
  \vspace{4pt}\noindent{\color{poporange}\rule{36pt}{3.6pt}}
  \end{minipage}\par\nopagebreak\vspace{8pt}
}
\newcounter{popsub}
\newcommand{\courssub}[1]{%
  \stepcounter{popsub}%
  \par\vspace{9pt}\noindent{\popxbold\fontsize{11.5}{13}\selectfont\color{popink}{\color{poporange}\thepopsec.\thepopsub}\hspace{6pt}#1}\par\nopagebreak\vspace{4pt}
}

% ============================================================
% Boîtes pédagogiques — même recette : fond blanc, bordure épaisse colorée,
% ombre dure encre, badge plein. Titre optionnel [#1].
% ============================================================
\tcbset{popbase/.style={breakable,enhanced,colback=white,boxrule=2.3pt,arc=9pt,
  shadow={3pt}{-3pt}{0pt}{popink},left=14pt,right=14pt,top=11pt,bottom=11pt,before skip=11pt,after skip=15pt,
  fontupper=\popsemi\fontsize{10.6}{14.5}\selectfont\color{popink},
  fontlower=\popsemi\fontsize{10.6}{14.5}\selectfont\color{popink}}}
% Coche dessinée (le glyphe ✓ n'existe pas dans les polices du document)
\newcommand{\popcheck}[1]{\tikz[baseline=-0.5ex]\draw[#1,line width=1.6pt,line cap=round,line join=round](-0.13,0.0)--(-0.04,-0.09)--(0.13,0.1);}
\providecommand{\coche}{\popcheck{popgreen}}
\providecommand{\resultat}[1]{\par\smallskip\noindent\fcolorbox{popgreen}{popgreenL}{\parbox{\dimexpr\linewidth-2\fboxsep-2\fboxrule\relax}{\textbf{Result.} #1}}\par\smallskip}
\newcommand{\popboxhead}[3]{\poppill{#1}{#2}{white}\ifx\relax#3\relax\else\hspace{6pt}{\popxbold\fontsize{10.6}{12}\selectfont\color{popink}#3}\fi\par\vspace{7pt}}

\newtcolorbox{aretenir}[1][]{popbase,colframe=popgreen,before upper={\popboxhead{Key points}{popgreen}{#1}}}
\newtcolorbox{exemplebox}[1][]{popbase,colframe=poporange,before upper={\popboxhead{Exemple}{poporange}{#1}}}
\newtcolorbox{definition}[1][]{popbase,colframe=popblue,before upper={\popboxhead{Definition}{popblue}{#1}}}
\newtcolorbox{propriete}[1][]{popbase,colframe=poppurple,before upper={\popboxhead{Property}{poppurple}{#1}}}
\newtcolorbox{methode}[1][]{popbase,colframe=popgold,before upper={\popboxhead{Method}{popgold}{#1}}}
\newtcolorbox{attention}[1][]{popbase,colframe=poppink,before upper={\popboxhead{Warning}{poppink}{#1}}}
\newtcolorbox{experience}[1][]{popbase,colframe=popblue,colback=popblueL,before upper={\popboxhead{Experiment}{popblue}{#1}}}
% « Le savais-tu ? » : carte violette pleine (contexte, culture, Cameroun)
\newtcolorbox{savaistu}[1][]{popbase,colback=poppurple,colframe=popink,
  fontupper=\popsemi\fontsize{10.4}{14.2}\selectfont\color{white},
  before upper={\poppill{Did you know?}{white}{poppurple}\ifx\relax#1\relax\else\hspace{6pt}{\popxbold\color{white}#1}\fi\par\vspace{7pt}}}
% Exercice résolu : énoncé en haut, \tcblower puis la solution sur fond crème
\newcounter{popexo}\newcounter{popexr}
\newtcolorbox{exoresolu}[1][]{popbase,colframe=poporange,colbacklower=popcream,
  segmentation style={popink,line width=1.2pt,dashed},
  before upper={\stepcounter{popexr}\popboxhead{Worked example \thepopexr}{poporange}{#1}},
  before lower={\poppill{Solution}{popink}{popcream}\par\vspace{6pt}}}
% Exercice (non résolu en place) — la correction est dans la partie Corrigés
\newtcolorbox{exercice}[1][]{popbase,colframe=popink,
  before upper={\stepcounter{popexo}\popboxhead{Exercise \thepopexo}{popink}{#1}}}
% Corrigé
\newtcolorbox{corrige}[1][]{popbase,colframe=popgreen,colback=popgreenL,
  before upper={\popboxhead{Solution}{popgreen}{#1}}}
% Objectifs / prérequis / bilan
\newtcolorbox{objectifs}{popbase,colframe=popink,colback=poporangeL,
  before upper={\popboxhead{Chapter objectives}{popink}{}}}
\newtcolorbox{prerequis}{popbase,colframe=popink,colback=popblueL,
  before upper={\popboxhead{Prerequisites}{popblue}{}}}
\newtcolorbox{bilan}{popbase,colframe=popink,colback=white,boxrule=2.8pt,
  before upper={\popboxhead{Fiche bilan}{poporange}{L'essentiel à connaître pour l'examen}}}
% Figure encadrée avec légende
\newtcolorbox{popfigure}[1][]{enhanced,breakable=false,colback=white,colframe=popline,boxrule=1.4pt,arc=8pt,
  left=10pt,right=10pt,top=10pt,bottom=8pt,before skip=10pt,after skip=12pt,halign=center,
  lower separated=false,halign lower=center,
  fontlower=\popsemi\fontsize{9}{11.5}\selectfont\color{popmuted},
  #1}
\newcommand{\legende}[1]{\par\vspace{6pt}{\popsemi\fontsize{9}{11.5}\selectfont\color{popmuted}#1}}

% ============================================================
% Signature OnBuch+
% ============================================================
\newcommand{\poplogoinline}[1]{{\poplogo\fontsize{#1}{#1}\selectfont\color{popink}OnBuch\color{poporange}+}}
\newcommand{\signatureonbuch}{%
  \begin{tcolorbox}[enhanced,colback=popink,colframe=popink,boxrule=2.2pt,arc=10pt,
      drop shadow=poporange,left=14pt,right=14pt,top=10pt,bottom=10pt,before skip=0pt,after skip=0pt]
    \tikz[baseline=-0.7ex]\node[circle,fill=popgreen,draw=popcream,line width=1.3pt,minimum size=22pt,inner sep=0]{\popcheck{white}};%
    \hspace{9pt}\begin{minipage}[c]{0.62\linewidth}
      {\popxbold\fontsize{12}{13}\selectfont\color{popcream}Signed\ \textbullet\ {\poplogo OnBuch\color{poporange}+}}\par\vspace{2pt}
      {\raggedright\popsemi\fontsize{8}{10.5}\selectfont\color{popline}\MakeUppercase{OnBuch+ teaching team}\\\MakeUppercase{Follows the official MINESEC syllabus}\par}
    \end{minipage}\hfill
    {\poplogo\fontsize{22}{22}\selectfont\color{popcream}OnBuch\color{poporange}+}
  \end{tcolorbox}%
}

% ============================================================
% Page de garde
% #1 titre · #2 matière · #3 niveau · #4 module · #5 n° leçon · #6 durée ·
% #7 description / objectifs (texte ou itemize)
% ============================================================
\newcommand{\pagegarde}[7]{%
  \thispagestyle{empty}
  \newgeometry{a4paper,margin=1.7cm,top=1.6cm,bottom=1.4cm}%
  \begin{tikzpicture}[remember picture,overlay]
    \fill[poporange!18] ([xshift=-3.2cm,yshift=-3.6cm]current page.north east) circle (4.6cm);
    \fill[poppurple!14] ([xshift=2.4cm,yshift=7.6cm]current page.south west) circle (3.8cm);
  \end{tikzpicture}%
  {\poplogo\fontsize{30}{30}\selectfont\color{popink}OnBuch\color{poporange}+}\hfill
  \raisebox{0.6ex}{\poppill{Signed course \textbullet\ Premium}{popink}{popcream}}\par
  \vspace{1pt}\popsquiggle{poppurple}\par
  \vspace{8pt}
  \begin{tcolorbox}[enhanced,colback=poporange,colframe=popink,boxrule=2.8pt,arc=13pt,
      drop shadow=popink,left=18pt,right=18pt,top=12pt,bottom=13pt,after skip=10pt]
    \poppill{#2 \textbullet\ #3}{popink}{popcream}\hspace{5pt}\poppill{#4}{white}{popink}\par\vspace{8pt}
    {\popdisplay\fontsize{14}{14}\selectfont\color{popink}CHAPTER #5}\par\vspace{4pt}
    {\raggedright\hyphenpenalty=10000\exhyphenpenalty=10000\popdisplay\fontsize{25}{27}\selectfont\color{popcream}\MakeUppercase{#1}\par}
    \vspace{8pt}
    {\popxbold\fontsize{9.5}{9.5}\selectfont\color{popink}\MakeUppercase{Suggested time: #6}}
  \end{tcolorbox}
  \begin{tcolorbox}[enhanced,colback=white,colframe=popink,boxrule=2.2pt,arc=10pt,
      drop shadow=popink,left=16pt,right=16pt,top=10pt,bottom=10pt,after skip=8pt]
    \poppill{In this chapter}{poppurple}{white}\par\vspace{5pt}
    \popsemi\fontsize{9.3}{12.6}\selectfont\color{popink}#7
  \end{tcolorbox}
  \noindent\begin{minipage}[t]{0.487\linewidth}
    \begin{tcolorbox}[enhanced,colback=popgreen,colframe=popink,boxrule=2.2pt,arc=10pt,
        drop shadow=popink,left=11pt,right=11pt,top=10pt,bottom=10pt,height=3.1cm,valign=center]
      \tcbox[colback=white,colframe=popink,boxrule=1.3pt,arc=5pt,boxsep=0pt,nobeforeafter,
        left=3pt,right=3pt,top=3pt,bottom=3pt,valign=center]{\qrcode[height=1.9cm]{https://play.google.com/store/apps/details?id=com.luvvix.onbuch&hl=en}}%
      \hspace{8pt}\begin{minipage}[c]{0.5\linewidth}
        {\popdisplay\fontsize{11}{12.5}\selectfont\color{white}\MakeUppercase{Download OnBuch}}\par\vspace{4pt}
        {\raggedright\popsemi\fontsize{8.4}{11}\selectfont\color{white}Free \textbullet\ Google Play\\AI tutor, past papers, results\par}
      \end{minipage}
    \end{tcolorbox}
  \end{minipage}\hfill
  \begin{minipage}[t]{0.487\linewidth}
    \begin{tcolorbox}[enhanced,colback=poppurple,colframe=popink,boxrule=2.2pt,arc=10pt,
        drop shadow=popink,left=11pt,right=11pt,top=10pt,bottom=10pt,height=3.1cm,valign=center]
      \tikz[baseline=-0.7ex]\node[circle,fill=white,draw=popink,line width=1.3pt,minimum size=1.6cm,inner sep=0]{\popdisplay\fontsize{24}{24}\selectfont\color{poppurple}+};%
      \hspace{8pt}\begin{minipage}[c]{0.6\linewidth}
        {\popdisplay\fontsize{11}{12.5}\selectfont\color{white}\MakeUppercase{Go OnBuch+}}\par\vspace{4pt}
        {\raggedright\popsemi\fontsize{8.4}{11}\selectfont\color{white}All signed courses, unlimited AI tutor, PDF sheets — OnBuch+ tab in the app\par}
      \end{minipage}
    \end{tcolorbox}
  \end{minipage}
  \vspace{8pt}
  \popcontenu
  \vfill
  \signatureonbuch
  \restoregeometry
  \newpage
}

% Rangée « Ce cours contient » (page de garde)
\newcommand{\poptile}[3]{% #1 couleur #2 gros texte #3 légende
  \begin{tcolorbox}[enhanced,colback=#1,colframe=popink,boxrule=2pt,arc=9pt,shadow={2.5pt}{-2.5pt}{0pt}{popink},
      left=6pt,right=6pt,top=9pt,bottom=9pt,halign=center,valign=center,height=1.75cm,width=\linewidth,nobeforeafter]
    {\popdisplay\fontsize{12.5}{13}\selectfont\color{white}\MakeUppercase{#2}}\par\vspace{3pt}
    {\centering\popsemi\fontsize{7.8}{9.5}\selectfont\color{white}#3\par}
  \end{tcolorbox}}
\newcommand{\popcontenu}{%
  \par\noindent{\popxboldls\fontsize{7.5}{7.5}\selectfont\color{popmuted}THIS SIGNED COURSE CONTAINS}\par\vspace{7pt}
  \noindent\begin{minipage}[t]{0.235\linewidth}\poptile{popblue}{Course}{complete, illustrated, step by step}\end{minipage}\hfill
  \begin{minipage}[t]{0.235\linewidth}\poptile{poporange}{Methods}{and worked examples}\end{minipage}\hfill
  \begin{minipage}[t]{0.235\linewidth}\poptile{poppink}{Exercises}{exam style + solutions}\end{minipage}\hfill
  \begin{minipage}[t]{0.235\linewidth}\poptile{popgreen}{Summary sheet}{the essentials to remember}\end{minipage}\par}

% Sommaire
\newtcolorbox{sommaire}{enhanced,colback=white,colframe=popink,boxrule=2.2pt,arc=10pt,
  drop shadow=popink,left=18pt,right=18pt,top=13pt,bottom=13pt,before skip=6pt,after skip=20pt,
  before upper={\poppill{Contents}{poppurple}{white}\par\vspace{9pt}\popsemi\fontsize{10.6}{14.5}\selectfont\color{popink}}}

% Signature de fin de cours — poussée tout en bas de la dernière page
\newcommand{\findecours}{%
  \par\vfill\nopagebreak
  \begin{center}\popsquiggle{poporange}\end{center}\vspace{4pt}
  \signatureonbuch
}
\AtBeginDocument{\def\llbracket{\mathopen{[\mkern-3.5mu[}}\def\rrbracket{\mathclose{]\mkern-3.5mu]}}} % intervalles d'entiers (glyphe ⟦ absent de la police)
% Glyphes absents de Fira Math : redéfinis à partir de glyphes disponibles
\AtBeginDocument{%
  \def\setminus{\mathbin{\backslash}}\let\smallsetminus\setminus
  \def\vdots{\mathord{\vbox{\baselineskip3.2pt\lineskiplimit0pt\kern4pt\hbox{.}\hbox{.}\hbox{.}}}}%
  \def\ddots{\mathinner{\mkern1mu\raise6.5pt\hbox{.}\mkern2mu\raise3.6pt\hbox{.}\mkern2mu\raise0.7pt\hbox{.}\mkern1mu}}%
  \def\triangle{\mathord{\scalebox{0.9}{$\symup{\Delta}$}}}%
  \def\nabla{\mathord{\raisebox{\depth}{\scalebox{1}[-1]{$\symup{\Delta}$}}}}%
  \def\checkmark{\popcheck{popdark}}%
  \def\star{\mathbin{\ast}}\def\bigstar{\mathord{\ast}}%
  \def\diamond{\mathbin{\scalebox{0.6}{$\Diamond$}}}%
  \def\bigcup{\mathop{\mathchoice{\raisebox{-0.3em}{\scalebox{1.7}{$\cup$}}}{\scalebox{1.25}{$\cup$}}{\cup}{\cup}}\displaylimits}%
  \def\bigcap{\mathop{\mathchoice{\raisebox{-0.3em}{\scalebox{1.7}{$\cap$}}}{\scalebox{1.25}{$\cap$}}{\cap}{\cap}}\displaylimits}%
  \def\lesseqqgtr{\mathrel{\lessgtr}}\def\bigoplus{\mathop{\oplus}\displaylimits}%
}
\providecommand{\ssi}{if and only if} % abréviation fréquente
\AtBeginDocument{\providecommand{\wideparen}{\overparen}} % arc de cercle

%% FIGFIX-BEGIN v5 — correctifs globaux des figures (docs/onbuch-generation/tools/figfix)
\usepackage{adjustbox}\usepackage{etoolbox}\tcbuselibrary{raster}
% toute figure TikZ/pgfplots plus large que la ligne est réduite à la largeur disponible
\BeforeBeginEnvironment{tikzpicture}{\begin{adjustbox}{max width=\linewidth}}
\AfterEndEnvironment{tikzpicture}{\end{adjustbox}}
% légendes pgfplots lisibles par-dessus les courbes
\pgfplotsset{every axis/.append style={legend style={fill=white,fill opacity=0.92,text opacity=1,draw=black!15,inner sep=3pt}}}
% étiquettes d'arêtes (syntaxe edge["..."]) avec fond blanc
\tikzset{every edge quotes/.append style={fill=white,inner sep=1.2pt}}
%% FIGFIX-END
\begin{document}

\pagegarde{Biotechnology}{Biology}{Upper Sixth}{Upper Sixth — Biology}{7}{6 periods}{Ce chapitre couvre l'ensemble du programme de biotechnologie pour la Terminale : des procédés traditionnels de fermentation aux technologies de l'ADN recombinant, en passant par la thérapie génique et les applications médicales, agricoles et environnementales. Il prépare l'élève à analyser, évaluer et concevoir des applications biotechnologiques pertinentes pour le contexte camerounais et mondial.
\vspace{4pt}
\begin{multicols}{2}\small
\begin{itemize}
\item Expliquer les principes biochimiques des fermentations traditionnelles et identifier les micro-organismes impliqués.
\item Décrire les mécanismes d'homéostasie et illustrer comment la biotechnologie intervient dans le soutien et la locomotion (prothèses, biomatériaux).
\item Analyser les stratégies de correction génique (thérapie génique, édition CRISPR) et en évaluer les avantages et limites.
\item Maîtriser les étapes clés de la technologie de l'ADN recombinant : enzymes de restriction, vecteurs, clonage, sélection et amplification par PCR.
\item Évaluer les applications majeures de la biotechnologie (médicaments, OGM, bioremédiation, biocarburants) sous l'angle scientifique, éthique et réglementaire.
\item Concevoir un protocole simple de clonage ou de détection d'un gène d'intérêt en utilisant les outils moléculaires standards.
\end{itemize}\end{multicols}}

\begin{sommaire}
\begin{multicols}{2}
\begin{enumerate}
\item 1. Biotechnologie traditionnelle : fermentation et transformation alimentaire
\item 2. Biotechnologie au service de l'homéostasie, du soutien et de la locomotion
\item 3. Correction et remédiation génique : thérapie génique et édition du génome
\item 4. Technologie de l'ADN recombinant et altération génique
\item 5. Applications majeures de la biotechnologie : santé, agriculture, environnement, industrie
\item Integration activity
\item Methods and worked examples
\item Exercises
\item Solutions to the exercises
\item Summary sheet
\end{enumerate}
\end{multicols}
\end{sommaire}

\begin{objectifs}
\begin{itemize}
\item Expliquer les principes biochimiques des fermentations traditionnelles et identifier les micro-organismes impliqués.
\item Décrire les mécanismes d'homéostasie et illustrer comment la biotechnologie intervient dans le soutien et la locomotion (prothèses, biomatériaux).
\item Analyser les stratégies de correction génique (thérapie génique, édition CRISPR) et en évaluer les avantages et limites.
\item Maîtriser les étapes clés de la technologie de l'ADN recombinant : enzymes de restriction, vecteurs, clonage, sélection et amplification par PCR.
\item Évaluer les applications majeures de la biotechnologie (médicaments, OGM, bioremédiation, biocarburants) sous l'angle scientifique, éthique et réglementaire.
\item Concevoir un protocole simple de clonage ou de détection d'un gène d'intérêt en utilisant les outils moléculaires standards.
\item Argumenter sur les enjeux de biosécurité et de propriété intellectuelle liés aux biotechnologies modernes.
\end{itemize}
\end{objectifs}

\begin{prerequis}
\begin{itemize}
\item Structure et fonction des cellules procaryotes et eucaryotes (Lower Sixth).
\item Bases de la génétique mendélienne et moléculaire (réplication, transcription, traduction).
\item Cinétique enzymatique et rôle des enzymes dans le métabolisme.
\item Notions de microbiologie : culture, identification, croissance bactérienne.
\item Principes de base de la chimie organique (glucides, acides aminés, acides nucléiques).
\end{itemize}
\end{prerequis}

\newpage

\coursec{Biotechnologie traditionnelle : fermentation et transformation alimentaire}

Dans les marchés de Yaoundé, l'odeur du vin de palme et du foutou fermenté rappelle que la biotechnologie fait partie du quotidien camerounais depuis des siècles. Aujourd'hui, les mêmes principes permettent de produire de l'insuline humaine dans des bactéries et de créer des plantes résistantes à la sécheresse. Comment passer de la fermentation artisanale à l'édition du génome ? Ce chapitre explore cette évolution, des méthodes traditionnelles aux outils moléculaires de pointe.

Le problème central de cette première section est le suivant : \emph{comment des micro-organismes invisibles transforment-ils nos aliments, et comment l'Homme a-t-il appris, d'abord empiriquement puis scientifiquement, à contrôler ces transformations ?}

\courssub{Définitions : biotechnologie traditionnelle et moderne}

Observez un vendeur de vin de palme à Buea : il récolte la sève du palmier, la laisse reposer quelques heures, et le liquide sucré devient pétillant et alcoolisé. Personne n'a ajouté de levure : les micro-organismes présents naturellement dans l'air et sur l'écorce ont fait le travail. C'est cela, la biotechnologie traditionnelle : une utilisation \textbf{empirique} du vivant, transmise de génération en génération, sans connaissance des mécanismes moléculaires sous-jacents.

\begin{definition}[Biotechnologie]
La \cle{biotechnologie} est l'utilisation d'\textbf{organismes vivants} (micro-organismes, cellules végétales ou animales) ou de leurs \textbf{composants} (enzymes, gènes) pour produire des biens et des services utiles à l'Homme : aliments, boissons, médicaments, procédés industriels.
\end{definition}

\begin{definition}[Biotechnologie traditionnelle vs moderne]
\begin{itemize}
\item La \cle{biotechnologie traditionnelle} repose sur des pratiques empiriques ancestrales : fermentation spontanée des aliments et boissons, sélection paysanne des semences, fabrication artisanale de fromages. Elle ne modifie pas directement le patrimoine génétique des organismes.
\item La \cle{biotechnologie moderne} utilise les outils de la biologie moléculaire : ADN recombinant, génie génétique, culture cellulaire, édition du génome. Elle permet de modifier de façon ciblée le génome des organismes (insuline humaine produite par \emph{Escherichia coli}, plantes transgéniques).
\end{itemize}
\end{definition}

\begin{aretenir}[Lien entre les deux]
La biotechnologie moderne n'est pas une rupture mais un \textbf{approfondissement} : elle repose sur les mêmes phénomènes biologiques (métabolisme microbien, reproduction, expression des gènes) que la tradition a exploités aveuglément pendant des millénaires.
\end{aretenir}

\courssub{La fermentation : une voie métabolique anaérobie}

Pourquoi la sève de palme devient-elle alcoolisée ? Parce que les levures, en l'absence d'oxygène, doivent tout de même produire de l'ATP pour survivre. Elles utilisent alors une voie de secours : la fermentation.

\begin{definition}[Fermentation]
La \cle{fermentation} est le \textbf{catabolisme anaérobie incomplet} d'un substrat organique (généralement un glucide) par des micro-organismes. Elle se déroule entièrement dans le cytoplasme, sans chaîne respiratoire, et ne produit que peu d'ATP ; le substrat n'est que partiellement oxydé, d'où la formation de produits finaux encore riches en énergie (éthanol, acide lactique).
\end{definition}

\textbf{Étape commune : la glycolyse.} Quel que soit le type de fermentation, le glucose est d'abord dégradé par la glycolyse :
\[ \text{Glucose} + 2\,\text{ADP} + 2\,\text{P}_i + 2\,\text{NAD}^+ \longrightarrow 2\,\text{Pyruvate} + 2\,\text{ATP} + 2\,\text{NADH,H}^+ + 2\,\text{H}_2\text{O} \]

Le problème est le suivant : la glycolyse consomme du NAD$^+$ et produit du NADH,H$^+$. Or la cellule ne possède qu'un stock limité de NAD$^+$. En anaérobiose, la chaîne respiratoire ne peut pas régénérer le NAD$^+$. La fermentation résout ce problème en \textbf{réoxydant le NADH,H$^+$ au niveau du pyruvate}.

\textbf{Fermentation alcoolique} (levures, ex. \emph{Saccharomyces cerevisiae}) :
\begin{align*}
\text{Pyruvate} &\xrightarrow{\text{pyruvate décarboxylase}} \text{Acétaldéhyde} + \text{CO}_2 \\
\text{Acétaldéhyde} + \text{NADH,H}^+ &\xrightarrow{\text{alcool déshydrogénase}} \text{Éthanol} + \text{NAD}^+
\end{align*}
Bilan global : $\text{C}_6\text{H}_{12}\text{O}_6 \rightarrow 2\,\text{C}_2\text{H}_5\text{OH} + 2\,\text{CO}_2 + 2\,\text{ATP}$.

\textbf{Fermentation lactique homolactique} (bactéries lactiques, ex. \emph{Lactobacillus}, \emph{Streptococcus}) :
\[ \text{Pyruvate} + \text{NADH,H}^+ \xrightarrow{\text{lactate déshydrogénase}} \text{Lactate} + \text{NAD}^+ \]
Bilan global : $\text{C}_6\text{H}_{12}\text{O}_6 \rightarrow 2\,\text{C}_3\text{H}_6\text{O}_3 + 2\,\text{ATP}$.

\begin{popfigure}
\begin{center}
\begin{tikzpicture}[font=\small, node distance=0.55cm and 1.6cm]
\node[draw=popblue, fill=popblueL, rounded corners, thick, minimum width=2.4cm] (glu) {Glucose};
\node[draw=popblue, fill=popblueL, rounded corners, thick, minimum width=2.4cm, below=1.1cm of glu] (pyr) {2 Pyruvate};
\draw[-latex, very thick, popblue] (glu) -- node[right, align=center, text width=3.2cm]{Glycolyse\\ \textbf{+2 ATP}, +2 NADH,H$^+$} (pyr);
\node[draw=poporange, fill=white, rounded corners, thick, minimum width=2.6cm, below left=1.2cm and 0.2cm of pyr] (acet) {2 Acétaldéhyde};
\node[draw=poporange, fill=white, rounded corners, thick, minimum width=2.6cm, below=1.1cm of acet] (eth) {2 Éthanol};
\node[draw=popgreen, fill=popgreenL, rounded corners, thick, minimum width=2.6cm, below right=1.2cm and 0.2cm of pyr] (lac) {2 Lactate};
\draw[-latex, thick, poporange] (pyr) -- node[above left, align=center, text width=3.4cm]{pyruvate\\ décarboxylase\\ $+$ CO$_2$} (acet);
\draw[-latex, thick, poporange] (acet) -- node[right, align=center, text width=3.6cm]{alcool déshydrogénase\\ NADH,H$^+$ $\rightarrow$ NAD$^+$} (eth);
\draw[-latex, thick, popgreen] (pyr) -- node[above right, align=center, text width=3.4cm]{lactate\\ déshydrogénase\\ NADH,H$^+$ $\rightarrow$ NAD$^+$} (lac);
\node[below=0.15cm of eth, text=popmuted, align=center, text width=3.4cm] {Fermentation alcoolique (levures)};
\node[below=0.15cm of lac, text=popmuted, align=center, text width=3.4cm] {Fermentation lactique (bactéries)};
\end{tikzpicture}
\end{center}
\legende{Voies métaboliques de la fermentation alcoolique et lactique à partir du glucose. Dans les deux cas, la régénération du NAD$^+$ permet à la glycolyse de se poursuivre.}
\end{popfigure}

\begin{propriete}[Rendement énergétique de la fermentation]
Une molécule de glucose ne fournit que \textbf{2 ATP} en fermentation (glycolyse seule), contre environ \textbf{30 à 32 ATP} en respiration aérobie complète. La fermentation est donc une voie énergétiquement peu rentable : les micro-organismes fermentaires doivent consommer de grandes quantités de substrat, ce qui explique l'accumulation rapide des produits finaux (alcool, acide) dans le milieu.
\end{propriete}

\begin{attention}[Erreur fréquente]
Ne dites jamais que la fermentation « produit de l'énergie sans oxygène donc sans ATP » : elle produit bien 2 ATP par glucose, par \textbf{phosphorylation au niveau du substrat} (glycolyse), et non par phosphorylation oxydative. Autre piège : le CO$_2$ n'est produit que par la fermentation \emph{alcoolique}, pas par la fermentation lactique homolactique.
\end{attention}

\courssub{Les grands procédés fermentaires camerounais}

Le Cameroun possède un patrimoine fermentaire remarquable, qui illustre parfaitement la diversité des micro-organismes et des substrats.

\begin{center}
\begin{tabularx}{\linewidth}{|l|X|X|}
\hline
\rowcolor{popblueL} \textbf{Produit} & \textbf{Substrat et principe} & \textbf{Micro-organismes dominants} \\
\hline
Vin de palme & Sève de palmier riche en sucres ; fermentation alcoolique spontanée rapide & Levures (\emph{Saccharomyces} spp.), puis bactéries acétiques si exposition à l'air \\
\hline
Foléré (bissap) & Calices d'hibiscus ; boisson fermentée légèrement acide & Bactéries lactiques, levures \\
\hline
Foutou / bâton de manioc & Tubercules de manioc fermentés en eau (réttage) puis pilés & Bactéries lactiques (\emph{Lactobacillus}, \emph{Leuconostoc}) ; détoxification de l'acide cyanhydrique \\
\hline
Ngondo & Graines de courge fermentées, condiment & Bactéries lactiques, moisissures de surface \\
\hline
Fromage de soja & Lait de soja caillé puis fermenté & Moisissures (\emph{Rhizopus}, \emph{Aspergillus}), bactéries lactiques \\
\hline
\end{tabularx}
\end{center}

Trois rôles essentiels de la fermentation alimentaire sont à retenir :
\begin{itemize}
\item \textbf{Conservation} : l'acide lactique et l'éthanol abaissent le pH et inhibent les germes de détérioration.
\item \textbf{Amélioration nutritionnelle et gustative} : arômes nouveaux, vitamines synthétisées par les micro-organismes, prédigestion des protéines et glucides.
\item \textbf{Détoxification} : la fermentation du manioc élimine les hétérosides cyanogénétiques (qui libèrent de l'acide cyanhydrique, toxique), rendant le tubercule consommable.
\end{itemize}

\begin{popfigure}
\begin{center}
\begin{tikzpicture}[font=\small,
box/.style={draw=popdark, fill=popblueL, rounded corners, thick, minimum width=3.6cm, minimum height=0.85cm, align=center},
arr/.style={-latex, very thick, poporange}]
\node[box, align=center] (a){Récolte et épluchage\\ du manioc};
\node[box, below=0.6cm of a, align=center] (b){Trempage en eau\\ 2 à 4 jours (réttage)};
\node[box, below=0.6cm of b, align=center] (c){Fermentation lactique\\ détoxification (HCN éliminé)};
\node[box, below=0.6cm of c, align=center] (d){Égouttage, essorage\\ et pressage};
\node[box, below=0.6cm of d, align=center] (e){Cuisson et pilage\\ $\rightarrow$ foutou};
\draw[arr] (a) -- (b);
\draw[arr] (b) -- (c);
\draw[arr] (c) -- (d);
\draw[arr] (d) -- (e);
\node[right=0.5cm of c, text=popmuted, align=left, text width=4.2cm] {Bactéries lactiques : baisse du pH, ramollissement des fibres};
\node[right=0.5cm of b, text=popmuted, align=left, text width=4.2cm] {Anaérobiose partielle dans l'eau de trempage};
\end{tikzpicture}
\end{center}
\legende{Diagramme de flux du procédé traditionnel de transformation du manioc en foutou.}
\end{popfigure}

\begin{savaistu}[Le manioc, un aliment rendu sûr par la biotechnologie traditionnelle]
Le manioc contient naturellement des composés cyanogénétiques. Sans fermentation ni trempage, sa consommation régulière peut provoquer des intoxications chroniques. Les procédés traditionnels camerounais (réttage, fermentation) constituent donc une véritable biotechnologie de \textbf{sécurité alimentaire}, mise au point empiriquement bien avant que la chimie n'explique le phénomène.
\end{savaistu}

\courssub{Facteurs influençant la fermentation}

La fermentation est une réaction enzymatique conduite par des êtres vivants : elle dépend donc étroitement des conditions du milieu.

\begin{itemize}
\item \textbf{Température} : chaque micro-organisme possède un optimum (environ 25--30 \textdegree C pour les levures du vin de palme). Une chaleur excessive dénature les enzymes ; le froid ralentit le métabolisme.
\item \textbf{pH} : les bactéries lactiques acidifient le milieu, ce qui finit par inhiber leur propre croissance (autolimitation) et protège le produit des contaminants.
\item \textbf{Concentration en substrat} : un excès de sucre crée une pression osmotique élevée qui peut inhiber les levures ; un déficit limite la production.
\item \textbf{Anaérobiose} : la fermentation alcoolique exige l'absence d'oxygène ; à l'air libre, l'éthanol est oxydé en acide acétique par les bactéries acétiques — c'est pourquoi le vin de palme « tourne au vinaigre » en un ou deux jours.
\item \textbf{Contamination} : des micro-organismes indésirables (moisissures toxigènes, bactéries pathogènes) peuvent s'installer si l'hygiène est mauvaise.
\end{itemize}

\begin{attention}[Risque de contamination]
En l'absence de contrôle hygiénique, des pathogènes dangereux peuvent se développer dans les aliments fermentés ou mal conservés : \emph{Clostridium botulinum} (toxine botulique, mortelle, en milieu anaérobie mal acidifié), \emph{Salmonella}, moisissures productrices d'aflatoxines. L'acidification rapide par les ferments lactiques est la principale barrière naturelle : un produit qui ne s'acidifie pas correctement ne doit \textbf{jamais} être consommé.
\end{attention}

\courssub{Du contrôle artisanal au contrôle quantitatif}

Traditionnellement, le producteur juge sa fermentation par les \textbf{organes des sens} : le goût (acidité, douceur résiduelle), l'odeur (arômes alcooliques ou acides), l'aspect (mousse, dégagement gazeux) et la densité perçue du liquide. Ces méthodes sont rapides mais subjectives et non reproductibles.

Les premières approches quantitatives ont introduit des instruments simples :
\begin{itemize}
\item le \textbf{densimètre} (ou pèse-mout) : la densité du moût diminue au fur et à mesure que le sucre est converti en éthanol ;
\item la \textbf{titration acide} : dosage de l'acidité totale par une solution de soude de concentration connue, en présence d'un indicateur coloré, pour suivre la fermentation lactique.
\end{itemize}

\begin{methode}[Dosage de l'éthanol par distillation suivie de densimétrie]
\begin{enumerate}
\item Prélever un volume connu $V$ de la boisson fermentée (ex. 100 mL de vin de palme).
\item Distiller : l'éthanol, plus volatil que l'eau, s'évapore en premier ; recueillir le distillat.
\item Mesurer la densité $\rho$ du distillat avec un densimètre (ou pycnomètre).
\item Convertir la densité en degré alcoolique à l'aide d'une table de correspondance éthanol--eau (tables officielles).
\item Calculer le titre alcoométrique volumique : $T = \dfrac{V_{\text{éthanol}}}{V_{\text{échantillon}}} \times 100$ (en \% vol).
\end{enumerate}
\end{methode}

\begin{exoresolu}[Calcul d'un titre alcoolique]
Un élève prélève 200 mL de vin de palme. Après distillation et densimétrie, il détermine que le distillat contient 8 mL d'éthanol pur rapportés à 100 mL de boisson. Quel est le titre alcoolique du vin de palme ?
\tcblower
Le titre alcoométrique volumique est le volume d'éthanol pur contenu dans 100 volumes de boisson :
\[ T = \frac{V_{\text{éthanol}}}{V_{\text{boisson}}} \times 100 = \frac{8}{100} \times 100 = 8\ \% \text{ vol.} \]
Le vin de palme titre donc 8 \% vol, valeur cohérente pour un vin de palme bien fermenté (la teneur augmente avec la durée de fermentation, puis l'acétification peut la faire baisser).
\end{exoresolu}

\courssub{De l'artisanat à l'industrie : standardisation et bioréacteurs}

Les procédés industriels modernes reprennent exactement les mêmes réactions, mais en les \textbf{contrôlant} :

\begin{center}
\begin{tabularx}{\linewidth}{|l|X|X|}
\hline
\rowcolor{popblueL} \textbf{Critère} & \textbf{Procédé traditionnel} & \textbf{Procédé industriel} \\
\hline
Micro-organismes & Flore spontanée, variable & Souches sélectionnées, pures, conservées en banques de souches \\
\hline
Conditions & Ambiantes, non contrôlées & Température, pH, oxygénation régulés en bioréacteur \\
\hline
Qualité du produit & Variable d'un lot à l'autre & Standardisée, contrôlée (densité, titration, chromatographie) \\
\hline
Rendement et échelle & Faible, familial ou artisanal & Élevé, cuves de plusieurs milliers de litres \\
\hline
Sécurité & Empirique & Contrôles microbiologiques systématiques \\
\hline
\end{tabularx}
\end{center}

\begin{aretenir}[À retenir de cette section]
\begin{itemize}
\item La biotechnologie traditionnelle exploite empiriquement le métabolisme microbien ; la biotechnologie moderne le pilote au niveau moléculaire.
\item La fermentation est un catabolisme anaérobie incomplet : 2 ATP par glucose, avec régénération du NAD$^+$ par réduction du pyruvate (éthanol + CO$_2$ chez les levures, lactate chez les bactéries lactiques).
\item Les produits camerounais fermentés (vin de palme, foutou, foléré, ngondo, fromage de soja) illustrent la conservation, l'amélioration gustative et la détoxification des aliments.
\item La maîtrise de la température, du pH, de l'anaérobiose et de l'hygiène conditionne la qualité et la sécurité du produit ; l'industrie ajoute souches sélectionnées et bioréacteurs.
\end{itemize}
\end{aretenir}

\begin{exercice}[Fermentation et rendement énergétique]
\begin{enumerate}
\item Une levure dégrade 5 moles de glucose par fermentation alcoolique. Combien de moles d'ATP, d'éthanol et de CO$_2$ sont produites ?
\item Comparer avec la dégradation complète de 5 moles de glucose par respiration aérobie (rendement : 30 ATP par glucose).
\item Expliquer pourquoi un moût très sucré peut ralentir la fermentation.
\end{enumerate}
\end{exercice}

\begin{corrige}[Exercice 1]
\begin{enumerate}
\item D'après le bilan $\text{C}_6\text{H}_{12}\text{O}_6 \rightarrow 2\,\text{C}_2\text{H}_5\text{OH} + 2\,\text{CO}_2 + 2\,\text{ATP}$ : pour 5 mol de glucose, on obtient $5 \times 2 = 10$ mol d'ATP, 10 mol d'éthanol et 10 mol de CO$_2$.
\item Respiration : $5 \times 30 = 150$ mol d'ATP, soit 15 fois plus que la fermentation. La fermentation est donc très peu rentable énergétiquement.
\item Une concentration élevée en sucre augmente la pression osmotique du milieu : l'eau sort des cellules de levure (plasmolyse), ce qui inhibe leur métabolisme et ralentit la fermentation.
\end{enumerate}
\end{corrige}

\coursec{Biotechnologie au service de l'homéostasie, du soutien et de la locomotion}

Dans la section précédente, nous avons vu comment l'être humain exploite des cellules vivantes depuis des millénaires pour fermenter et transformer les aliments. La biotechnologie traditionnelle utilise des organismes entiers dans leur état naturel. La biotechnologie moderne va bien plus loin : elle ne se contente pas de transformer la nourriture, elle peut \emph{réparer, remplacer ou assister le corps humain lui-même}. Lorsqu'un patient diabétique à Yaoundé s'injecte de l'insuline produite par des bactéries génétiquement modifiées, lorsqu'un amputé à Douala remarcher grâce à une prothèse myoélectrique, ou lorsqu'un ligament déchiré est reconstruit sur un échafaudage de polymère biodégradable, la biotechnologie agit directement sur trois fonctions fondamentales étudiées plus tôt dans ce cours : l'\cle{homéostasie}, le \cle{soutien} et la \cle{locomotion}. Cette section montre comment.

\courssub{Rappel sur l'homéostasie : les boucles de rétroaction négative}

\textbf{Intuition première.} Imaginez le système de climatisation d'une salle de classe à Yaoundé. Lorsque la température dépasse la valeur de consigne, un capteur la détecte, un contrôleur déclenche le refroidissement, et la température redescend. Lorsqu'elle devient trop basse, le refroidissement s'arrête. Le système \emph{s'oppose} à tout écart. Votre corps fonctionne exactement de la même manière pour la glycémie, la température corporelle et le pH sanguin.

\begin{definition}[Homéoestasie]
L'\cle{homéostasie} est le maintien d'un milieu intérieur stable (le \emph{milieu intérieur} : plasma sanguin et liquide interstitiel) dans des limites étroites, malgré les variations du milieu extérieur ou de l'activité métabolique. Elle est réalisée par des mécanismes d'autorégulation, principalement des \cle{boucles de rétroaction négative}, impliquant des récepteurs (capteurs), un centre de coordination et des effecteurs.
\end{definition}

\begin{propriete}[Rétroaction négative]
Dans une \cle{boucle de rétroaction négative}, la réponse \emph{réduit ou annule le stimulus initial} : la sortie du système inhibe sa propre production. Exemple : une hausse de la glycémie déclenche la sécrétion d'insuline ; l'insuline abaisse la glycémie ; la baisse de la glycémie stoppe alors la sécrétion d'insuline. (Le mécanisme détaillé est exigible à l'examen et doit être décrit avec précision au niveau Advanced Level.)
\end{propriete}

\textbf{Trois boucles à maîtriser parfaitement.}

\begin{itemize}
\item \textbf{Glycémie (taux de glucose sanguin).} La glycémie à jeun normale est d'environ $0,8$--$1,0\ \mathrm{g\,L^{-1}}$ (soit environ $4,4$--$5,5\ \mathrm{mmol\,L^{-1}}$). Après un repas riche en glucides (pensez à un plat de \emph{ndolé} avec du miondo), la glycémie augmente : les cellules $\beta$ des îlots de Langerhans sécrètent de l'\cle{insuline}, qui stimule la capture du glucose par les muscles et le tissu adipeux et son stockage sous forme de glycogène dans le foie (glycogenèse). Au cours du jeûne, les cellules $\alpha$ sécrètent du \cle{glucagon}, qui déclenche la glycogenolyse et la néoglucogenèse, remontant la glycémie. Les deux hormones sont \emph{antagonistes}.
\item \textbf{Température corporelle.} L'hypothalamus agit comme un thermostat. La surchauffe déclenche la vasodilatation des artérioles cutanées et la transpiration (l'évaporation évacue la chaleur) ; le refroidissement déclenche la vasoconstriction, les frissons (contractions musculaires libérant de la chaleur) et une augmentation de la thyroxine et de l'adrénaline qui élève le métabolisme de base.
\item \textbf{pH sanguin.} Le pH sanguin est maintenu près de $7,4$ par des systèmes tampons (le couple \ce{H2CO3}/\ce{HCO3-} est le principal), par les poumons (l'expiration de \ce{CO2} élimine de l'acide) et par les reins (excrétion de \ce{H+} et réabsorption de \ce{HCO3-}).
\end{itemize}

\begin{popfigure}
\begin{tikzpicture}[font=\small, >=stealth, node distance=1.2cm]
\tikzset{
  box/.style={draw=popblue, thick, rounded corners=2pt, fill=popblueL, text width=3.0cm, align=center, minimum height=1.0cm},
  eff/.style={draw=popgreen, thick, rounded corners=2pt, fill=popgreenL, text width=3.2cm, align=center, minimum height=1.0cm},
  arr/.style={->, thick, popdark}
}
\node[box, align=center] (stim){\textbf{Stimulus}\\ Glycémie $> 1,0\ \mathrm{g\,L^{-1}}$};
\node[box, right=2.8cm of stim, align=center] (rec){\textbf{Récepteur}\\ Cellules $\beta$ des îlots de Langerhans};
\node[eff, right=2.8cm of rec, align=center] (eff){\textbf{Réponse effectrice}\\ Sécrétion d'insuline};
\node[eff, below=1.5cm of eff, align=center] (target){\textbf{Organes cibles}\\ Foie, muscle, tissu adipeux};
\node[box, below=1.5cm of stim, align=center] (res){\textbf{Résultat}\\ Glycémie revient à la normale};
\draw[arr] (stim) -- (rec);
\draw[arr] (rec) -- (eff);
\draw[arr] (eff) -- (target);
\draw[arr] (target.west) -| (res.east);
\draw[->, thick, poporange, dashed] (res.north) -- (stim.south) node[midway, right, text=poporange, align=left, text width=2.5cm]{\footnotesize rétroaction : stimulus supprimé};
\node[align=left, text width=4.8cm, anchor=west, fill=poppink!15, draw=poppink, rounded corners=2pt, inner sep=3pt] at ($(target)+(2.5,0.2)$) {\footnotesize \textbf{Jeûne :} cellules $\alpha$ $\rightarrow$ \textbf{glucagon} $\rightarrow$ glycogenolyse $\rightarrow$ glycémie remonte (boucle antagoniste).};
\end{tikzpicture}
\legende{Rétroaction négative du contrôle de la glycémie. L'insuline abaisse la glycémie ; le glucagon la remonte. La réponse s'oppose toujours à l'écart initial.}
\end{popfigure}

\begin{attention}[Confusion entre les deux hormones]
Une erreur classique au GCE est d'écrire que le glucagon abaisse la glycémie, ou que l'insuline est sécrétée pendant le jeûne. Retenir : l'\textbf{ins}uline met le glucose \textbf{in}side les cellules et en réserve (glycémie baisse) ; le \textbf{gluc}agon est libéré quand le glucose est parti (\textbf{gone} en anglais, glycémie monte). De plus, la rétroaction négative ne signifie pas que la variable est constante : elle \emph{oscille} autour d'un point de consigne.
\end{attention}

\courssub{Insuline humaine recombinante : du gène à la protéine thérapeutique}

Lorsque la boucle de l'insuline fait défaut, le résultat est le \cle{diabète sucré}. Dans le \textbf{diabète de type 1}, les cellules $\beta$ sont détruites (réaction auto-immune) et aucune insuline n'est produite ; la survie impose des injections quotidiennes d'insuline. Pendant des décennies, l'insuline était extraite de pancréas de porcs ou de bovins provenant d'abattoirs. L'insuline animale fonctionne, mais elle diffère légèrement de l'insuline humaine (l'insuline bovine diffère de trois acides aminés), provoquant des réactions allergiques et une production d'anticorps chez certains patients, et l'approvisionnement était limité.

La technologie de l'ADN recombinant (détaillée en Section 4) a résolu ce problème. La logique est simple : donner à une bactérie le gène humain de l'insuline, et la bactérie devient une usine microscopique.

\begin{methode}[Production d'insuline humaine recombinante]
\begin{enumerate}
\item \textbf{Isoler le gène.} Le gène humain de l'insuline est obtenu (ou synthétisé à partir de l'ARNm des cellules $\beta$ grâce à la transcriptase inverse, donnant un ADNc sans introns).
\item \textbf{Insérer dans un vecteur.} Le gène est coupé et ligaturé dans un plasmide à l'aide d'enzymes de restriction et d'ADN ligase, placé sous le contrôle d'un promoteur bactérien fort.
\item \textbf{Transformer l'hôte.} Le plasmide recombinant est introduit dans \emph{Escherichia coli} (ou la levure \emph{Saccharomyces cerevisiae}).
\item \textbf{Sélectionner et cultiver.} Les cellules transformées sont sélectionnées (marqueur de résistance aux antibiotiques) et cultivées en grands fermenteurs industriels sous température, pH et aération contrôlés.
\item \textbf{Extraire et purifier.} Les chaînes d'insuline sont récoltées, les deux chaînes (A et B) sont repliées correctement et reliées par des ponts disulfure, puis purifiées et formulées pour injection.
\end{enumerate}
\end{methode}

\begin{savaistu}[Insuline et Cameroun]
Le diabète est un problème de santé publique croissant au Cameroun, favorisé par l'urbanisation et le changement des régimes alimentaires. L'insuline humaine recombinante, produite industriellement depuis le début des années 1980, a été la \emph{première} protéine thérapeutique jamais obtenue par génie génétique. Comme elle est identique à l'hormone humaine, elle provoque beaucoup moins de réactions immunitaires que l'insuline animale, et sa production peut être montée à l'échelle pour répondre à la demande mondiale — un avantage décisif là où l'accès au traitement est déjà difficile.
\end{savaistu}

\courssub{Biomatériaux pour le soutien tissulaire}

L'homéostasie n'est pas la seule fonction que la biotechnologie peut restaurer. Os, cartilage, ligaments et peau assurent le \cle{soutien} ; lorsqu'ils sont détruits par un traumatisme, une infection ou un cancer, les ingénieurs les remplacent aujourd'hui par des \textbf{biomatériaux}.

\begin{definition}[Biomatériau]
Un \cle{biomatériau} est un matériau (naturel ou synthétique) conçu pour interagir avec un système biologique à des fins thérapeutiques : remplacer, soutenir ou régénérer un tissu ou un organe, tout en étant toléré par l'organisme (\emph{biocompatibilité}).
\end{definition}

\textbf{Trois familles à connaître.}

\begin{itemize}
\item \textbf{Hydrogels.} Réseaux de polymères hydrophiles gonflés d'eau (jusqu'à $90\ \%$ d'eau). Leur souplesse et leur teneur en eau rappellent le tissu vivant, ils servent donc pour les lentilles de contact, les pansements qui maintiennent les brûlures humides, et comme échafaudages libérant des cellules ou des facteurs de croissance vers du cartilage lésé.
\item \textbf{Échafaudages de collagène.} Le collagène est la protéine structurelle la plus abondante de l'os, de la peau et des tendons. Du collagène purifié (souvent d'origine bovine, traité pour éliminer les parties antigéniques) est façonné en éponges ou membranes qui servent de \emph{squelette temporaire} : les propres cellules du patient le colonisent, sécrètent leur matrice, et l'échafaudage est progressivement remplacé par du tissu naturel. C'est le principe de l'\cle{ingénierie tissulaire}.
\item \textbf{Polymères biodégradables : PLA et PGA.} L'acide polylactique (PLA) et l'acide polyglycolique (PGA) sont des polyestères synthétiques qui s'hydrolysent dans l'organisme en acide lactique et acide glycolique — métabolites naturels éliminés par la respiration et les urines. Ils servent pour les fils de suture résorbables, les broches et vis de fixation osseuse, et les échafaudages poreux pour la régénération osseuse : l'implant soutient le tissu pendant la cicatrisation, puis disparaît, évitant une seconde chirurgie pour l'enlever.
\end{itemize}

\begin{attention}[Rejet et infection à l'interface]
Aucun biomatériau n'est parfaitement inerte. Les deux grands dangers à l'interface biomatériau-tissu sont : (i) le \textbf{rejet immunitaire ou l'inflammation chronique} si le matériau est reconnu comme étranger (l'organisme l'isole par du tissu fibreux, altérant la fonction) ; (ii) l'\textbf{infection}, car les bactéries peuvent former un \emph{biofilm} à la surface de l'implant que les antibiotiques pénètrent mal. D'où la nécessité absolue de tests de biocompatibilité, de conditions chirurgicales stériles, et parfois d'immunosuppression. N'écrivez jamais en examen qu'un matériau artificiel est ``accepté automatiquement'' par le corps.
\end{attention}

\courssub{Prothèses et exosquelettes : restaurer la locomotion}

Le soutien donne au corps son cadre ; la \cle{locomotion} résulte de muscles tirant sur ce cadre sous contrôle nerveux. Lorsqu'un membre est perdu — après des accidents de la route, hélas fréquents sur les routes camerounaises, ou après amputation pour gangrène diabétique — la biotechnologie propose des \textbf{prothèses myoélectriques}.

\textbf{Principe.} Lorsqu'un muscle se contracte, même un muscle résiduel de moignon, il génère un minuscule signal électrique (quelques millivolts) : le \cle{signal électromyographique (EMG)}. Des électrodes de surface placées sur la peau du moignon détectent ce signal ; un microprocesseur (le contrôleur) l'interprète et commande des moteurs électriques (actionneurs) qui bougent l'articulation artificielle. Le patient commande la prothèse simplement en contractant les muscles qui actionnaient autrefois le membre manquant. Les systèmes plus avancés ajoutent une \textbf{interface neurale-machine} : électrodes connectées aux nerfs périphériques, voire implantées dans le cortex moteur, qui lisent l'intention de mouvement directement, et un retour sensoriel (capteurs de pression stimulant la peau) qui restaure une forme de toucher. Les \textbf{exosquelettes} appliquent la même logique à l'extérieur : un cadre motorisé attaché aux jambes d'un patient paraplégique détecte les mouvements du tronc et commande les articulations de hanche et de genou pour reproduire la marche.

\begin{popfigure}
\begin{tikzpicture}[font=\small, >=stealth]
\tikzset{
  part/.style={draw=popblue, thick, rounded corners=2pt, fill=popblueL, align=center},
  lbl/.style={align=left, text width=3.8cm, font=\footnotesize}
}
% socket
\draw[part, fill=popgreenL, draw=popgreen] (0,4.2) ellipse (1.1 and 0.55);
\node at (0,4.2) {emboîture (socket)};
% EMG electrodes
\fill[poppink] (-0.9,4.6) circle (0.08); \fill[poppink] (0.9,4.6) circle (0.08);
% controller
\node[part, minimum width=1.6cm, minimum height=0.7cm] (ctrl) at (0,3.1) {contrôleur};
% battery
\node[part, minimum width=1.2cm, minimum height=0.6cm, fill=popgold!25, draw=popgold] (bat) at (0,2.2) {batterie};
% knee actuator
\node[part, fill=poporange!20, draw=poporange, minimum width=1.4cm, minimum height=1.0cm] (knee) at (0,1.0) {moteur genou};
% shank
\draw[popdark, very thick] (0,0.5) -- (0,-0.9);
% force sensor + foot
\node[part, minimum width=1.0cm, minimum height=0.4cm, fill=poppurple!15, draw=poppurple] (fs) at (0,-1.1) {};
\draw[part, fill=popmuted!20] (-0.9,-1.6) -- (1.0,-1.6) -- (1.0,-1.35) -- (-0.2,-1.35) -- (-0.9,-1.45) -- cycle;
\node at (0.35,-1.75) {\footnotesize pied};
% labels
\node[lbl, anchor=west] (l1) at (2.4,4.6) {\textbf{Électrodes EMG} détectent signaux muscles moignon};
\draw[->, poppink] (l1.west) -- (0.9,4.6);
\node[lbl, anchor=west] (l2) at (2.4,3.1) {\textbf{Microprocesseur} décode intention mouvement};
\draw[->, popblue] (l2.west) -- (ctrl.east);
\node[lbl, anchor=west] (l3) at (2.4,1.0) {\textbf{Actionneur} (moteur) commande flexion/extension genou};
\draw[->, poporange] (l3.west) -- (knee.east);
\node[lbl, anchor=east] (l4) at (-2.4,4.2) {\textbf{Emboîture} : ajustée sur mesure au moignon};
\draw[->, popgreen] (l4.east) -- (-1.1,4.2);
\node[lbl, anchor=east] (l5) at (-2.4,2.2) {\textbf{Batterie} : alimente moteurs et électronique};
\draw[->, popgold] (l5.east) -- (bat.west);
\node[lbl, anchor=east] (l6) at (-2.4,-1.2) {\textbf{Capteur force} : mesure contact sol pour retour};
\draw[->, poppurple] (l6.east) -- (fs.west);
\end{tikzpicture}
\legende{Vue éclatée d'une prothèse de genou myoélectrique. Les signaux EMG des muscles du moignon sont décodés par le contrôleur, qui commande l'actionneur du genou ; les capteurs de force au pied fournissent un retour pour une marche stable.}
\end{popfigure}

\courssub{Évaluation fonctionnelle : le dispositif marche-t-il vraiment ?}

Une prothèse ou un implant n'a de valeur que s'il améliore la vie du patient \emph{mesurablement}. Biotechnologues et kinésithérapeutes utilisent donc des tests fonctionnels standardisés.

\begin{methode}[Protocole d'évaluation d'une prothèse de membre inférieur]
\begin{enumerate}
\item \textbf{Test de marche de six minutes (6MWT).} Le patient marche le plus loin possible en 6 minutes sur un couloir plat ; la distance totale (en mètres) mesure l'endurance. Le test est répété avant l'appareillage et après rééducation pour quantifier le progrès.
\item \textbf{Analyse de la marche.} Caméras et marqueurs réfléchissants enregistrent les angles articulaires ; une plateforme de force intégrée au sol mesure les forces de réaction du sol, révélant les asymétries entre le membre sain et la prothèse.
\item \textbf{Capteurs de force.} Semelles instrumentées mesurent la distribution des pressions plantaires, détectant les zones de surcharge pouvant causer des ulcères cutanés.
\item \textbf{Score de satisfaction.} Un questionnaire standardisé (confort, douleur, facilité d'utilisation, image corporelle) est noté, car un dispositif techniquement parfait que le patient refuse de porter est un échec.
\item \textbf{Comparer et décider.} Les résultats sont comparés aux valeurs de référence et à la ligne de base du patient ; les réglages du contrôleur sont alors ajustés.
\end{enumerate}
\end{methode}

\begin{exoresolu}[Interprétation d'un test de marche de six minutes]
Un patient appareillé d'une prothèse de genou myoélectrique parcourt $210\ \mathrm{m}$ au 6MWT lors du premier appareillage. Après 8 semaines de rééducation à la marche, il parcourt $330\ \mathrm{m}$. (a) Calculer le pourcentage d'amélioration. (b) Sa vitesse moyenne lors du second test, en $\mathrm{m\,s^{-1}}$. (c) Pourquoi un questionnaire de satisfaction reste-t-il nécessaire ?
\tcblower
(a) Amélioration $= \dfrac{330-210}{210}\times 100 = \dfrac{120}{210}\times 100 \approx 57\ \%$.\\
(b) Vitesse $= \dfrac{\text{distance}}{\text{temps}} = \dfrac{330\ \mathrm{m}}{360\ \mathrm{s}} \approx 0,92\ \mathrm{m\,s^{-1}}$.\\
(c) La distance seule ne dit rien sur la douleur, le confort de l'emboîture ou la confiance. Un patient peut marcher plus loin tout en souffrant d'irritations cutanées ou en trouvant l'appareil trop lourd ; seule l'évaluation combinée objective (distance, symétrie de la marche) \emph{et} subjective (score de satisfaction) valide la prothèse.
\end{exoresolu}

\courssub{Considérations éthiques et réglementaires}

Ces technologies touchent directement le corps humain, elles soulèvent donc des questions éthiques spécifiques que l'examinateur du GCE attend que vous discutiez de manière équilibrée.

\begin{itemize}
\item \textbf{Consentement éclairé.} Le patient doit comprendre les bénéfices, les risques (infection, rejet, défaillance du dispositif) et les alternatives avant de recevoir un implant ou de rejoindre un essai clinique. Le consentement doit être libre, éclairé et révocable.
\item \textbf{Accessibilité et équité.} Les prothèses myoélectriques et l'insuline recombinante sont coûteuses. Est-il acceptable qu'un dispositif changeant la vie soit accessible aux patients aisés mais pas à un paysan de la région du Nord-Ouest ? La biotechnologie doit s'accompagner de politiques (subventions, production locale, insulines génériques) qui réduisent les inégalités.
\item \textbf{Brevets sur les dispositifs médicaux.} Les brevets récompensent l'innovation et financent la recherche, mais ils peuvent maintenir les prix élevés pendant des années. Le débat entre protection des inventeurs et protection des patients est central en biotechnologie de santé.
\item \textbf{Sécurité et régulation.} Tout biomatériau ou dispositif doit passer des tests de biocompatibilité, des essais cliniques et l'approbation des autorités de régulation avant mise sur le marché ; une surveillance à long terme se poursuit après commercialisation.
\item \textbf{Augmentation humaine.} Si un exosquelette peut un jour permettre à une personne saine de courir plus vite, où se situe la frontière entre \emph{thérapie} (restaurer une fonction perdue) et \emph{augmentation} (améliorer une fonction normale) ? La société, pas les seuls scientifiques, doit répondre.
\end{itemize}

\begin{aretenir}[Points clés de cette section]
\begin{itemize}
\item L'homéostasie (glycémie, température, pH sanguin) est maintenue par des \cle{boucles de rétroaction négative} : la réponse annule le stimulus.
\item L'insuline humaine recombinante, produite par des bactéries ou levures génétiquement modifiées, a restauré une boucle homéostatique défaillante et transformé le traitement du diabète de type 1.
\item Les \cle{biomatériaux} (hydrogels, échafaudages de collagène, polymères biodégradables PLA/PGA) remplacent ou régénèrent les tissus de soutien ; les risques majeurs sont le rejet immunitaire et l'infection.
\item Les prothèses myoélectriques et les exosquelettes convertissent les \cle{signaux EMG} en mouvement via électrodes, contrôleur et actionneurs ; les interfaces neurales vont plus loin en lisant directement les signaux nerveux.
\item La fonction est évaluée objectivement (test de marche 6 min, analyse de la marche, capteurs de force) et subjectivement (scores de satisfaction).
\item Consentement, accès équitable, brevets et régulation de sécurité encadrent l'usage éthique de ces technologies.
\end{itemize}
\end{aretenir}

\begin{exercice}[Vérifiez votre compréhension]
\begin{enumerate}
\item Dessinez et annotez la boucle de rétroaction négative qui restaure la glycémie après un repas riche en glucides.
\item Expliquez pourquoi l'insuline humaine recombinante provoque moins de réactions immunitaires que l'insuline bovine.
\item Donnez un avantage et un risque de l'utilisation d'une vis en PLA biodégradable au lieu d'une vis métallique pour fixer une fracture osseuse.
\item Tracez le chemin de l'information dans une prothèse myoélectrique, de l'intention du patient au mouvement du genou.
\item Un patient parcourt $250\ \mathrm{m}$ au 6MWT avant rééducation et $325\ \mathrm{m}$ après. Calculez le pourcentage d'amélioration et commentez pourquoi l'analyse de la marche doit compléter ce résultat.
\end{enumerate}
\end{exercice}

\begin{corrige}[Exercice 2]
\begin{enumerate}
\item Repas $\rightarrow$ glycémie monte $\rightarrow$ cellules $\beta$ îlots Langerhans stimulées $\rightarrow$ insuline sécrétée $\rightarrow$ capture glucose muscle/adipose et glycogenèse foie $\rightarrow$ glycémie baisse $\rightarrow$ sécrétion insuline stoppée (rétroaction).
\item L'insuline recombinante a exactement la séquence d'acides aminés de l'insuline humaine, alors que l'insuline bovine diffère de trois acides aminés ; ces séquences étrangères peuvent être reconnues comme antigènes, déclenchant production d'anticorps et allergie.
\item Avantage : elle s'hydrolyse en acide lactique, métabolite naturel, donc pas de seconde opération pour l'enlever. Risque : inflammation ou infection à l'interface tissu-matériau, ou rupture mécanique avant consolidation complète.
\item Intention $\rightarrow$ contraction muscles moignon $\rightarrow$ signal EMG détecté électrodes surface $\rightarrow$ microprocesseur décode signal $\rightarrow$ commande envoyée actionneur genou (moteur) $\rightarrow$ mouvement articulation ; capteurs force fournissent retour pour ajustement.
\item Amélioration $= \dfrac{325-250}{250}\times 100 = 30\ \%$. Le 6MWT mesure l'endurance seulement ; l'analyse de la marche révèle la \emph{qualité} de la marche (symétrie, angles articulaires, boiterie), que la distance seule ne montre pas.
\end{enumerate}
\end{corrige}

\coursec{Correction et remédiation génique : thérapie génique et édition du génome}

In the previous section, we explored how biotechnology provides external support for homeostasis, mobility and structural repair — from biosensors that monitor blood glucose to engineered prostheses and tissue scaffolds. We now turn to a more radical approach: instead of managing the symptoms of a genetic defect, we correct the defect at its source, the DNA itself. This is the promise of \cle{gene therapy} and \cle{genome editing}. Imagine a patient with sickle cell disease, a condition highly prevalent in Cameroon, whose red blood cells sickle because of a single nucleotide mutation in the \emph{HBB} gene. What if we could precisely edit that mutation back to the normal sequence, or introduce a functional copy of the gene? This section unpacks the molecular tools, the clinical realities, and the ethical boundaries of rewriting the human genome.

\courssub{Monogenic Diseases and Therapeutic Strategies}

Many debilitating inherited disorders are \cle{monogenic diseases} — caused by mutations in a single gene. Classic examples include \emph{sickle cell disease} (mutation in \emph{HBB}, $\beta$-globin), \emph{cystic fibrosis} (mutation in \emph{CFTR}, chloride channel), \emph{Duchenne muscular dystrophy} (\emph{DMD}, dystrophin), and \emph{Leber congenital amaurosis} (\emph{RPE65}, retinal pigment epithelium). The therapeutic logic is straightforward: if a disease results from a loss-of-function mutation, restoring functional protein should alleviate the phenotype. Two broad strategies exist:

\begin{definition}[Gene Therapy]
Gene therapy is the introduction of nucleic acids (DNA or RNA) into the cells of a patient to treat or prevent a disease. It aims to modify the genetic information of target cells to achieve a therapeutic effect.
\end{definition}

\begin{itemize}
    \item \textbf{Gene addition (augmentation):} A functional copy of the defective gene (\cle{transgene}) is delivered into the nucleus, often remaining episomal (non-integrating) or integrating randomly. The endogenous mutant allele is untouched. This suits recessive loss-of-function disorders.
    \item \textbf{Gene editing (correction):} The endogenous mutated sequence is directly modified — corrected, disrupted, or rewritten — at its native chromosomal locus. This preserves physiological regulation and avoids insertional mutagenesis. The breakthrough tool is \cle{CRISPR-Cas9}.
\end{itemize}

\courssub{Gene Addition Therapy: Viral and Non-Viral Vectors}

Delivering a transgene (often 2–5 kb of cDNA plus promoter and polyA signal) into the nucleus of specific cells in a living patient is a formidable delivery problem. The vehicle is called a \cle{vector}.

\begin{definition}[Viral Vector]
A viral vector is a virus that has been genetically modified to remove its pathogenic genes and replication capacity, while retaining its ability to enter target cells and deliver a therapeutic transgene cassette.
\end{definition}

\begin{popfigure}
\begin{tikzpicture}[>=stealth, node distance=1.2cm, font=\small]
% Styles
\tikzset{
    proc/.style={rectangle, rounded corners, draw=popblue, thick, fill=popblueL, text width=2.8cm, align=center, minimum height=1.2cm},
    arrow/.style={->, thick, draw=popdark},
    labelnode/.style={rectangle, draw=poporange, fill=poporange!10, text width=3cm, align=center, font=\tiny}
}
% Nodes
\node[proc, align=center] (entry){1. Binding \& Entry\\ (Receptor-mediated\\ endocytosis)};
\node[proc, right=of entry, align=center] (endosome){2. Endosomal\\ Escape\\ (Low pH triggers\\ capsid change)};
\node[proc, right=of endosome, align=center] (cyto){3. Cytoplasmic\\ Trafficking\\ (Microtubule\\ transport)};
\node[proc, right=of cyto, align=center] (nuclear){4. Nuclear Entry\\ (Nuclear pore\\ complex)};
\node[proc, right=of nuclear, align=center] (uncoat){5. Uncoating \&\\ 2nd Strand\\ Synthesis};
\node[proc, right=of uncoat, align=center] (episome){6. Episomal\\ Concatemer\\ Formation};
\node[proc, right=of episome, align=center] (expr){7. Transgene\\ Expression\\ (mRNA $\rightarrow$ Protein)};

% Arrows
\draw[arrow] (entry) -- (endosome);
\draw[arrow] (endosome) -- (cyto);
\draw[arrow] (cyto) -- (nuclear);
\draw[arrow] (nuclear) -- (uncoat);
\draw[arrow] (uncoat) -- (episome);
\draw[arrow] (episome) -- (expr);

% Annotations
\node[labelnode, above=0.3cm of entry, align=center] (a1){AAV capsid\\ binds receptor\\(e.g. AAV2: HSPG)};
\node[labelnode, above=0.3cm of endosome, align=center] (a2){VP1 PLA2 domain\\ exposes, membrane\\ penetration};
\node[labelnode, above=0.3cm of cyto, align=center] (a3){Dynein motor\\ on microtubules};
\node[labelnode, above=0.3cm of nuclear, align=center] (a4){Importin-mediated\\ through NPC};
\node[labelnode, above=0.3cm of uncoat, align=center] (a5){Host DNA pol\\ synthesizes (+) strand};
\node[labelnode, above=0.3cm of episome, align=center] (a6){Head-to-tail\\ circles, stable\\ in non-dividing cells};
\node[labelnode, above=0.3cm of expr, align=center] (a7){Promoter drives\\ therapeutic protein};

\draw[dashed, thin, popmuted] (entry.north) -- (a1.south);
\draw[dashed, thin, popmuted] (endosome.north) -- (a2.south);
\draw[dashed, thin, popmuted] (cyto.north) -- (a3.south);
\draw[dashed, thin, popmuted] (nuclear.north) -- (a4.south);
\draw[dashed, thin, popmuted] (uncoat.north) -- (a5.south);
\draw[dashed, thin, popmuted] (episome.north) -- (a6.south);
\draw[dashed, thin, popmuted] (expr.north) -- (a7.south);
\end{tikzpicture}
\legende{Life cycle of an AAV vector in a target cell. AAV remains predominantly episomal, making it ideal for non-dividing cells (neurons, muscle, retina) but diluted in dividing tissues.}
\end{popfigure}

\begin{center}
\begin{tabularx}{\linewidth}{|l|X|X|X|}\hline
\rowcolor{popblueL}
\textbf{Vector} & \textbf{Genome / Capacity} & \textbf{Integration} & \textbf{Key Features / Use Cases} \\ \hline
\textbf{Adenovirus} & dsDNA, $\sim$36 kb\\(up to 8–10 kb insert) & Episomal (non-integrating) & High immunogenicity (strong innate/adaptive response). Transient expression. Used in vaccines (e.g. COVID-19 AstraZeneca, J\&J) and some cancer therapies. \\ \hline
\textbf{AAV (Adeno-Associated Virus)} & ssDNA, $\sim$4.7 kb\\(max $\sim$4.5 kb insert) & Mostly episomal (concateners); low frequency integration at \emph{AAVS1} & Low immunogenicity, many serotypes for tissue tropism (AAV2-retina, AAV8/9-liver, AAV9-CNS). Gold standard for \emph{in vivo} gene therapy (Luxturna, Zolgensma). \\ \hline
\textbf{Lentivirus (HIV-1 based)} & ssRNA, $\sim$9 kb\\(up to 8 kb insert) & Integrating (random in active genes) & Transduces dividing \& non-dividing cells. Stable expression in progeny. Used \emph{ex vivo} for HSPCs (Strimvelis, Zynteglo, Casgevy). Risk of insertional oncogenesis (mitigated by SIN LTRs, insulators). \\ \hline
\textbf{Non-viral (Liposomes, LNPs, Nanoparticles)} & DNA/RNA, flexible size & Episomal (transient) or integrating (if transposase co-delivered) & Low immunogenicity, scalable manufacturing, no strict size limit. LNPs revolutionized mRNA vaccines. Lower transduction efficiency \emph{in vivo} for DNA; improving with targeted ligands. \\ \hline
\end{tabularx}
\end{center}

\begin{methode}[Choosing a Vector for Gene Therapy]
\textbf{Step 1: Transgene size.} $>4.5$ kb $\rightarrow$ exclude AAV; consider Lentivirus, Adenovirus, or non-viral.
\textbf{Step 2: Target tissue \& division status.} Non-dividing (neurons, muscle, retina) $\rightarrow$ AAV ideal. Dividing (HSPCs, T cells) $\rightarrow$ Lentivirus (integrating) or \emph{ex vivo} editing.
\textbf{Step 3: Duration of expression needed.} Lifelong $\rightarrow$ Integrating (Lentivirus) or stable episome in non-dividing cells (AAV). Transient $\rightarrow$ Adenovirus, mRNA-LNP.
\textbf{Step 4: Immunogenicity constraints.} Pre-existing neutralizing antibodies (NAbs) against AAV serotypes are common (30–60\% of humans). Screen patients; consider alternate serotypes or immunosuppression.
\textbf{Step 5: Safety profile.} Insertional mutagenesis risk (Lentivirus) vs. immune toxicity (Adenovirus) vs. off-target editing (CRISPR). Regulatory precedent matters.
\end{methode}

\courssub{Genome Editing: The CRISPR-Cas9 Revolution}

While gene addition adds a functional copy, \cle{genome editing} rewrites the existing sequence. The discovery of the bacterial adaptive immune system CRISPR-Cas (Clustered Regularly Interspaced Short Palindromic Repeats, CRISPR-associated proteins) and its repurposing by Doudna, Charpentier, and Zhang (2012–2013) transformed biology.

\begin{popfigure}
\begin{tikzpicture}[>=stealth, font=\small, scale=0.9]
% DNA
\draw[popdark, very thick] (-4,1.5) -- (4,1.5) node[right] {5'};
\draw[popdark, very thick] (-4,-1.5) -- (4,-1.5) node[right] {3'};
% Base pairs
\foreach \x in {-3.5,-3,-2.5,-2,-1.5,-1,-0.5,0,0.5,1,1.5,2,2.5,3,3.5} {
    \draw[popmuted, thin] (\x,1.5) -- (\x,-1.5);
}
% PAM site (NGG) on non-target strand (top)
\node[draw=poporange, thick, fill=poporange!20, rounded corners, inner sep=2pt] at (2.5,1.8) {PAM: 5'-NGG-3'};
\draw[poporange, thick, <->] (2.5,1.6) -- (2.5,1.5);
% gRNA
\draw[poppurple, very thick] (-3.5,3) .. controls (-3.5,2.5) and (-1,2.5) .. (-1,2) node[left, text width=2cm, align=right] {gRNA: scaffold};
\draw[popgreen, very thick] (-1,2) -- (2,2) node[midway, above] {spacer (20 nt)};
\draw[popgreen, thick, decorate, decoration={brace, amplitude=3pt}] (-1,2.2) -- (2,2.2);
% Cas9
\draw[popblue, thick, fill=popblueL, rounded corners=5pt] (-0.5,0.5) rectangle (2.5,-0.5) node[midway, text width=2.5cm, align=center] {Cas9\\ (RuvC + HNH domains)};
% Recognition
\draw[<->, popred, thick] (-1,1.5) -- (-1,2);
\node[popred, above] at (-1,2.3) {Base pairing};
% Cleavage sites
\draw[popred, very thick, dashed] (1.5,1.5) -- (1.5,-1.5);
\node[popred, right] at (1.5,1.5) {Cut (3 bp upstream PAM)};
% Labels
\node[align=center] at (-3.5, -2.2) {Target DNA\\ (Double-stranded)};
\node[align=center, poppurple] at (-3.5, 3.3) {gRNA\\ (crRNA + tracrRNA)};
\node[align=center, popblue] at (1, -2.2) {Cas9 Protein};
\end{tikzpicture}
\legende{CRISPR-Cas9 mechanism. The gRNA spacer (20 nt) base-pairs with the target DNA strand adjacent to a PAM (NGG for \emph{S. pyogenes} Cas9). Cas9 undergoes conformational change; HNH domain cuts the target (complementary) strand, RuvC domain cuts the non-target strand, generating a blunt double-strand break (DSB) 3 bp upstream of the PAM.}
\end{popfigure}

\begin{propriete}[DNA Repair Pathways After CRISPR-Cas9 Cleavage]
A double-strand break (DSB) is repaired by two main cellular pathways:
\begin{enumerate}
    \item \textbf{NHEJ (Non-Homologous End Joining):} Error-prone, ligates ends directly. Frequent in all cell types, especially G1. Creates small insertions/deletions (indels) $\rightarrow$ gene knockout (frameshift). \textbf{High frequency, imprecise.}
    \item \textbf{HDR (Homology-Directed Repair):} Uses a homologous template (sister chromatid or exogenous donor DNA). Precise correction/insertion possible. \textbf{Restricted to S/G2 phases; low frequency in vivo; requires donor template.}
\end{enumerate}
\textbf{Examinable:} You must be able to explain why HDR is inefficient in non-dividing cells (neurons, muscle) and why NHEJ dominates \emph{in vivo}.
\end{propriete}

\begin{attention}[Off-Target Effects of CRISPR]
CRISPR-Cas9 can cleave DNA at sites with sequence similarity to the target (mismatches tolerated, especially in the PAM-distal "seed" region). These \cle{off-target mutations} are random, potentially disrupting tumor suppressor genes or activating oncogenes. \textbf{Mitigation:} High-fidelity Cas9 variants (eSpCas9, HiFi Cas9), truncated gRNAs (17–18 nt), RNP delivery (short exposure), and rigorous off-target screening (GUIDE-seq, CIRCLE-seq, DISCOVER-seq) in preclinical studies.
\end{attention}

\courssub{Clinical Case Studies}

\begin{exemplebox}[Leber Congenital Amaurosis Type 2 (LCA2) — \emph{RPE65} Gene Addition]
\textbf{Disease:} Autosomal recessive blindness from infancy. \emph{RPE65} encodes an isomerase essential for the visual cycle. Mutations $\rightarrow$ no functional visual pigment regeneration.
\textbf{Therapy:} \cle{Voretigene neparvovec (Luxturna\textsuperscript{\textregistered})}, AAV2 vector carrying human \emph{RPE65} cDNA. Subretinal injection (detaches retina locally, delivers vector to RPE cells).
\textbf{Outcome:} Improved light sensitivity, navigational ability in low light. First FDA-approved \emph{in vivo} gene therapy (2017). Demonstrates: AAV + non-dividing target + subretinal immune privilege = success.
\end{exemplebox}

\begin{exemplebox}[Sickle Cell Disease (SCD) / \emph{$\beta$-Thalassemia} — \cle{CTX001 (Casgevy\textsuperscript{\textregistered})} — Genome Editing]
\textbf{Disease:} $\beta$-globin mutation (HbS: Glu6Val). Polymerization of deoxy-HbS $\rightarrow$ sickling, vaso-occlusion, anemia.
\textbf{Strategy:} Not correcting \emph{HBB} directly (hard in HSPCs). Instead: \textbf{Reactivate fetal hemoglobin (HbF)}. HbF ($\alpha_2\gamma_2$) does not sickle. The \emph{BCL11A} erythroid enhancer represses $\gamma$-globin. Disrupt this enhancer $\rightarrow$ HbF $\uparrow$.
\textbf{Procedure (Ex vivo):}
\begin{enumerate}
    \item Mobilize CD34+ hematopoietic stem/progenitor cells (HSPCs) from patient (G-CSF + plerixafor).
    \item Electroporate Cas9-gRNA RNP targeting \emph{BCL11A} erythroid enhancer.
    \item Myeloablative conditioning (busulfan) to make niche space.
    \item Autologous infusion of edited HSPCs.
\end{enumerate}
\textbf{Results (CLIMB-111/121 trials):} $>90\%$ patients free of vaso-occlusive crises (VOCs) or transfusion independence for $\beta$-thalassemia. HbF $\sim$20–40\%. Approved UK/US/EU 2023–2024. First CRISPR therapy approved.
\textbf{Significance:} Proof that \emph{ex vivo} editing of HSPCs with NHEJ-mediated disruption is curative for hemoglobinopathies.
\end{exemplebox}

\begin{exoresolu}[Designing a CRISPR Strategy for a Hypothetical Mutation]
\textbf{Statement:} A patient has a recessive disease caused by a nonsense mutation (CAG $\rightarrow$ TAG, Gln $\rightarrow$ Stop) in exon 4 of gene \emph{X}. The genomic target sequence (non-target strand) is: 5'-\textbf{GGACTG\textcolor{popred}{TAG}}GGTACC-3' (PAM underlined). Design a gRNA spacer (20 nt) to correct this via HDR using a single-stranded oligodeoxynucleotide (ssODN) donor. Explain the PAM requirement and the donor design.

\textbf{Solution:}
\textbf{1. Identify PAM and Protospacer.} \emph{S. pyogenes} Cas9 requires 5'-NGG-3' on the \textbf{non-target strand}. The sequence shows \textbf{TGG} at positions 10–12 (if we number from 1). Wait, the mutation is at position 7–9 (TAG). We need a PAM \emph{downstream} (3') of the target site on the non-target strand.
Let's scan 3' of the mutation: ...TAG\textbf{GGT}... No GG. ...TAGG\textbf{TA}... No. ...TAGGTA\textbf{CC}... No.
Scan 5' (upstream): \textbf{GGA}CTG... No. \textbf{GAC}TG... No. \textbf{ACT}G... No. \textbf{CTG}... No. \textbf{TGA}... No.
\textbf{Problem:} There may be no suitable NGG PAM near the mutation for SpCas9.
\textbf{Alternative:} Use a different Cas ortholog (e.g. \emph{S. aureus} Cas9: NNGRRT; \emph{Cpf1/Cas12a}: TTTV on 5' end).
\textbf{Assume we find a PAM:} Suppose sequence is 5'-...CTGTAGGGT\textbf{AGG}...-3' (PAM = AGG at end). The protospacer is the 20 nt immediately 5' of PAM: 5'-\textbf{CTGTAGGGTAGGCTGAA???}-3' (need 20 nt total).
\textbf{2. gRNA Spacer Sequence:} The gRNA spacer is identical to the non-target strand (except U for T). If protospacer = 5'-ATGCTGTAGGGTAGGCTGAA-3' (example 20 nt), gRNA = 5'-AUGUUGUAGGGUAGGCUGAA-3'? Wait: DNA: A T G C T G T A G G G T A G G C T G A A. RNA: A U G C U G U A G G G U A G G C U G A A.
\textbf{3. Donor ssODN Design:}
\begin{itemize}
    \item Length: $\sim$100–200 nt total, symmetric homology arms ($\sim$50–90 nt each side).
    \item Correction: Change TAG (Stop) back to CAG (Gln).
    \item \textbf{Silent PAM/Seed mutation:} Introduce a synonymous mutation in the PAM or seed region (e.g. AGG $\rightarrow$ AGA, still Arg) to prevent re-cutting of the corrected allele by Cas9-gRNA.
    \item Strand: Use the strand complementary to the non-target strand (i.e. same as target strand) for better HDR efficiency with Cas9.
\end{itemize}
\textbf{Key Point:} PAM availability dictates targetability. Base editors (CBEs/ABEs) or Prime Editing are often preferred for point mutations as they don't require DSBs or donor templates.
\end{exoresolu}

\courssub{Evaluation of Efficacy and Safety}

Moving from bench to bedside requires rigorous biomarkers.

\begin{itemize}
    \item \textbf{Vector Copy Number (VCN) / Transduction Efficiency:} qPCR or ddPCR on genomic DNA from target tissue (or blood for HSPCs) measuring transgene vs. single-copy reference gene. Target VCN: 1–3 copies/cell for lentiviral; vector genomes (vg) per cell for AAV.
    \item \textbf{Transgene Expression:} mRNA (RT-qPCR, RNA-seq) and protein (ELISA, Western blot, flow cytometry, functional assay e.g. Factor IX activity for hemophilia B).
    \item \textbf{Off-Target Analysis (Editing):} Targeted deep sequencing of predicted off-target sites (in silico prediction + GUIDE-seq). Whole-genome sequencing (WGS) for unbiased detection. Threshold: typically $<0.1\%$ indel frequency at off-targets.
    \item \textbf{Insertional Site Analysis (Integrating vectors):} LAM-PCR or nrLAM-seq to map integration sites. Check for clonal dominance (enrichment near oncogenes like \emph{LMO2}, \emph{CCND2}).
    \item \textbf{Immunogenicity:} Neutralizing antibodies (NAbs) against capsid (AAV) or transgene product (if non-self). T-cell responses to capsid peptides (monitor liver enzymes for AAV-hepatotropism).
\end{itemize}

\courssub{Ethical and Legal Framework}

\begin{savaistu}[Bioethics in Cameroon and Internationally]
\textbf{International:} The \textbf{Declaration of Helsinki} (WMA, 1964, revised 2013) is the cornerstone for human research: informed consent, risk/benefit assessment, independent ethics committee (IEC/IRB) review. The \textbf{UNESCO Universal Declaration on the Human Genome and Human Rights (1997)} states the genome is "heritage of humanity" and opposes germline modifications "contrary to human dignity".
\textbf{Europe/West:} Strict distinction: \textbf{Somatic gene therapy} (treats patient, not inherited) = allowed with oversight. \textbf{Germline editing} (embryos, gametes, heritable) = banned (EU Convention on Human Rights and Biomedicine, Oviedo 1997).
\textbf{Cameroon:} \textbf{Law No. 2002/003 of 19 April 2002} relating to the protection of persons undergoing biomedical research. Requires authorization from the \textbf{National Ethics Committee for Human Health Research (CNERSH)}. Gene therapy trials would fall under this. No specific "bioethics law" on germline editing yet, but the Constitution protects human dignity and the Helsinki principles are ratified. The \textbf{Ministry of Public Health (MINSANTÉ)} oversees clinical trials.
\textbf{He Jiankui Affair (2018):} Chinese scientist created CRISPR-edited babies (CCR5 $\Delta$32 for HIV resistance). Universal condemnation: premature, inadequate consent, off-target risks, germline modification. Highlights need for global governance (WHO Expert Advisory Committee on Developing Global Standards for Governance and Oversight of Human Genome Editing, 2019–2021).
\end{savaistu}

\begin{aretenir}[Key Points for Revision]
\begin{itemize}
\item \textbf{Gene Therapy} = Nucleic acid delivery for therapeutic effect. Two paradigms: \textbf{Addition} (viral/non-viral vectors) vs \textbf{Editing} (CRISPR-Cas9, Base/Prime editors).
\item \textbf{Vectors:} \textbf{AAV} (episomal, non-dividing cells, low immunogenicity, $\sim$4.5 kb limit). \textbf{Lentivirus} (integrating, dividing/non-dividing, $\sim$8 kb, \emph{ex vivo} HSPCs). \textbf{Non-viral/LNP} (scalable, transient, improving).
\item \textbf{CRISPR-Cas9:} gRNA (spacer + scaffold) guides Cas9 to target DNA adjacent to PAM (NGG). DSB $\rightarrow$ NHEJ (indels, knockout) or HDR (precise edit, needs donor, low efficiency).
\item \textbf{Clinical successes:} Luxturna (AAV-\emph{RPE65}, retinal), Zolgensma (AAV9-\emph{SMN1}, SMA), Casgevy (CRISPR-\emph{BCL11A} enhancer, SCD/Thal).
\item \textbf{Safety:} Off-target edits (oncogenic risk), insertional mutagenesis (lentivirus), immune responses (AAV capsid, transgene). Monitoring: VCN, expression, off-target seq, integration site analysis.
\item \textbf{Ethics:} Somatic = therapeutic, regulated. Germline = heritable, internationally prohibited. Cameroon: Law 2002/003, CNERSH oversight, Helsinki Declaration.
\end{itemize}
\end{aretenir}

\begin{exercice}[Gene Therapy Design Challenge]
A 6-year-old Cameroonian boy is diagnosed with X-linked Severe Combined Immunodeficiency (SCID-X1) due to a null mutation in \emph{IL2RG} (common $\gamma$ chain). He has no matched sibling donor for HSCT. You are the molecular biologist designing a gene therapy protocol.
\begin{enumerate}
    \item Choose the optimal vector (AAV, Lentivirus, Adenovirus, LNP) and justify your choice based on: target cell type (HSPCs), need for lifelong correction, transgene size (\emph{IL2RG} cDNA $\sim$1.2 kb), and division status of target cells.
    \item Outline the \emph{ex vivo} manufacturing steps from leukapheresis to cryopreserved drug product.
    \item Identify the major historical safety concern for integrating vectors in SCID-X1 trials (early 2000s) and how modern vectors mitigate it.
    \item Propose a monitoring plan for the first 2 years post-infusion (efficacy + safety biomarkers).
\end{enumerate}
\end{exercice}

\begin{corrige}[Exercise: Gene Therapy Design Challenge]
\begin{enumerate}
    \item \textbf{Vector: Lentivirus (SIN, self-inactivating).} Justification: Target = CD34+ HSPCs (dividing, long-lived). Need stable integration for lifelong multi-lineage reconstitution. Transgene small (fits easily). AAV is episomal $\rightarrow$ lost on HSPC division. Adenovirus immunogenic, transient. LNP DNA delivery to HSPCs inefficient.
    \item \textbf{Steps:} (1) Mobilization (G-CSF/plerixafor) $\rightarrow$ Leukapheresis. (2) CD34+ selection (CliniMACS). (3) Activation/Prestimulation (cytokines: SCF, TPO, Flt3L, IL-6). (4) Transduction (LV vector, MOI 5–20, retronectin/fibronectin). (5) Wash, formulation (cryoprotectant DMSO/albumin). (6) QC: Sterility, VCN (ddPCR), viability, potency (colony-forming units), replication-competent lentivirus (RCL) test. (7) Cryopreservation. (8) Patient conditioning (busulfan). (9) Thaw $\rightarrow$ Infusion.
    \item \textbf{Historical concern:} Insertional oncogenesis. First SCID-X1 trials (gamma-retrovirus, LTR-driven) $\rightarrow$ T-cell leukemia in $\sim$25\% patients due to LTR enhancer activating \emph{LMO2} oncogene. \textbf{Mitigation:} SIN-LTRs (deleted U3 enhancer), internal promoters (e.g. \emph{EFS}, \emph{PGK}) weaker/physiological, chromatin insulators (cHS4), integration-deficient lentiviral vectors (IDLV) for non-dividing targets (not here).
    \item \textbf{Monitoring:} \textbf{Months 1–3:} Engraftment (chimerism), lymphocyte subsets (CD3, CD4, CD8, CD19, NK), Ig levels, VCN in PBMCs (qPCR/ddPCR), integration site analysis (clonal tracking), RCL testing. \textbf{Months 6–24:} Immune reconstitution (vaccine responses, thymic output TRECs), clinical infections, VCN stability, clonal diversity (no dominant clone $>30\%$), oncogene proximity screening. Long-term: Annual follow-up 15 years (FDA/EMA).
\end{enumerate}
\end{corrige}

\coursec{Technologie de l'ADN recombinant et altération génique}

In Section 3 we saw how a defective gene can be corrected or replaced inside a patient --- the logic of gene therapy and genome editing. But before any gene can be corrected, delivered or edited, molecular biologists must first be able to \emph{cut} DNA at precise positions, \emph{paste} fragments together, \emph{copy} a gene millions of times, and \emph{check} that the final construct is exactly what was intended. These four abilities --- cut, paste, amplify, verify --- are the heart of \cle{recombinant DNA technology}. This section presents the enzymatic toolbox (restriction enzymes, ligases, polymerases), the vehicles that carry foreign DNA into cells (cloning vectors), the classical cloning workflow, the polymerase chain reaction (PCR), modern editing tools such as CRISPR--Cas9, and the quality-control methods that validate every construct.

\courssub{Restriction enzymes: the molecular scissors}

\textbf{Observation.} Bacteria are constantly attacked by viruses (bacteriophages). Many bacteria defend themselves by producing enzymes that chop the invading viral DNA into pieces while leaving their own DNA untouched (their own recognition sites are protected by methylation). These enzymes, purified and used in the laboratory, became the first great tool of genetic engineering.

\begin{definition}[Recombinant DNA]
\cle{Recombinant DNA} is a DNA molecule formed by joining together fragments of DNA from different sources (for example a human insulin gene inserted into a bacterial plasmid). Its construction requires enzymes able to cut DNA at defined sequences and to reseal the sugar--phosphate backbone.
\end{definition}

\begin{definition}[Restriction endonuclease]
A \cle{restriction enzyme} (restriction endonuclease) is an enzyme that recognises a specific short DNA sequence, the \cle{recognition site}, and cuts both strands of the double helix. Type II enzymes, the most useful in the laboratory, cut \emph{within or immediately next to} their recognition site.
\end{definition}

\textbf{Nomenclature.} The name encodes the organism of origin: the first letter of the genus, the first two letters of the species, the strain, and a Roman numeral for the order of discovery. Thus \cle{EcoRI} comes from \emph{Escherichia coli}, strain R, enzyme I; \emph{Hin}dIII from \emph{Haemophilus influenzae} strain d; \emph{Bam}HI from \emph{Bacillus amyloliquefaciens} strain H.

\textbf{Palindromic recognition sites.} Most recognition sites are \cle{palindromes}: the sequence reads identically in the $5' \to 3'$ direction on both strands. For example EcoRI recognises:
\[
\begin{array}{ccccccc}
5' & G & A & A & T & T & C\ 3'\\
3' & C & T & T & A & A & G\ 5'
\end{array}
\]
Reading the top strand $5'\to3'$ gives GAATTC; reading the bottom strand $5'\to3'$ also gives GAATTC. This symmetry matches the dimeric structure of the enzyme.

\textbf{Sticky versus blunt ends.} Depending on where the two strands are cut, two kinds of ends are produced:
\begin{itemize}
\item \cle{Cohesive (sticky) ends}: the cuts are staggered, leaving short single-stranded overhangs (EcoRI leaves a $5'$--AATT overhang). Two fragments cut with the same enzyme have complementary overhangs and can pair by hydrogen bonding --- this makes ligation efficient and directional cloning possible when two different enzymes are used.
\item \cle{Blunt ends}: both strands are cut at the same position (e.g.\ \emph{Sma}I, CCC$\downarrow$GGG). Any blunt end can join any other blunt end, but ligation is far less efficient and orientation cannot be controlled.
\end{itemize}

\begin{propriete}[Type II restriction enzymes]
Type II restriction endonucleases cut DNA within or very close to their recognition sequence, generating predictable, reproducible ends (cohesive or blunt). Consequently, digestion of a given DNA molecule by a given enzyme always yields the same set of fragments --- the basis of restriction mapping and cloning. (Statement examinable; no formal proof required.)
\end{propriete}

\begin{attention}[Star activity and methylation]
Under non-optimal conditions (high glycerol, wrong pH or ionic strength), some enzymes show ``star activity'' and cut at sites resembling the true site, giving extra bands on gels. Also, DNA extracted from dam$^+$ \emph{E. coli} strains is methylated and may resist enzymes such as \emph{Cla}I or \emph{Xba}I. Always check the enzyme data sheet before designing a digest.
\end{attention}

\courssub{Cloning vectors: vehicles for foreign DNA}

A restriction fragment alone cannot replicate inside a cell. It must be joined to a \cle{cloning vector}: a small, autonomously replicating DNA molecule that carries the insert into a host cell and amplifies it.

\begin{definition}[Cloning vector]
A \cle{cloning vector} is a DNA molecule (plasmid, phage-derived or artificial chromosome) capable of independent replication in a host cell, used to carry and amplify a foreign DNA insert. Every useful vector possesses three essential elements:
\begin{enumerate}
\item an \cle{origin of replication} (ori), recognised by the host replication machinery;
\item a \cle{selectable marker} (e.g.\ an antibiotic-resistance gene such as \emph{amp}$^R$) allowing survival of only those cells that received the vector;
\item a \cle{multiple cloning site} (MCS or polylinker): a short region containing many unique restriction sites where the insert is placed.
\end{enumerate}
\end{definition}

\begin{center}
\begin{tabularx}{\linewidth}{|l|X|X|}
\hline
\rowcolor{popblueL} \textbf{Vector} & \textbf{Insert capacity} & \textbf{Typical use} \\
\hline
Plasmid (pUC19, pBR322) & up to $\sim 10$ kb & Routine cloning, subcloning, expression \\
\hline
Cosmid (plasmid + phage $\lambda$ \emph{cos} sites) & $35$--$45$ kb & Genomic libraries of moderate-size genomes \\
\hline
BAC (bacterial artificial chromosome) & $100$--$300$ kb & Large genomic fragments, genome projects \\
\hline
YAC (yeast artificial chromosome) & $250$ kb -- $1$ Mb & Very large fragments; contains yeast ori, centromere and telomeres \\
\hline
\end{tabularx}
\end{center}

The plasmid \cle{pUC19} is the classic teaching example: it carries the \emph{ori}, the \emph{amp}$^R$ gene (ampicillin resistance), and the MCS inserted inside the \emph{lacZ}$\alpha$ gene --- an arrangement exploited by blue/white screening (Section 4.3).

\begin{popfigure}
\begin{tikzpicture}[scale=1.0]
% plasmid circle
\draw[thick, popblue] (0,0) circle (2.6);
\draw[thick, popblueL, line width=4pt] (90:2.6) arc (90:150:2.6);
\draw[thick, popgreenL, line width=4pt] (150:2.6) arc (150:235:2.6);
\draw[thick, poporange!60!white, line width=4pt] (235:2.6) arc (235:330:2.6);
\draw[thick, poppink!70!white, line width=4pt] (330:2.6) arc (330:90:2.6);
% labels
\node[align=center, font=\small] at (120:3.6) {\textbf{ori}};
\draw[->, popmuted] (120:3.3) -- (120:2.7);
\node[align=center, font=\small] at (192:3.9) {\textbf{\emph{amp}$^R$}};
\draw[->, popmuted] (192:3.3) -- (192:2.7);
\node[align=center, font=\small] at (282:3.9) {\textbf{\emph{lacZ}$\alpha$}};
\draw[->, popmuted] (282:3.3) -- (282:2.7);
\node[align=center, font=\small] at (30:3.9) {\textbf{MCS inside \emph{lacZ}$\alpha$}};
\draw[->, popmuted] (30:3.3) -- (30:2.7);
% MCS restriction sites
\node[font=\footnotesize, align=left, popdark] at (5.6,1.6) {MCS sites (5'$\to$3'):};
\node[font=\footnotesize, align=left, popdark] at (5.6,1.1) {\emph{Eco}RI -- \emph{Sac}I -- \emph{Kpn}I};
\node[font=\footnotesize, align=left, popdark] at (5.6,0.6) {\emph{Sma}I -- \emph{Bam}HI -- \emph{Xba}I};
\node[font=\footnotesize, align=left, popdark] at (5.6,0.1) {\emph{Hin}dIII -- \emph{Pst}I -- \emph{Sph}I};
\node[font=\small, popblue] at (0,0) {\textbf{pUC19}};
\node[font=\footnotesize, popmuted] at (0,-0.5) {2686 bp};
\end{tikzpicture}
\legende{Simplified map of the cloning plasmid pUC19: origin of replication (ori), ampicillin-resistance marker (\emph{amp}$^R$), and the multiple cloning site (MCS) embedded in the \emph{lacZ}$\alpha$ reporter gene.}
\end{popfigure}

\courssub{The classical cloning workflow}

\begin{methode}[Classical gene cloning: the five steps]
\begin{enumerate}
\item \textbf{Digestion.} Cut both the vector and the source DNA with the same restriction enzyme(s) so that compatible ends are generated. Dephosphorylate the vector ends (alkaline phosphatase) to prevent self-ligation.
\item \textbf{Ligation.} Mix insert and vector with T4 DNA ligase, which reseals the sugar--phosphate backbone (see the ligation protocol box below).
\item \textbf{Transformation.} Introduce the recombinant plasmid into competent \emph{E. coli} cells, either by \cle{heat shock} (cells made competent with CaCl$_2$, chilled on ice, then pulsed at $42\ ^\circ$C for $\sim 45$ s) or by \cle{electroporation} (a brief electric pulse opens transient pores in the membrane; higher efficiency).
\item \textbf{Selection.} Spread the cells on agar containing ampicillin: only bacteria carrying a plasmid (with \emph{amp}$^R$) form colonies.
\item \textbf{Screening.} Distinguish recombinant from non-recombinant plasmids:
\begin{itemize}
\item \cle{Blue/white screening}: the MCS lies inside \emph{lacZ}$\alpha$. An intact \emph{lacZ}$\alpha$ produces $\beta$-galactosidase, which cleaves X-gal into a blue product. An insert disrupts \emph{lacZ}$\alpha$, so recombinant colonies stay \textbf{white} while empty-vector colonies turn \textbf{blue}.
\item \cle{Colony PCR}: a tiny amount of each colony is used directly as PCR template with primers flanking the MCS; the size of the product reveals the presence and size of the insert.
\end{itemize}
\end{enumerate}
\end{methode}

\begin{methode}[Ligation protocol]
\begin{enumerate}
\item Prepare the reaction with an \cle{insert : vector molar ratio of 3:1} (excess insert favours productive joining over vector self-circularisation). Compute masses with
\[
m_{\text{insert}} = 3 \times m_{\text{vector}} \times \frac{\text{size of insert (bp)}}{\text{size of vector (bp)}} .
\]
\item Add T4 DNA ligase and its ATP-containing buffer.
\item Incubate at $16\ ^\circ$C overnight (a compromise: ends anneal best at low temperature, but the ligase works faster at higher temperature).
\item Transform $1$--$2\ \mu$L of the ligation mixture into competent cells.
\end{enumerate}
\end{methode}

\begin{exoresolu}[Calculating insert mass for a ligation]
We wish to clone a $750$ bp insert into plasmid pUC19 ($2686$ bp). We use $50$ ng of vector. What mass of insert gives a 3:1 molar ratio?
\tcblower
\textbf{Solution.} Equal numbers of molecules require masses proportional to length. For a 3:1 insert:vector molar ratio:
\[
m_{\text{insert}} = 3 \times 50\ \text{ng} \times \frac{750}{2686} \approx 150 \times 0.279 \approx 42\ \text{ng}.
\]
We therefore mix about $42$ ng of insert with $50$ ng of cut, dephosphorylated vector before adding T4 DNA ligase.
\end{exoresolu}

\begin{attention}[Antibiotic-resistance markers and biosafety]
Resistance markers such as \emph{amp}$^R$ are convenient, but their release into the environment could spread resistance genes to pathogenic bacteria --- a real public-health concern, including in Cameroon where antibiotic resistance is already a serious problem. Good practice: never discard recombinant cultures untreated (autoclave or bleach them), and prefer alternative selection systems where possible --- \cle{auxotrophic complementation} (the vector restores a metabolic gene the host lacks), \cle{fluorescent markers} (GFP) or toxin--antitoxin systems.
\end{attention}

\courssub{DNA amplification by PCR}

Cloning amplifies DNA \emph{inside} cells. The \cle{polymerase chain reaction} (PCR), conceived by Kary Mullis in the 1980s, amplifies a chosen DNA region \emph{in vitro}, in a tube, in a few hours.

\begin{definition}[PCR]
\cle{PCR} is a cyclic in vitro reaction that doubles a target DNA sequence at each cycle, using: (i) a \textbf{template} DNA; (ii) two short \cle{primers} (amorces) flanking the target, one per strand; (iii) the four \textbf{dNTPs}; (iv) a \cle{thermostable DNA polymerase} (Taq polymerase, from the hot-spring bacterium \emph{Thermus aquaticus}) that survives repeated heating; (v) a buffer containing Mg$^{2+}$.
\end{definition}

Each cycle has three steps:
\begin{enumerate}
\item \textbf{Denaturation} ($94$--$98\ ^\circ$C, $20$--$30$ s): the double helix separates into single strands.
\item \textbf{Annealing} (typically $50$--$65\ ^\circ$C, $20$--$40$ s, depending on the primers): primers hybridise to their complementary sequences.
\item \textbf{Extension (elongation)} ($72\ ^\circ$C, $\sim 1$ min per kb): Taq polymerase synthesises the new strands from the primers in the $5'\to3'$ direction.
\end{enumerate}
After $n$ cycles the target is amplified theoretically $2^n$-fold: 30 cycles give $\sim 10^9$ copies.

\begin{popfigure}
\begin{tikzpicture}[scale=1.0]
\draw[->, thick, popmuted] (0,0) -- (10.6,0) node[below left, font=\footnotesize]{time (min)};
\draw[->, thick, popmuted] (0,0) -- (0,5.6) node[above left, font=\footnotesize, align=center]{T ($^\circ$C)};
% temperature levels
\draw[dashed, popline] (0,4.8) -- (10.4,4.8) node[right, font=\footnotesize, popdark]{94};
\draw[dashed, popline] (0,3.6) -- (10.4,3.6) node[right, font=\footnotesize, popdark]{72};
\draw[dashed, popline] (0,1.6) -- (10.4,1.6) node[right, font=\footnotesize, popdark]{55};
% initial denaturation
\draw[very thick, popblue] (0,0.4) -- (0,4.8) -- (1.2,4.8);
% cycle 1
\draw[very thick, popblue] (1.2,4.8) -- (1.5,1.6) -- (2.3,1.6) -- (2.6,3.6) -- (3.9,3.6) -- (4.2,4.8) -- (4.8,4.8);
% cycle 2
\draw[very thick, popblue] (4.8,4.8) -- (5.1,1.6) -- (5.9,1.6) -- (6.2,3.6) -- (7.5,3.6) -- (7.8,4.8) -- (8.4,4.8);
% cycle 3 (partial)
\draw[very thick, popblue] (8.4,4.8) -- (8.7,1.6) -- (9.5,1.6) -- (9.8,3.6) -- (10.4,3.6);
% labels
\node[font=\footnotesize, popdark, align=center] at (0.6,5.2) {initial\\ denat.};
\node[font=\footnotesize, popdark, align=center] at (1.9,1.1) {annealing\\ 30 s};
\node[font=\footnotesize, popdark, align=center] at (3.25,3.1) {extension\\ 60 s};
\node[font=\footnotesize, popdark, align=center] at (4.5,4.3) {denat.\\ 25 s};
\node[font=\footnotesize, popmuted] at (9.2,5.2) {$\times$ 30 cycles};
% braces for one cycle
\draw[decorate, decoration={brace, amplitude=4pt}, poporange] (4.2,-0.3) -- (1.2,-0.3) node[midway, below=6pt, font=\footnotesize, popdark]{one cycle};
\end{tikzpicture}
\legende{Temperature chronogram of a typical PCR programme: initial denaturation, then repeated cycles of denaturation ($94\ ^\circ$C), primer annealing ($55\ ^\circ$C) and extension by Taq polymerase ($72\ ^\circ$C).}
\end{popfigure}

\textbf{Applications.} PCR is used for \cle{genotyping} (e.g.\ detecting the sickle-cell allele Hb$^S$ from a drop of blood), diagnosis of infections (malaria parasites, HIV, SARS-CoV-2), forensic identification, and cloning. In \cle{quantitative PCR (qPCR)}, a fluorescent dye or probe reports product accumulation after each cycle: the earlier the fluorescence crosses a threshold (the $C_q$ value), the more target was initially present --- allowing precise quantification of viral load or gene expression.

\begin{exoresolu}[How many copies after 25 cycles?]
A PCR starts from a single DNA molecule. Assuming perfect efficiency, how many copies of the target exist after 25 cycles? If the efficiency is only $90\%$ (each cycle multiplies by $1.9$), how many copies result?
\tcblower
\textbf{Solution.} At $100\%$ efficiency, $N = 2^{25} = 33\,554\,432 \approx 3.4\times 10^{7}$ copies.\\
At $90\%$ efficiency, $N = 1.9^{25}$. Taking logarithms: $\log_{10} N = 25 \times \log_{10} 1.9 = 25 \times 0.2788 = 6.97$, so $N \approx 10^{6.97} \approx 9.3 \times 10^{6}$ copies.\\
\textbf{Comment.} A $10\%$ drop in efficiency reduces the final yield almost four-fold over 25 cycles --- this is why primer design and Mg$^{2+}$ optimisation matter, and why qPCR efficiency must be measured before quantification.
\end{exoresolu}

\courssub{Modern genome-editing tools}

PCR and cloning manipulate DNA in a tube; \cle{genome editing} modifies a chosen sequence \emph{inside} the living genome.

\textbf{CRISPR--Cas9.} Derived from a bacterial immune system, the \cle{CRISPR--Cas9} complex uses a \cle{guide RNA (gRNA)} of about 20 nucleotides, complementary to the target, to lead the Cas9 nuclease to the matching genomic site, which must sit next to a short \cle{PAM} motif (NGG for the \emph{Streptococcus pyogenes} enzyme). Cas9 cuts both strands; the cell then repairs the break either by error-prone non-homologous end joining (gene knockout) or, if a donor template is provided, by homology-directed repair (precise correction --- the strategy behind the sickle-cell therapies discussed in Section 3).

\textbf{Base editors and prime editing.} \cle{Base editors} fuse a catalytically impaired Cas9 to a deaminase and convert one base into another (C$\to$T or A$\to$G) \emph{without} cutting both strands. \cle{Prime editing} uses a Cas9 nickase fused to a reverse transcriptase and an extended guide RNA (pegRNA) that both specifies the target and encodes the desired edit, allowing virtually any small substitution, insertion or deletion with minimal off-target damage.

\begin{methode}[Designing a good guide RNA]
\begin{enumerate}
\item Choose a $20$-nt target sequence adjacent to a PAM (NGG) within the region to be edited.
\item Check specificity with bioinformatics tools such as \cle{CRISPOR} or \cle{Benchling}: reject any gRNA whose sequence shows more than $\sim 80\%$ homology with another genomic site, especially sites differing by fewer than 3 mismatches.
\item Prefer guides with a GC content of $40$--$60\%$ and avoid long poly-T stretches (which can terminate transcription).
\item Order the top-scoring candidates and validate them experimentally (e.g.\ by a T7E1 assay or targeted sequencing) before committing to the full experiment.
\end{enumerate}
\end{methode}

\courssub{Quality control of constructs}

No construct is trusted until it has been verified by independent methods:

\begin{itemize}
\item \cle{Agarose gel electrophoresis}: DNA fragments migrate towards the positive electrode at a speed inversely related to their size; staining (ethidium bromide or safer dyes) reveals bands whose sizes are compared with a molecular-weight ladder. A quick check of digestions, ligations and PCR products.
\item \cle{Restriction analysis}: the purified plasmid is re-digested with diagnostic enzymes; the observed band pattern must match the pattern predicted from the map.
\item \cle{Sanger sequencing}: the gold standard. A primer anneals to the construct and a polymerase extends it in the presence of dNTPs plus small amounts of chain-terminating dideoxynucleotides, each labelled with a different fluorescent colour. The resulting fragments, separated by capillary electrophoresis, give the exact sequence, base by base, confirming that the insert is present, correctly oriented and free of mutations.
\end{itemize}

\begin{aretenir}[Key points of Section 4]
\begin{itemize}
\item Type II restriction enzymes cut at palindromic sites, giving predictable sticky or blunt ends; T4 DNA ligase reseals compatible ends (insert:vector $= 3:1$, $16\ ^\circ$C overnight).
\item A cloning vector needs an \emph{ori}, a selectable marker and an MCS; capacity increases from plasmids to cosmids, BACs and YACs.
\item Cloning workflow: digestion $\to$ ligation $\to$ transformation (heat shock or electroporation) $\to$ antibiotic selection $\to$ screening (blue/white, colony PCR).
\item PCR amplifies a target exponentially ($2^n$) using primers and a thermostable polymerase; qPCR quantifies the starting amount.
\item CRISPR--Cas9, base editors and prime editing modify genomes at chosen sites; gRNAs must be designed with tools such as CRISPOR or Benchling to avoid off-targets ($>80\%$ homology rejected).
\item Every construct is validated by gel electrophoresis, restriction analysis and Sanger sequencing.
\end{itemize}
\end{aretenir}

\begin{exercice}[Restriction mapping]
A linear DNA fragment of $6000$ bp is digested with EcoRI alone, giving fragments of $4000$ bp and $2000$ bp. Digestion with BamHI alone gives $4500$ bp and $1500$ bp. The double digest EcoRI + BamHI gives three fragments: $2000$, $2500$ and $1500$ bp. Draw the restriction map showing the relative positions of the EcoRI (E) and BamHI (B) sites.
\end{exercice}

\begin{corrige}[Restriction mapping]
Let the left end be position 0 and the right end 6000 bp.
EcoRI alone gives 4000 and 2000 bp $\Rightarrow$ EcoRI site (E) is at 2000 bp or 4000 bp from the left.
BamHI alone gives 4500 and 1500 bp $\Rightarrow$ BamHI site (B) is at 1500 bp or 4500 bp from the left.
Double digest gives three fragments: 2000, 2500, 1500 bp (sum = 6000). These are the intervals between 0, E, B, 6000.

Test the four combinations:
\begin{itemize}
\item E=2000, B=1500: intervals 1500, 500, 4000 $\neq$ given.
\item E=2000, B=4500: intervals 2000, 2500, 1500 $\checkmark$ matches double digest.
\item E=4000, B=1500: intervals 1500, 2500, 2000 $\checkmark$ same set, but order differs.
\item E=4000, B=4500: intervals 4000, 500, 1500 $\neq$ given.
\end{itemize}
Both (E=2000, B=4500) and (E=4000, B=1500) produce the same three fragment sizes. However, the single digests distinguish them:
\begin{itemize}
\item If E=2000, B=4500: EcoRI fragments = 2000 and 4000 $\checkmark$; BamHI fragments = 4500 and 1500 $\checkmark$.
\item If E=4000, B=1500: EcoRI fragments = 4000 and 2000 $\checkmark$; BamHI fragments = 1500 and 4500 $\checkmark$.
\end{itemize}
Both are consistent with all data; the map is symmetric. We can draw either orientation. Conventionally, we place the smaller coordinate first:
\[
0 \ \xrightarrow{\ 2000\ }\ \mathrm{E}\ \xrightarrow{\ 2500\ }\ \mathrm{B}\ \xrightarrow{\ 1500\ }\ 6000
\]
Check: EcoRI at 2000 $\to$ 2000 + 4000. BamHI at 4500 $\to$ 4500 + 1500. Double digest: 0--E = 2000, E--B = 2500, B--6000 = 1500. All consistent.
\end{corrige}

\coursec{Applications majeures de la biotechnologie : santé, agriculture, environnement, industrie}

In the previous sections we mastered the \emph{tools} of modern biotechnology: restriction enzymes, vectors, PCR, recombinant DNA, and gene editing with CRISPR--Cas9. A tool, however, is only as valuable as what it builds. In this final section we answer the question every examiner loves: \emph{what has biotechnology actually done for humanity?} We shall tour the four great fields of application --- \cle{health}, \cle{agriculture}, \cle{environment} and \cle{industry/energy} --- and finish with the regulatory framework that governs all of this in Cameroon.

\courssub{Pharmaceutical applications: medicines made by living cells}

\textbf{The intuition.} A diabetic patient needs insulin every day. Before 1982, insulin was extracted from the pancreases of pigs and cattle slaughtered in abattoirs: it was expensive, scarce, and sometimes triggered allergic reactions because animal insulin differs slightly from human insulin in its amino-acid sequence. The biotechnological idea is disarmingly simple: \emph{why not give the human insulin gene to a microbe, and let the microbe --- which doubles every 20 minutes --- manufacture the hormone for us?}

\begin{exemplebox}[Recombinant human insulin]
The gene coding for human insulin is inserted into a plasmid and introduced into \emph{Escherichia coli} (or the yeast \emph{Saccharomyces cerevisiae}). The bacteria are grown in huge fermenters; they produce the insulin chains, which are extracted, folded, purified and formulated. Since 1982 (first commercial product: Humulin), virtually all insulin used worldwide is recombinant human insulin: cheaper, unlimited in supply, and immunologically identical to the patient's own hormone.
\end{exemplebox}

The same logic gives us \cle{recombinant human growth hormone} (somatotropin), used to treat growth hormone deficiency in children. Previously it had to be extracted from the pituitary glands of human cadavers --- a practice abandoned after it transmitted a fatal prion disease (Creutzfeldt--Jakob disease) to some treated patients. Recombinant production eliminated this risk entirely. Other recombinant pharmaceuticals include \cle{factor VIII} for haemophiliacs (formerly prepared from pooled blood, with a risk of HIV and hepatitis transmission), \cle{erythropoietin} (EPO) for anaemia, and \cle{hepatitis B vaccine}, produced as a recombinant surface antigen in yeast.

\begin{popfigure}
\begin{tikzpicture}[
  box/.style={draw=popblue, thick, rounded corners=2pt, fill=popblueL, text width=2.6cm, align=center, minimum height=1.1cm, font=\small},
  arr/.style={-{Stealth[length=2.5mm]}, very thick, popdark}]
\node[box] (g) at (0,0) {Human insulin gene isolated or synthesised};
\node[box, right=0.9cm of g] (p) {Gene inserted into plasmid vector};
\node[box, right=0.9cm of p] (e) {Transformation of \emph{E. coli}; fermentation};
\node[box, below=1.0cm of e] (ib) {Inclusion bodies recovered (insoluble protein)};
\node[box, left=0.9cm of ib] (rf) {Refolding and chain assembly};
\node[box, left=0.9cm of rf] (pu) {Chromatographic purification, formulation};
\draw[arr] (g) -- (p);
\draw[arr] (p) -- (e);
\draw[arr] (e) -- (ib);
\draw[arr] (ib) -- (rf);
\draw[arr] (rf) -- (pu);
\end{tikzpicture}
\legende{Production flow of recombinant human insulin: from the human gene to the purified pharmaceutical product.}
\end{popfigure}

\textbf{Monoclonal antibodies.} A \cle{monoclonal antibody} (mAb) is a single, identical antibody species directed against one precise epitope, produced industrially. Classical production uses \emph{hybridoma} cells, obtained by fusing an antibody-producing B lymphocyte with an immortal myeloma (cancer) cell: the hybrid combines specific antibody production with unlimited division. Modern industrial production uses genetically engineered \cle{CHO cells} (Chinese Hamster Ovary cells), which perform correct human-like protein folding and glycosylation. mAbs treat cancers (e.g.\ trastuzumab against certain breast cancers), autoimmune diseases, and serve in diagnostics. A promising low-cost alternative is the production of antibodies in \emph{transgenic plants} (``plantibodies''), for example in tobacco.

\begin{propriete}[mRNA vaccines]
An \cle{mRNA vaccine} delivers, inside \cle{lipid nanoparticles}, a messenger RNA molecule that codes for a pathogen's antigen (e.g.\ the spike protein of SARS-CoV-2). The patient's own cells translate this mRNA and display the antigen, triggering immunity. The mRNA works in the \emph{cytoplasm}, is rapidly degraded, and \textbf{never enters the nucleus nor integrates into the host genome}. Advantages: extremely rapid design (days once the sequence is known) and no risk of infection since no live pathogen is used.
\end{propriete}

\begin{attention}[A common misconception]
mRNA vaccines do \textbf{not} modify your DNA. mRNA cannot integrate into the genome: it lacks the enzymes (integrase, reverse transcriptase) required, and it never crosses the nuclear envelope. Confusing mRNA vaccines with gene therapy is a classic exam error.
\end{attention}

\courssub{Agricultural applications: GM crops and gene editing}

\begin{definition}[GMO]
A \cle{Genetically Modified Organism} (GMO) is an organism whose genetic material has been modified by genetic engineering techniques (introduction, deletion or targeted modification of one or more genes), in a way that does not occur naturally by mating or natural recombination.
\end{definition}

Three flagship examples dominate world agriculture:
\begin{itemize}
\item \cle{Bt cotton} and Bt maize carry a gene from the bacterium \emph{Bacillus thuringiensis} coding for a protein toxic to insect larvae (bollworms, stem borers) but harmless to humans. Result: drastic reduction of insecticide spraying --- a major benefit for smallholder farmers, including in West and Central Africa.
\item \cle{Herbicide-tolerant maize and soybean} carry a gene (e.g.\ a modified EPSPS enzyme) allowing them to survive glyphosate, so weeds can be controlled without killing the crop.
\item \cle{Golden Rice} is engineered to produce $\beta$-carotene (provitamin A) in the grain, to fight vitamin A deficiency, which causes blindness in hundreds of thousands of children yearly in developing countries.
\end{itemize}

Beyond transgenesis, \cle{gene editing} (CRISPR--Cas9) now allows precise modification of a plant's \emph{own} genes, without introducing foreign DNA. For example, editing rice genes that control stomatal closure in response to the stress hormone ABA (abscisic acid) improves \cle{drought tolerance} --- a crucial trait as rainfall becomes erratic in the Sudano-Sahelian zone of northern Cameroon.

\begin{popfigure}
\begin{tikzpicture}[scale=0.9]
% Simplified world map: schematic continents
\fill[popgreenL, draw=popdark, thick] (-5.2,0.2) .. controls (-5.6,1.6) and (-3.4,2.2) .. (-2.6,1.4) .. controls (-2.0,0.8) and (-2.6,-0.4) .. (-3.6,-0.6) .. controls (-4.6,-0.8) and (-4.9,-0.6) .. (-5.2,0.2); % N America
\fill[popgreenL, draw=popdark, thick] (-3.4,-1.0) .. controls (-2.8,-0.8) and (-2.2,-1.4) .. (-2.6,-2.4) .. controls (-3.0,-3.2) and (-3.8,-3.0) .. (-3.9,-2.2) .. controls (-4.0,-1.4) and (-3.8,-1.1) .. (-3.4,-1.0); % S America
\fill[popblueL, draw=popdark, thick] (-0.6,1.0) .. controls (0.4,1.8) and (1.6,1.6) .. (1.4,0.8) .. controls (1.2,0.2) and (0.0,0.2) .. (-0.6,1.0); % Europe
\fill[popgold!40, draw=popdark, thick] (-0.4,0.2) .. controls (0.8,0.4) and (1.4,-0.4) .. (1.2,-1.6) .. controls (1.0,-2.6) and (0.2,-3.0) .. (-0.2,-2.2) .. controls (-0.6,-1.2) and (-1.0,-0.4) .. (-0.4,0.2); % Africa
\fill[popgreenL, draw=popdark, thick] (1.8,0.4) .. controls (3.0,1.8) and (5.2,1.8) .. (5.6,0.6) .. controls (5.9,-0.3) and (4.6,-0.9) .. (3.6,-0.6) .. controls (2.6,-0.3) and (1.4,-0.4) .. (1.8,0.4); % Asia
\fill[popblueL, draw=popdark, thick] (4.4,-2.2) .. controls (5.4,-1.8) and (5.8,-2.8) .. (5.0,-3.2) .. controls (4.2,-3.4) and (3.9,-2.6) .. (4.4,-2.2); % Australia
\node[font=\scriptsize, align=center] at (-3.9,1.0) {USA, Canada\\major GM\\adopters};
\node[font=\scriptsize, align=center] at (-3.2,-2.1) {Brazil,\\Argentina};
\node[font=\scriptsize, align=center] at (3.7,0.9) {India, China\\(Bt cotton)};
\node[font=\scriptsize, align=center] at (0.5,-1.2) {Africa:\\limited,\\growing};
\node[font=\scriptsize, align=center] at (0.5,1.1) {EU:\\restricted};
% Legend
\fill[popgreenL, draw=popdark] (-5.4,-3.6) rectangle (-5.0,-3.3); \node[font=\scriptsize, right] at (-5.0,-3.45) {Large-scale GM cultivation};
\fill[popgold!40, draw=popdark] (-1.6,-3.6) rectangle (-1.2,-3.3); \node[font=\scriptsize, right] at (-1.2,-3.45) {Limited / field trials};
\fill[popblueL, draw=popdark] (1.6,-3.6) rectangle (2.0,-3.3); \node[font=\scriptsize, right] at (2.0,-3.45) {Restricted / import only};
\end{tikzpicture}
\legende{Simplified world map of GM crop adoption. The Americas dominate cultivation; Africa (including Cameroon) proceeds cautiously under biosafety law.}
\end{popfigure}

\begin{methode}[Risk assessment of a GMO]
Before any release, a GMO is evaluated in three steps:
\begin{enumerate}
\item \textbf{Comparative analysis} (principle of \emph{substantial equivalence}): compare the GM plant with its conventional counterpart --- composition, nutrients, toxins, allergens. Any significant difference triggers deeper study.
\item \textbf{Environmental impact study}: gene flow to wild relatives, effects on non-target organisms (e.g.\ pollinators), persistence in soil.
\item \textbf{Management plan}: monitoring after release, refuge zones to delay resistance, traceability and labelling.
\end{enumerate}
\end{methode}

\begin{attention}[Resistance evolution]
Widespread planting of Bt crops exerts strong \emph{selection pressure}: rare resistant caterpillars survive, reproduce, and can produce fully resistant populations within years. The same applies to weeds exposed repeatedly to one herbicide. This is \emph{evolution by natural selection} in action --- the remedy is the \emph{refuge strategy} (non-Bt plots where susceptible insects survive, diluting resistance genes). Never write that ``the insecticide creates the mutation'': mutations arise randomly; the insecticide only \emph{selects} them.
\end{attention}

\courssub{Environmental applications: cleaning up with life}

\begin{definition}[Bioremediation]
\cle{Bioremediation} is the use of living organisms --- micro-organisms (bacteria, fungi) or plants --- to degrade, transform or immobilise pollutants in contaminated soils, water or sediments, restoring the environment.
\end{definition}

\textbf{How it works.} Many bacteria possess enzymes that oxidise hydrocarbons and use them as carbon sources. After an oil spill, one can either seed specialised bacteria (\emph{bioaugmentation}) or add nutrients (nitrogen, phosphorus) to stimulate indigenous oil-degrading bacteria such as \emph{Pseudomonas} species (\emph{biostimulation}). Other micro-organisms transform toxic heavy metals: for example, certain bacteria reduce the highly toxic and mobile Cr(VI) to the less toxic, less mobile Cr(III).

\textbf{Phytoremediation.} Plants can also clean up: \emph{hyperaccumulator} plants (e.g.\ some \emph{Brassica} species) absorb heavy metals such as cadmium, zinc or nickel from contaminated soils into their harvestable shoots; the plants are then collected and treated. This is slow but cheap and aesthetically acceptable --- relevant for mining areas in Cameroon.

\textbf{Bioindicators.} Living organisms also \emph{reveal} pollution: the absence of sensitive aquatic larvae (e.g.\ mayflies) in a stream, or the accumulation of pollutants in lichens, signals environmental degradation before chemical instruments detect it.

\courssub{Energy applications: biofuels of the new generations}

First-generation biofuels (ethanol from sugarcane or maize starch) compete with food crops. Biotechnology now targets:
\begin{itemize}
\item \cle{Second-generation bioethanol}: lignocellulosic biomass (straw, bagasse, wood waste) is broken down by engineered enzymes (cellulases) into fermentable sugars --- no competition with food.
\item \cle{Biodiesel from microalgae}: microalgae can accumulate oils (up to about 50\% of dry mass in some species) and grow on non-arable land with far higher yields per hectare than oil palms.
\item \cle{Biohydrogen}: certain algae and bacteria produce \ce{H2} by photosynthesis or fermentation, a clean fuel whose combustion releases only water.
\end{itemize}

\courssub{Diagnostic applications: seeing the invisible}

\begin{itemize}
\item \cle{Real-time PCR (qPCR)}: amplifies and quantifies pathogen DNA/RNA as the reaction proceeds (fluorescence measured each cycle). It detects viruses (HIV viral load, SARS-CoV-2) with extreme sensitivity.
\item \cle{ELISA}: enzyme-linked immunosorbent assay; an antigen--antibody reaction coupled to an enzyme producing a coloured signal proportional to the target (used for HIV screening).
\item \cle{Lateral flow rapid tests}: the principle of pregnancy tests and of the malaria RDTs used daily in Cameroonian health centres --- a drop of blood migrates along a strip; a coloured line appears if the target antigen (e.g.\ \emph{Plasmodium} proteins) is present. Result in about 15 minutes, no laboratory needed.
\item \cle{Next-generation sequencing (NGS)}: millions of DNA fragments sequenced simultaneously, allowing identification of new pathogens, tracking of variants, and personalised medicine.
\end{itemize}

\begin{exoresolu}[Interpreting a qPCR result]
In a real-time PCR test for a virus, the $C_t$ (cycle threshold) is the cycle number at which fluorescence exceeds the background. Sample A gives $C_t = 20$, sample B gives $C_t = 30$. Each cycle doubles the DNA. Which patient has the higher initial viral load, and by approximately what factor?
\tcblower
\textbf{Solution.} A \emph{lower} $C_t$ means the signal was reached earlier, so the starting material was \emph{more} abundant: patient A has the higher viral load. The difference is $\Delta C_t = 30 - 20 = 10$ cycles. Since each cycle doubles the DNA:
\[
\text{factor} = 2^{\Delta C_t} = 2^{10} = 1024.
\]
Patient A therefore started with approximately $1024$ times more viral genetic material than patient B. \textbf{Common mistake:} thinking a high $C_t$ means high viral load --- it is the opposite.
\end{exoresolu}

\courssub{The Cameroonian and international regulatory framework}

Powerful technologies require governance. Cameroon is a Party to the \cle{Cartagena Protocol on Biosafety} (an international agreement under the Convention on Biological Diversity) which regulates the transboundary movement of living modified organisms, based on the \emph{precautionary principle} and on Advance Informed Agreement: an exporting country must inform and obtain consent from the importing country.

At national level, \cle{Law No. 2003/006 of 21 April 2003} lays down the safety rules governing modern biotechnology in Cameroon. It established the \cle{National Biosafety Committee}, a multidisciplinary body (scientists, administrators, civil society) that examines applications for contained use, field trials or import of GMOs, conducts risk assessments, and advises the government before any authorisation. Thus, while Cameroon develops biotechnological capacity (e.g.\ molecular diagnosis of malaria and COVID-19, research on improved crop varieties), every step is framed by biosafety rules protecting human health and our exceptional biodiversity.

\begin{aretenir}[Key points of Section 5]
\begin{itemize}
\item Recombinant DNA technology produces human insulin, growth hormone, factor VIII, monoclonal antibodies (hybridomas, CHO cells, plants) and mRNA vaccines (lipid nanoparticles, no genomic integration).
\item GM crops (Bt cotton, herbicide-tolerant maize, Golden Rice) and CRISPR-edited crops (drought-tolerant rice) improve yields and nutrition, but require rigorous risk assessment and management of resistance evolution.
\item Bioremediation and phytoremediation use microbes and plants to clean polluted environments; bioindicators monitor pollution.
\item Advanced biofuels (cellulosic ethanol, algal biodiesel, biohydrogen) avoid competition with food.
\item Diagnostics (qPCR, ELISA, lateral flow tests, NGS) revolutionise disease detection --- vital for malaria and emerging diseases in Cameroon.
\item In Cameroon, biotechnology is regulated by Law No. 2003/006 and the National Biosafety Committee, under the Cartagena Protocol.
\end{itemize}
\end{aretenir}

\begin{exercice}[Practice questions]
\begin{enumerate}
\item Explain why recombinant human insulin is preferred to insulin extracted from animal pancreases (three reasons).
\item A farmer plants Bt cotton on his entire farm for ten consecutive years. Predict, with justification, the likely evolution of the bollworm population, and propose a strategy to prevent it.
\item Distinguish biostimulation from bioaugmentation in the treatment of an oil-polluted mangrove.
\item In a qPCR test, sample X has $C_t = 18$ and sample Y has $C_t = 24$. Calculate the ratio of their initial template quantities.
\end{enumerate}
\end{exercice}

\begin{corrige}[Exercise 5]
\textbf{1.} Recombinant insulin is (i) identical to human insulin, so no allergic reactions; (ii) available in unlimited quantities from fermenters, independent of animal supply; (iii) cheaper and free of animal pathogens.

\textbf{2.} Continuous exposure to the Bt toxin selects rare resistant caterpillars; they survive and reproduce, so the frequency of resistance alleles rises until the population becomes largely resistant (directional selection). Prevention: maintain refuge plots of non-Bt cotton nearby, where susceptible insects survive and mate with resistant ones, diluting resistance alleles.

\textbf{3.} Biostimulation = adding nutrients (N, P) to stimulate indigenous pollutant-degrading microbes already present. Bioaugmentation = introducing specialised degrading micro-organisms from outside into the polluted site.

\textbf{4.} $\Delta C_t = 24 - 18 = 6$; ratio $= 2^6 = 64$. Sample X contained about 64 times more initial template than sample Y.
\end{corrige}

\newpage

\coursec{Integration activity}

\coursec{Integration Activity: Cassava Valorisation by Recombinant Yeast}

\courssub{Aims of the activity}

This integration activity invites you to act as a junior biotechnologist in a Cameroonian agro-food research team. By the end of the activity, you should be able to:

\begin{itemize}
\item \cle{Analyse} the scientific and technical limitations of a traditional biotechnological process (artisanal cassava fermentation) using your knowledge of microbial metabolism;
\item \cle{Design} a recombinant DNA construct (promoter -- gene of interest -- terminator -- selectable marker) suitable for expressing a thermostable $\alpha$-amylase in the yeast \emph{Saccharomyces cerevisiae};
\item \cle{Plan} a screening and assay strategy (starch-iodine test, reducing-sugar assay) to select and characterise transformed clones;
\item \cle{Evaluate} a pilot-scale bioprocess quantitatively (yield, productivity) and critically (biosafety, cost, social acceptability);
\item \cle{Communicate} scientific results through a poster and an oral defence, as expected at the GCE Advanced Level.
\end{itemize}

The activity is carried out in \textbf{groups of four} and mobilises the core lessons of this chapter: traditional biotechnology, recombinant DNA technology, gene alteration, screening methods, and the applications and limits of biotechnology.

\courssub{Situation}

\textbf{Cassava in Cameroon: an opportunity and a bottleneck.}\par
Cassava (\emph{Manihot esculenta}) is one of the most important food crops in Cameroon. In the Littoral, Centre and South regions, tubers are processed by artisanal fermentation into products such as \emph{bâton de manioc} (bobolo/miondo), gari and \emph{kumkum}. Fermentation is carried out by naturally occurring micro-organisms (mainly lactic acid bacteria and yeasts) present on the tubers and in the environment. This is a typical example of \cle{traditional biotechnology}: it softens the tissue, develops flavour and, crucially, helps to degrade the \cle{cyanogenic glucosides} (linamarin and lotaustralin) which release toxic hydrogen cyanide (HCN) when the tissue is damaged.

However, a field survey conducted by a (fictional) cooperative in Édéa, ``Coopérative Manioc Plus'', revealed the following problems:

\begin{center}
\begin{tabularx}{\linewidth}{|l|X|}\hline
\rowcolor{popblueL} \textbf{Problem observed} & \textbf{Data from the survey} \\ \hline
Residual cyanide & Some batches of bobolo still contain HCN levels considered unsafe when fermentation is too short (less than 3 days). \\ \hline
Variable quality & Fermentation time varies from 3 to 7 days depending on ambient temperature and the chance micro-organisms present; taste and texture are inconsistent. \\ \hline
Low starch valorisation & The cooperative wants to produce \textbf{glucose syrup} from cassava starch, but the starch must first be hydrolysed. Industrial hydrolysis uses purified $\alpha$-amylase bought abroad, which is expensive. \\ \hline
Yield & Only about $60\ \mathrm{kg}$ of marketable product is obtained per $100\ \mathrm{kg}$ of fresh tubers. \\ \hline
\end{tabularx}
\end{center}

\textbf{The project.}\par
The cooperative contacts your team, based at a university biotechnology laboratory in Douala. Your mission: design a \textbf{mini-project} to improve cassava valorisation using modern biotechnology. The central idea is to construct a strain of the baker's yeast \emph{Saccharomyces cerevisiae} that secretes a \cle{thermostable $\alpha$-amylase} from the bacterium \emph{Bacillus licheniformis}. Such a yeast could, in a single step, liquefy cassava starch at high temperature ($90$--$95\ ^{\circ}\mathrm{C}$ for the enzyme step) and then ferment the released sugars, reducing the need to import purified enzyme.

\textbf{Data provided to your team.}

\begin{itemize}
\item The \emph{B. licheniformis} $\alpha$-amylase gene (\emph{amyL}) is about $1.5\ \mathrm{kb}$ long; its sequence is available in public databases.
\item You have a \textbf{shuttle plasmid} (replicating in both \emph{E. coli} and yeast) carrying: an origin of replication for \emph{E. coli}, a yeast $2\mu$ origin, a strong constitutive yeast promoter (P$_{\text{PGK}}$), a multiple cloning site (MCS), a yeast terminator (T$_{\text{CYC1}}$), an ampicillin-resistance gene (\emph{bla}, for selection in \emph{E. coli}) and the \emph{URA3} gene (selection marker in yeast: allows growth on medium lacking uracil).
\item Restriction sites available in the MCS: \emph{Eco}RI, \emph{Bam}HI, \emph{Hind}III. The \emph{amyL} gene can be amplified by PCR with \emph{Bam}HI and \emph{Hind}III sites added at its two ends.
\item Starch-iodine test: iodine gives a blue-black colour with starch; a clear halo around a colony on starch agar indicates starch hydrolysis.
\item Reducing sugars can be assayed colorimetrically (DNS method): the absorbance at $540\ \mathrm{nm}$ is proportional to the glucose concentration.
\item Pilot trial target: convert $10\ \mathrm{kg}$ of cassava starch per batch; native (non-transformed) yeast hydrolyses less than $5\ \%$ of the starch, while a good engineered strain should hydrolyse at least $80\ \%$.
\end{itemize}

\courssub{Tasks}

Answer the following questions progressively. Questions 1--3 concern traditional biotechnology; questions 4--7 concern recombinant DNA technology; questions 8--10 concern evaluation and critical discussion.

\begin{enumerate}
\item Explain why artisanal cassava fermentation is described as \emph{biotechnology}, and identify the type of metabolism (respiration or fermentation) mainly used by the micro-organisms involved.
\item Using the survey data, state \textbf{three} limitations of the artisanal process and, for each one, explain its biological or technological cause.
\item Explain the biochemical role of fermentation in the detoxification of cassava (cyanogenic glucosides $\rightarrow$ HCN). Why is an uncontrolled, spontaneous fermentation risky from a public-health point of view?
\item The team chooses \emph{S. cerevisiae} as the host rather than \emph{E. coli}. Give \textbf{two} reasons that justify this choice for a food application.
\item Draw a labelled diagram of the final expression cassette as it will appear in the shuttle plasmid, showing: promoter, \emph{amyL} gene, terminator, yeast selectable marker, and the two restriction sites used for cloning. Indicate the orientation of each element.
\item Describe, in numbered steps, how you would obtain the recombinant plasmid, from the \emph{B. licheniformis} chromosome to a transformed yeast colony. Name the key enzymes used at each step.
\item Explain the role of the \emph{URA3} marker during the selection of transformed yeasts. What culture medium would you use, and why can untransformed yeasts not grow on it?
\item Describe the \textbf{screening test} on starch agar that distinguishes yeast clones secreting active $\alpha$-amylase from negative clones. What result do you expect for a positive clone, and why?
\item In the pilot trial, a transformed strain hydrolyses $82\ \%$ of the $10\ \mathrm{kg}$ of starch per batch, and $1\ \mathrm{kg}$ of starch yields theoretically $1.1\ \mathrm{kg}$ of glucose. Calculate the mass of glucose expected per batch, then the mass actually recovered if the measured recovery is $90\ \%$ of the theoretical value. Compare with the native strain ($5\ \%$ hydrolysis) and conclude.
\item Discuss \textbf{three} points that must be examined before the engineered yeast is used at village scale: one biosafety point, one economic point, and one social/ethical point. Suggest a precaution for each.
\end{enumerate}

\courssub{Model answer}

\textbf{Question 1.} Artisanal cassava fermentation is biotechnology because it uses \cle{living micro-organisms} (lactic acid bacteria and yeasts naturally present on the tubers) to transform a raw material into a useful product with improved properties (texture, flavour, safety, preservation). The micro-organisms mainly use \cle{anaerobic metabolism}: yeasts perform alcoholic fermentation (glucose $\rightarrow$ ethanol $+$ \ce{CO2}) and lactic acid bacteria perform lactic fermentation (glucose $\rightarrow$ lactic acid), which acidifies the product and inhibits spoilage organisms.

\textbf{Question 2.} Three limitations:
\begin{itemize}
\item \emph{Residual cyanide:} detoxification depends on the enzyme linamarase, produced by the plant tissue and by some fermenting microbes; if fermentation is too short, linamarin is not fully degraded and HCN remains in the food.
\item \emph{Variable quality:} the process is \emph{spontaneous} --- the microbial community is not controlled, so its composition, and therefore the duration (3--7 days) and the sensory result, vary with temperature and chance contamination.
\item \emph{Low valorisation of starch:} the artisanal microbes do not produce enough amylase to hydrolyse starch into fermentable or marketable sugars, so glucose-syrup production is impossible without imported enzymes, which raises costs.
\end{itemize}

\textbf{Question 3.} When cassava cells are crushed, the glucoside linamarin comes into contact with the enzyme linamarase ($\beta$-glucosidase), which hydrolyses it into glucose and acetone cyanohydrin; the cyanohydrin then decomposes spontaneously, releasing \ce{HCN}. Prolonged fermentation, combined with pressing and drying/cooking, allows the volatile HCN to escape. A spontaneous, uncontrolled fermentation is risky because neither the duration nor the microbial enzymatic activity is guaranteed: a short or failed fermentation leaves dangerous residual cyanide, which can cause acute poisoning or, chronically, neurological disorders (e.g. konzo).

\textbf{Question 4.} Two reasons for choosing \emph{S. cerevisiae}:
\begin{enumerate}
\item It has a long history of safe use in food (bread, beer, palm wine) and is \cle{Generally Recognised As Safe (GRAS)}, unlike \emph{E. coli}, which is associated with disease and may produce endotoxins (lipopolysaccharides) undesirable in food.
\item As a eukaryote, it performs proper protein folding, post-translational modifications and \cle{secretion} via the ER-Golgi pathway; moreover, it naturally ferments sugars to ethanol, so the same organism can both help hydrolyse starch and ferment the released sugars in a consolidated bioprocess.
\end{enumerate}

\textbf{Question 5.} Expression cassette in the shuttle plasmid (orientation $5' \rightarrow 3'$):

\begin{popfigure}
\begin{tikzpicture}[scale=1.1, every node/.style={font=\small}]
% plasmid backbone circle
\draw[thick, popmuted] (0,0) circle (2.8);
\node[popmuted] at (0,-2.3) {shuttle plasmid};
\node[popmuted] at (-2.0,1.8) {\emph{bla} (Amp$^R$)};
\node[popmuted] at (2.0,1.8) {ori \emph{E. coli} / $2\mu$};
% cassette blocks (linear representation on the circle)
\draw[thick, popblue, fill=popblueL] (-2.8,-0.7) rectangle (-1.0,0.7);
\node at (-1.9,0) {P$_{\text{PGK}}$};
\draw[thick, popgreen, fill=popgreenL] (-1.0,-0.7) rectangle (1.4,0.7);
\node at (0.2,0) {\emph{amyL} ($\alpha$-amylase)};
\draw[thick, poporange, fill=white] (1.4,-0.7) rectangle (2.8,0.7);
\node at (2.1,0) {T$_{\text{CYC1}}$};
% restriction sites labels below cassette
\draw[thick, popdark] (-1.0,-1.0) -- (-1.0,-1.3);
\node[below, popdark] at (-1.0,-1.3) {\emph{Bam}HI};
\draw[thick, popdark] (1.4,-1.0) -- (1.4,-1.3);
\node[below, popdark] at (1.4,-1.3) {\emph{Hind}III};
% URA3 marker elsewhere on plasmid
\draw[thick, poppurple, fill=white] (-0.7,2.3) rectangle (0.7,3.1);
\node at (0,2.7) {\emph{URA3}};
% transcription arrow
\draw[-latex, thick, popblue] (-1.9,0.95) -- (-0.5,0.95);
\node[above, popblue] at (-1.2,0.95) {transcription};
\end{tikzpicture}
\legende{Map of the recombinant shuttle plasmid: the \emph{amyL} gene is inserted between the \emph{Bam}HI and \emph{Hind}III sites of the MCS, under the control of the yeast promoter P$_{\text{PGK}}$ and terminator T$_{\text{CYC1}}$; \emph{URA3} selects transformed yeasts.}
\end{popfigure}

\textbf{Question 6.} Steps to obtain the recombinant plasmid:
\begin{enumerate}
\item Extract chromosomal DNA from \emph{B. licheniformis}.
\item Amplify the \emph{amyL} gene by \cle{PCR} using \emph{Taq} DNA polymerase and primers carrying a \emph{Bam}HI site (forward) and a \emph{Hind}III site (reverse).
\item Digest the purified PCR product and the shuttle plasmid separately with the \cle{restriction enzymes} \emph{Bam}HI and \emph{Hind}III; the two different sticky ends ensure directional cloning (correct orientation).
\item Join the insert and the vector with \cle{T4 DNA ligase}.
\item Transform competent \emph{E. coli} (e.g. by heat shock or electroporation); select on LB-ampicillin medium (only bacteria carrying the plasmid grow).
\item Screen colonies (colony PCR or restriction analysis of extracted plasmids) to confirm the $1.5\ \mathrm{kb}$ insert; verify by DNA sequencing.
\item Amplify the verified plasmid in \emph{E. coli}, purify it (miniprep), and transform \emph{S. cerevisiae} (e.g. by lithium acetate/PEG method).
\item Plate the yeasts on synthetic defined medium \textbf{lacking uracil} (SD-Ura): only transformed clones grow.
\end{enumerate}

\textbf{Question 7.} \emph{URA3} encodes orotidine-5'-phosphate decarboxylase, an enzyme required for uracil biosynthesis. The host yeast strain is \emph{ura3}$^{-}$ (auxotrophic for uracil): it cannot synthesise its own uracil. On a medium lacking uracil, only yeasts that have taken up the plasmid (and therefore a functional \emph{URA3}) can synthesise uracil and grow; untransformed yeasts die. This is \cle{selection} (survival based on a vital function), distinct from \cle{screening} (detection of a specific phenotype like enzyme activity).

\textbf{Question 8.} Screening on starch agar: replica-plate or streak the selected yeast colonies on agar containing soluble starch, incubate at $30\ ^{\circ}\mathrm{C}$ for 24--48 h, then flood the plate with iodine solution (iodine-potassium iodide). Starch forms a blue-black complex with iodine everywhere \emph{except} around colonies that secrete active $\alpha$-amylase: there, starch has been hydrolysed to maltose/glucose, so a \cle{clear halo} appears. Positive clone $=$ colony surrounded by a clear halo; the larger the halo, the higher the amylase activity. The best clones are then confirmed by the quantitative DNS assay of reducing sugars (absorbance at $540\ \mathrm{nm}$).

\textbf{Question 9.} Calculations for one $10\ \mathrm{kg}$ batch of starch:
\begin{align*}
m_{\text{starch hydrolysed}} &= 10 \times 0.82 = 8.2\ \mathrm{kg}\\
m_{\text{glucose, theoretical}} &= 8.2 \times 1.1 = 9.02\ \mathrm{kg}\\
m_{\text{glucose, recovered}} &= 9.02 \times 0.90 = 8.118\ \mathrm{kg} \approx 8.1\ \mathrm{kg}
\end{align*}
For the native strain ($5\ \%$ hydrolysis, same recovery):
\[
m_{\text{glucose, native}} = 10 \times 0.05 \times 1.1 \times 0.90 = 0.495\ \mathrm{kg} \approx 0.50\ \mathrm{kg}
\]
The engineered strain therefore recovers about \textbf{16 times more glucose} per batch ($8.1\ \mathrm{kg}$ versus $0.50\ \mathrm{kg}$) and exceeds the $80\ \%$ hydrolysis target: the pilot trial is successful.

\textbf{Question 10.} Discussion before village-scale use:
\begin{itemize}
\item \emph{Biosafety:} the engineered yeast must be used in closed vessels to avoid uncontrolled release into the environment; the final food strain carries no antibiotic-resistance gene (\emph{bla} is only for the \emph{E. coli} cloning step) and \emph{URA3} is a harmless metabolic gene. \textbf{Precaution:} contain the process physically and inactivate the yeast by heat treatment ($>70\ ^{\circ}\mathrm{C}$) before the product is sold.
\item \emph{Economic:} producing glucose syrup locally with the engineered yeast removes the recurring cost of imported purified $\alpha$-amylase, but the upstream laboratory steps (PCR, sequencing, strain verification) are capital-intensive. \textbf{Precaution:} once the strain is built and validated, maintain it as a stable starter culture that the cooperative can multiply cheaply on-site (like a bread-yeast or palm-wine starter), avoiding repeated lab interventions.
\item \emph{Social/ethical:} consumers may fear ``genetically modified'' food and reject the product. \textbf{Precaution:} engage the community early through the cooperative, explain that the yeast is a \cle{processing aid} removed/inactivated before consumption (similar in spirit to traditional fermentation starters), and seek regulatory approval from the relevant Cameroonian authorities (e.g. MINSANTE, MINADER) for transparent labelling.
\end{itemize}

\begin{attention}[Common mistakes to avoid]
Do not confuse \textbf{selection} (a marker that keeps only transformed cells alive, e.g.\ \emph{URA3} on uracil-free medium) with \textbf{screening} (a test that reveals which survivors actually express the desired activity, e.g.\ the starch-iodine halo). In the cloning scheme, the insert and the vector must be cut with the \emph{same} restriction enzymes, and the promoter must be placed \emph{upstream} ($5'$) of the gene, in the correct orientation. In yield calculations, always apply the percentages sequentially (hydrolysis $\times$ theoretical yield $\times$ recovery) and state your units at every step.
\end{attention}

\begin{aretenir}[What this activity consolidates]
Traditional fermentation is real biotechnology but suffers from variability and safety limits; recombinant DNA technology (PCR, restriction enzymes, ligase, shuttle vector, selectable marker) allows a precise, heritable improvement of a food micro-organism; screening assays verify enzyme activity; and a responsible biotechnologist always evaluates yield \emph{and} biosafety, cost and social acceptability before scaling up.
\end{aretenir}

\begin{savaistu}[Link to medical biotechnology (Homeostasis)]
The same recombinant DNA principles used here to produce a thermostable enzyme are applied to produce human proteins for therapy: \cle{human insulin} (diabetes homeostasis), \cle{growth hormone}, \cle{clotting factors} and \cle{vaccines} (e.g. hepatitis B surface antigen in yeast). These are major applications of biotechnology in the service of human homeostasis, directly addressing Lesson 84 of the syllabus.
\end{savaistu}

\newpage

\coursec{Methods and worked examples}

\coursec{Biotechnology}

\courssub{Skill 1: Analysing a traditional fermentation process}

\begin{methode}[How to describe and explain a traditional fermentation]
When a GCE question asks you to describe a traditional biotechnological process (palm wine, \emph{fufu} water, \emph{burukutu}, yoghurt, bread), proceed in numbered steps:
\begin{enumerate}
\item \textbf{Identify the substrate} (the raw material: palm sap, milk, cassava, dough) and the \textbf{micro-organism} responsible (yeast \emph{Saccharomyces cerevisiae}, lactic acid bacteria \emph{Lactobacillus} spp., \emph{Acetobacter} spp.).
\item \textbf{State the type of respiration}: fermentation is \cle{anaerobic} (occurs without oxygen).
\item \textbf{Write the balanced summary equation}. For alcoholic fermentation: $\ce{C6H12O6 -> 2C2H5OH + 2CO2}$ + small amount of energy (2 ATP). For lactic fermentation: $\ce{C6H12O6 -> 2 CH3CHOHCOOH}$ + 2 ATP.
\item \textbf{Explain the conditions}: warmth (optimum about $25$--$35\ ^\circ\mathrm{C}$), absence of air, presence of fermentable sugar, time, pH near neutral initially.
\item \textbf{Link the products to the observed changes}: bubbles = $\ce{CO2}$; sour taste = lactic acid; alcohol smell = ethanol; preservation = acid/alcohol inhibits spoilage microbes.
\item \textbf{Conclude with the definition}: traditional biotechnology is the use of living micro-organisms or their enzymes to transform raw materials into useful products, \textbf{without altering their genetic material}.
\end{enumerate}
\textbf{Reflex:} always mention that fermentation yields only \textbf{2 ATP per glucose} (versus about 36--38 in aerobic respiration), because the electron transport chain is not used and glucose is only partially oxidised.
\end{methode}

\begin{popfigure}
\begin{tikzpicture}[scale=0.9, every node/.style={font=\small}]
% Palm wine fermentation setup
\draw[popdark, thick] (0,0) -- (6,0) -- (6,5) -- (0,5) -- cycle;
\draw[popblue, thick] (0.5,0.5) -- (5.5,0.5) -- (5.5,3.5) -- (0.5,3.5) -- cycle;
\draw[popblue!30, fill=popblue!10] (0.5,0.5) rectangle (5.5,3.5);
\node[popdark] at (3,4.2) {Calabash with fresh palm sap};
\node[popdark] at (3,2.5) {Sucrose, glucose, fructose};
\node[popdark] at (3,1.8) {Natural yeasts (\emph{Saccharomyces})};
\node[popdark] at (3,1.1) {Lactic acid bacteria (\emph{Lactobacillus})};
% Bubbles
\foreach \x/\y in {1.2,1.0/1.5,2.5/1.2,3.8/1.6,4.5/1.0} {
    \draw[popblue, thick] (\x,\y) ellipse (0.15 and 0.1);
    \draw[popblue, thick] (\x+0.05,\y+0.2) ellipse (0.1 and 0.07);
}
\draw[->, poporange, thick] (6.5,2.5) -- (8.5,2.5);
\node[poporange, right, text width=3cm, align=center] at (9.5,2.5) {Products: Ethanol + CO$_2$ + Lactic acid + Acetic acid (later)};
% Equations
\node[popdark, left, text width=5cm, align=left] at (-0.5, -0.5) {
\textbf{Alcoholic fermentation:}\\
$\ce{C6H12O6 -> 2 C2H5OH + 2 CO2}$ (2 ATP)\\
\textbf{Lactic fermentation:}\\
$\ce{C6H12O6 -> 2 CH3CHOHCOOH}$ (2 ATP)\\
\textbf{Acetic oxidation (aerobic):}\\
$\ce{C2H5OH + O2 -> CH3COOH + H2O}$
};
\end{tikzpicture}
\legende{Traditional palm wine fermentation in a calabash: spontaneous mixed-culture fermentation.}
\end{popfigure}

\begin{exoresolu}[Fermentation of palm wine in the South-West Region]
A student in Buea collects fresh palm sap in a clean calabash. After 6 hours the sap is sweet and slightly fizzy; after 24 hours it tastes strongly alcoholic and sour. (a) Name the micro-organism mainly responsible for the first stage. (b) Write the equation of the reaction occurring. (c) Explain why the sap becomes fizzy. (d) Explain why the drink becomes sour after 24 hours. (e) State two reasons why palm wine spoils faster than bottled beer.
\tcblower
\textbf{(a)} The first stage is carried out by \textbf{yeasts} (mainly \emph{Saccharomyces} species) naturally present on the palm tree bark, in the air and in the calabash. No starter culture is added: this is a \textbf{spontaneous fermentation}, characteristic of traditional biotechnology.

\textbf{(b)} The yeasts perform \textbf{alcoholic fermentation} of the sugars (mainly sucrose, glucose and fructose) in the sap:
\[ \ce{C6H12O6 -> 2 C2H5OH + 2 CO2} + \text{energy (2 ATP)} \]
This is anaerobic respiration: glucose is only partially oxidised, so most of the chemical energy remains locked in the ethanol.

\textbf{(c)} The fizziness is caused by the \textbf{carbon dioxide} released by the fermentation. The gas dissolves in the sap and escapes as bubbles, exactly as in bread dough rising or in the brewing of \emph{burukutu} (millet beer) in the North.

\textbf{(d)} After about 24 hours, \textbf{lactic acid bacteria} (\emph{Lactobacillus}) and \textbf{acetic acid bacteria} (\emph{Acetobacter}) multiply in the wine. Lactic acid bacteria convert remaining sugars to lactic acid (lactic fermentation), and acetic acid bacteria oxidise ethanol to ethanoic (acetic) acid in the presence of air:
\[ \ce{C2H5OH + O2 -> CH3COOH + H2O} \]
The accumulation of these acids lowers the pH and gives the sour taste. This is why palm wine left overnight is used like vinegar in some households.

\textbf{(e)} Palm wine spoils faster than bottled beer because:
\begin{itemize}
\item it is \textbf{not pasteurised or filtered}, so the mixed microbial population (yeasts, lactic and acetic bacteria) remains alive and active;
\item it is kept in \textbf{open or non-sterile containers} at ambient tropical temperature ($25$--$30\ ^\circ\mathrm{C}$), which favours rapid microbial growth and exposure to oxygen (needed by acetic acid bacteria).
\end{itemize}
\textbf{Examiner's tip:} in part (b) do not forget to state that the process is anaerobic and releases only 2 ATP; this is the marking point most often missed.
\end{exoresolu}

\courssub{Skill 2: Interpreting a fermentation experiment (data analysis)}

\begin{methode}[How to analyse gas-production or growth data in biotechnology]
\begin{enumerate}
\item \textbf{Read the axes and units first}; identify the independent variable (time, temperature, sugar concentration, pH) and the dependent variable (volume of $\ce{CO2}$, optical density, alcohol content, pH change).
\item \textbf{Describe the curve in phases}: lag phase (slow start, cells adapting), exponential/active phase (rapid gas production or growth), stationary phase (substrate exhausted or product inhibiting), death/decline phase.
\item \textbf{Explain each phase biologically}, not just descriptively: e.g.\ stationary phase because sugar is finished or ethanol has reached a toxic concentration for the yeast.
\item \textbf{Compare treatments} by quoting figures (rates = gradient of the curve in the exponential phase).
\item \textbf{Draw a conclusion} that answers the aim, and state one \textbf{control variable} that had to be kept constant (temperature, volume, yeast mass, sugar concentration, pH, anaerobiosis).
\end{enumerate}
\textbf{Reflex:} rate $= \dfrac{\text{change in dependent variable}}{\text{time}}$; always give units (e.g.\ $\mathrm{cm^3\,min^{-1}}$).
\end{methode}

\begin{popfigure}
\begin{tikzpicture}[scale=0.8]
\begin{axis}[
    width=12cm, height=7cm,
    xlabel={Time (hours)}, ylabel={Volume of CO$_2$ (cm$^3$)},
    xmin=0, xmax=24, ymin=0, ymax=100,
    xtick={0,4,8,12,16,20,24}, ytick={0,20,40,60,80,100},
    grid=major, grid style={popline},
    legend pos=north west,
    legend style={font=\small, fill=white, draw=popdark},
]
\addplot[popblue, thick, mark=*] coordinates {(0,0) (2,5) (4,20) (6,45) (8,70) (10,85) (12,92) (14,95) (16,96) (18,96) (20,95) (22,93) (24,90)};
\addlegendentry{37°C (optimum)}
\addplot[poporange, thick, mark=square*] coordinates {(0,0) (2,2) (4,8) (6,15) (8,22) (10,28) (12,32) (14,34) (16,35) (18,35) (20,34) (22,33) (24,32)};
\addlegendentry{25°C}
\addplot[popgreen, thick, mark=triangle*] coordinates {(0,0) (2,0.5) (4,1) (6,2) (8,3) (10,4) (12,4.5) (14,4.8) (16,5) (18,5) (20,4.8) (22,4.5) (24,4)};
\addlegendentry{10°C}
\addplot[poppink, thick, mark=diamond*] coordinates {(0,0) (2,0) (4,0) (6,0) (8,0) (10,0) (12,0) (14,0) (16,0) (18,0) (20,0) (22,0) (24,0)};
\addlegendentry{70°C (denatured)}
% Phase labels
\node[popdark, font=\small, rotate=90] at (2,50) {Lag};
\node[popdark, font=\small, rotate=90] at (6,50) {Exponential};
\node[popdark, font=\small, rotate=90] at (14,50) {Stationary};
\node[popdark, font=\small, rotate=90] at (20,50) {Decline};
\end{axis}
\end{tikzpicture}
\legende{Typical fermentation curves at different temperatures. At 37°C: rapid exponential phase, high yield. At 10°C: slow throughout. At 70°C: no activity (enzymes denatured).}
\end{popfigure}

\begin{exoresolu}[Effect of temperature on yeast fermentation]
Equal masses of yeast and glucose solution were placed in four water baths and the volume of $\ce{CO2}$ collected after 20 minutes was measured:
\begin{center}
\begin{tabular}{|l|c|c|c|c|}
\hline
\rowcolor{popblueL} Temperature ($^\circ$C) & 10 & 25 & 37 & 70 \\
\hline Volume of $\ce{CO2}$ (cm$^3$) & 4 & 22 & 35 & 0 \\
\hline
\end{tabular}
\end{center}
(a) Calculate the rate of $\ce{CO2}$ production at $37\ ^\circ$C. (b) Explain the results at $10\ ^\circ$C and at $70\ ^\circ$C. (c) State two variables that must be controlled. (d) Suggest the optimum temperature and justify.
\tcblower
\textbf{(a)} Rate $= \dfrac{\text{volume of } \ce{CO2}}{\text{time}} = \dfrac{35\ \mathrm{cm^3}}{20\ \mathrm{min}} = 1.75\ \mathrm{cm^3\,min^{-1}}$.

\textbf{(b)} At $10\ ^\circ$C the rate is low ($0.2\ \mathrm{cm^3\,min^{-1}}$) because the \textbf{enzymes} of the yeast (zymase complex) have low kinetic energy: few enzyme--substrate collisions occur per unit time, so fermentation is slow but the yeast is \emph{not dead}. At $70\ ^\circ$C no gas is produced because the high temperature has \textbf{denatured the enzymes}: the tertiary structure and the active site are destroyed irreversibly, and the yeast cells are killed. No fermentation can occur.

\textbf{(c)} Variables to control: the \textbf{mass/concentration of yeast}, the \textbf{concentration and volume of glucose solution}, the \textbf{duration} of the experiment, \textbf{pH}, and \textbf{absence of oxygen} (anaerobic conditions).

\textbf{(d)} The optimum is around $37\ ^\circ$C (possibly slightly below or above), because this temperature gave the greatest volume of gas in the same time, i.e.\ the highest rate. To confirm, the experiment should be repeated at $30$, $35$, $40$ and $45\ ^\circ$C to narrow down the peak.

\textbf{Examiner's tip:} ``the yeast is killed by cold'' is a classic error: cold only \emph{slows} enzyme activity; only heat \emph{denatures} enzymes irreversibly.
\end{exoresolu}

\courssub{Skill 3: Biotechnology at the service of homeostasis, support and locomotion}

\begin{definition}[Biotechnology for homeostasis, support and locomotion]
This refers to the application of biotechnological products and techniques to maintain \textbf{internal stability (homeostasis)}, provide \textbf{structural support}, and enable \textbf{movement (locomotion)} in humans and animals. Key examples include recombinant hormones, biomaterials, tissue engineering, and assistive devices.
\end{definition}

\begin{propriete}[Major applications]
\begin{itemize}
\item \textbf{Glucose homeostasis:} Recombinant human \textbf{insulin} (produced in \emph{E.~coli} or yeast) and \textbf{glucagon} for diabetes mellitus treatment. Continuous glucose monitors (CGMs) and insulin pumps (biotechnology + engineering).
\item \textbf{Calcium/phosphate homeostasis:} Recombinant \textbf{calcitonin} and \textbf{parathyroid hormone} analogues for osteoporosis and hypercalcaemia.
\item \textbf{Growth and metabolism:} Recombinant \textbf{human growth hormone (somatropin)} for growth hormone deficiency, Turner syndrome, chronic renal failure.
\item \textbf{Red blood cell homeostasis:} Recombinant \textbf{erythropoietin (EPO)} for anaemia in chronic kidney disease and chemotherapy.
\item \textbf{Support (skeletal repair):} \textbf{Biomaterials} -- hydroxyapatite coatings, collagen scaffolds, biodegradable polymers (PLA, PGA) for bone grafts and fracture fixation. \textbf{Tissue engineering} of cartilage and bone using mesenchymal stem cells on 3D scaffolds.
\item \textbf{Locomotion (joint/muscle repair):} \textbf{Viscosupplementation} with hyaluronic acid (produced by bacterial fermentation) for osteoarthritis. \textbf{Monoclonal antibodies} (e.g.\ anti-TNF$\alpha$ like adalimumab) for rheumatoid arthritis. \textbf{Myostatin inhibitors} (experimental) for muscular dystrophies.
\item \textbf{Prosthetics and interfaces:} Osseointegrated implants (titanium with bioactive coatings), neural interfaces for bionic limbs.
\end{itemize}
\end{propriete}

\begin{exemplebox}[Insulin and glucose homeostasis]
In type 1 diabetes, pancreatic $\beta$-cells are destroyed $\rightarrow$ no insulin $\rightarrow$ chronic hyperglycaemia. Before 1982, insulin was extracted from pig/cow pancreas (slightly different amino acid sequence $\rightarrow$ allergic reactions, limited supply). Recombinant DNA technology allowed production of \textbf{identical human insulin} in bacteria/yeast: unlimited supply, no immune reaction, consistent quality. This is a direct application of biotechnology to restore \textbf{glucose homeostasis}.
\end{exemplebox}

\begin{experience}[Testing the effect of insulin on blood glucose (simulation)]
\textbf{Aim:} Demonstrate the role of insulin in glucose uptake by cells.
\textbf{Materials:} Yeast suspension (model for cells), glucose solution, recombinant insulin, Benedict's reagent, water bath.
\textbf{Method:} Set up three tubes: (1) yeast + glucose, (2) yeast + glucose + insulin, (3) yeast + water (control). Incubate at $37\ ^\circ\mathrm{C}$ for 30 min. Test supernatant for remaining glucose with Benedict's reagent.
\textbf{Expected:} Tube 2 shows greatest glucose uptake (least reducing sugar) because insulin stimulates glucose transporters. Tube 1 shows moderate uptake. Tube 3 shows no uptake.
\textbf{Conclusion:} Insulin enhances cellular glucose uptake, lowering blood glucose -- a homeostatic mechanism restored by biotechnological insulin.
\end{experience}

\courssub{Skill 4: Correction and gene remediation -- gene therapy and genome editing}

\begin{methode}[How to structure an answer on gene correction (gene therapy / CRISPR)]
\begin{enumerate}
\item \textbf{Define}: gene therapy is the introduction of a \textbf{normal functional allele} into a patient's cells to correct a genetic defect; genome editing (e.g.\ \textbf{CRISPR--Cas9}) \textbf{directly modifies} the faulty DNA sequence.
\item \textbf{Distinguish the two types}: \textbf{somatic} gene therapy (body cells; affects only the patient; not inherited) and \textbf{germ-line} therapy (gametes/early embryo; heritable; ethically controversial and generally prohibited internationally).
\item \textbf{Give the steps for somatic therapy}: isolate the normal gene $\rightarrow$ insert into a \textbf{vector} (often a disabled virus: adenovirus, adeno-associated virus AAV, lentivirus/retrovirus, or non-viral liposomes/nanoparticles) $\rightarrow$ deliver to target cells (ex vivo or in vivo) $\rightarrow$ gene expressed $\rightarrow$ functional protein produced.
\item \textbf{For CRISPR--Cas9}: a \textbf{guide RNA (gRNA)} directs the \textbf{Cas9 nuclease} to the target sequence; Cas9 cuts both strands; the cell repairs the break via NHEJ (error-prone) or HDR (homology-directed repair with a correct template), allowing the faulty base to be corrected.
\item \textbf{Evaluate}: advantages (treats the cause, not symptoms; potential cure) and limitations (delivery efficiency, immune reaction to vector, risk of off-target cuts/insertional mutagenesis, high cost, ethical issues for germ-line).
\end{enumerate}
\textbf{Reflex:} always state whether the correction is heritable or not -- this is the key marking point.
\end{methode}

\begin{popfigure}
\begin{tikzpicture}[scale=0.85, every node/.style={font=\small}]
% Somatic gene therapy flow
\node[draw, rounded corners, fill=popblue!10, thick] (start) at (0,4) {Normal $\beta$-globin gene isolated};
\node[draw, rounded corners, fill=popgreen!10, thick] (vector) at (4,4) {Inserted into viral vector (e.g. lentivirus)};
\node[draw, rounded corners, fill=poporange!10, thick] (delivery) at (8,4) {Ex vivo: patient's haematopoietic stem cells harvested};
\node[draw, rounded corners, fill=poppink!10, thick] (transduce) at (12,4) {Cells transduced with vector};
\node[draw, rounded corners, fill=poppurple!10, thick] (infuse) at (16,4) {Corrected stem cells infused back into patient};
\node[draw, rounded corners, fill=popgold!10, thick] (result) at (20,4) {Normal HbA produced in RBCs};
\draw[->, thick, popdark] (start) -- (vector);
\draw[->, thick, popdark] (vector) -- (delivery);
\draw[->, thick, popdark] (delivery) -- (transduce);
\draw[->, thick, popdark] (transduce) -- (infuse);
\draw[->, thick, popdark] (infuse) -- (result);

% CRISPR
\node[draw, rounded corners, fill=popblue!10, thick] (c1) at (0,0) {Guide RNA (gRNA) designed for HbS mutation};
\node[draw, rounded corners, fill=popgreen!10, thick] (c2) at (4,0) {gRNA + Cas9 protein complex formed};
\node[draw, rounded corners, fill=poporange!10, thick] (c3) at (8,0) {Complex enters nucleus, finds target DNA};
\node[draw, rounded corners, fill=poppink!10, thick] (c4) at (12,0) {Cas9 cuts both strands at mutation site};
\node[draw, rounded corners, fill=poppurple!10, thick] (c5) at (16,0) {HDR repair with correct template $\rightarrow$ HbA restored};
\node[draw, rounded corners, fill=popgold!10, thick] (c6) at (20,0) {Edited cells (somatic or embryo)};
\draw[->, thick, popdark] (c1) -- (c2);
\draw[->, thick, popdark] (c2) -- (c3);
\draw[->, thick, popdark] (c3) -- (c4);
\draw[->, thick, popdark] (c4) -- (c5);
\draw[->, thick, popdark] (c5) -- (c6);

% Labels
\node[popdark, font=\bfseries\small] at (-2,4) {Somatic gene therapy:};
\node[popdark, font=\bfseries\small] at (-2,0) {CRISPR-Cas9 editing:};
\end{tikzpicture}
\legende{Comparison of somatic gene therapy (ex vivo lentiviral) and CRISPR-Cas9 genome editing for sickle cell disease.}
\end{popfigure}

\begin{exoresolu}[Gene therapy for sickle cell disease]
Sickle cell disease is common in Cameroon. Scientists propose two strategies: (A) introducing a normal $\beta$-globin gene into the patient's bone-marrow stem cells using a modified virus; (B) editing the faulty gene directly in embryos using CRISPR--Cas9. (a) Classify each strategy as somatic or germ-line. (b) Explain why bone-marrow stem cells are chosen in strategy A. (c) State the role of the guide RNA in strategy B. (d) Give one ethical objection to strategy B.
\tcblower
\textbf{(a)} Strategy A targets the \textbf{body (somatic) cells} of an already-born patient: it is \textbf{somatic gene therapy}. Its effects are limited to the treated individual and are \textbf{not transmitted} to offspring. Strategy B modifies an \textbf{embryo}, whose cells include the future germ line: it is \textbf{germ-line editing}, and the change would be \textbf{heritable}.

\textbf{(b)} Bone-marrow \textbf{haematopoietic stem cells} are chosen because they are the cells that \textbf{continuously divide and differentiate} to produce all red blood cells throughout life. Correcting the gene in these stem cells means that \emph{all} red blood cells produced afterwards will contain normal haemoglobin, giving a long-lasting (potentially permanent) treatment rather than a temporary one.

\textbf{(c)} The \textbf{guide RNA (gRNA)} has a sequence \textbf{complementary to the target DNA} (the mutated region of the $\beta$-globin gene, specifically the A$\rightarrow$T transversion at codon 6). It binds the target by base pairing and thereby \textbf{directs the Cas9 enzyme} to exactly the right position, where Cas9 cuts both DNA strands. The cell's repair machinery then fixes the break using a correct template (provided or endogenous), restoring the normal sequence (GAG $\rightarrow$ Glu instead of GTG $\rightarrow$ Val).

\textbf{(d)} Ethical objections to embryo editing include: the modification is \textbf{heritable} and affects future generations without their consent; \textbf{off-target cuts} could introduce new harmful mutations into the human gene pool; and it opens the way to \textbf{``designer babies''} (selection for non-medical traits), which raises issues of justice, equity and human dignity.

\textbf{Examiner's tip:} note the link with malaria: the $\mathrm{Hb}^S$ allele is maintained in Central Africa because heterozygotes ($\mathrm{Hb}^A\mathrm{Hb}^S$) are more resistant to severe malaria -- an example of \textbf{balanced polymorphism} that examiners love to combine with this topic.
\end{exoresolu}

\courssub{Skill 5: Describing the steps of recombinant DNA technology}

\begin{methode}[How to present the production of a transgenic organism (e.g.\ insulin)]
Learn the sequence as a numbered chain and be able to name the \textbf{enzyme or tool at each step}:
\begin{enumerate}
\item \textbf{Isolate the gene} of interest (e.g.\ human insulin gene) from donor DNA, or obtain its mRNA and make cDNA using \textbf{reverse transcriptase} (avoids introns).
\item \textbf{Cut} the donor DNA (or cDNA) and the \textbf{vector} (usually a bacterial \textbf{plasmid}) with the same \textbf{restriction endonuclease}, producing complementary \textbf{sticky ends} (or blunt ends).
\item \textbf{Join} the gene into the plasmid using \textbf{DNA ligase} $\rightarrow$ \textbf{recombinant DNA} (rDNA).
\item \textbf{Introduce} the recombinant plasmid into a host cell (e.g.\ \emph{E.~coli}) by \textbf{transformation} (heat shock or electroporation).
\item \textbf{Select} the transformed cells using a \textbf{marker gene} (e.g.\ antibiotic resistance, fluorescence) carried by the plasmid.
\item \textbf{Screen/identify} clones containing the correct insert (e.g.\ blue-white screening, PCR, restriction digest).
\item \textbf{Culture} the host cells in a \textbf{fermenter} (large-scale cloning = gene amplification); the cells express the gene (transcription + translation).
\item \textbf{Extract and purify} the product (e.g.\ insulin) for medical use (downstream processing).
\end{enumerate}
\textbf{Reflex:} restriction enzyme = ``molecular scissors''; ligase = ``molecular glue''; plasmid = ``vector/vehicle''; host bacterium = ``factory''; marker gene = ``selection tool''.
\end{methode}

\begin{popfigure}
\begin{tikzpicture}[scale=0.8, every node/.style={font=\small}]
% Recombinant DNA steps
\node[draw, rounded corners, fill=popblue!10, thick] (s1) at (0,7) {1. Isolate human insulin gene (cDNA via reverse transcriptase)};
\node[draw, rounded corners, fill=popgreen!10, thick] (s2) at (0,5.5) {2. Cut gene and plasmid with same restriction enzyme (e.g. EcoRI) $\rightarrow$ sticky ends};
\node[draw, rounded corners, fill=poporange!10, thick] (s3) at (0,4) {3. Mix gene + plasmid + DNA ligase $\rightarrow$ recombinant plasmid};
\node[draw, rounded corners, fill=poppink!10, thick] (s4) at (0,2.5) {4. Transform into E. coli (heat shock / electroporation)};
\node[draw, rounded corners, fill=poppurple!10, thick] (s5) at (0,1) {5. Plate on ampicillin agar $\rightarrow$ only transformed colonies grow};
\node[draw, rounded corners, fill=popgold!10, thick] (s6) at (0,-0.5) {6. Screen colonies (PCR / restriction digest) for insulin insert};
\node[draw, rounded corners, fill=popblue!20, thick] (s7) at (0,-2) {7. Scale up in fermenter: optimal T, pH, O2, nutrients $\rightarrow$ protein expression};
\node[draw, rounded corners, fill=popgreen!20, thick] (s8) at (0,-3.5) {8. Harvest, lyse cells, purify insulin (chromatography)};
\draw[->, thick, popdark] (s1) -- (s2);
\draw[->, thick, popdark] (s2) -- (s3);
\draw[->, thick, popdark] (s3) -- (s4);
\draw[->, thick, popdark] (s4) -- (s5);
\draw[->, thick, popdark] (s5) -- (s6);
\draw[->, thick, popdark] (s6) -- (s7);
\draw[->, thick, popdark] (s7) -- (s8);

% Side annotations
\node[popdark, right, text width=4cm, align=left] at (6,7) {Reverse transcriptase makes intron-free cDNA};
\node[popdark, right, text width=4cm, align=left] at (6,5.5) {EcoRI cuts GAATTC $\rightarrow$ sticky ends AATT};
\node[popdark, right, text width=4cm, align=left] at (6,4) {Ligase forms phosphodiester bonds};
\node[popdark, right, text width=4cm, align=left] at (6,2.5) {Competent cells take up plasmid};
\node[popdark, right, text width=4cm, align=left] at (6,1) {Amp$^R$ = selectable marker};
\node[popdark, right, text width=4cm, align=left] at (6,-0.5) {Confirm insert orientation/size};
\node[popdark, right, text width=4cm, align=left] at (6,-2) {Fed-batch culture; induction (IPTG)};
\node[popdark, right, text width=4cm, align=left] at (6,-3.5) {Downstream processing};
\end{tikzpicture}
\legende{Step-by-step production of recombinant human insulin in E. coli.}
\end{popfigure}

\begin{exoresolu}[Producing human insulin in bacteria]
(a) Explain why the same restriction enzyme must be used to cut both the human DNA and the plasmid. (b) Name the enzyme that seals the insulin gene into the plasmid. (c) The plasmid carries a gene for resistance to the antibiotic ampicillin. Explain how this is used to identify bacteria that took up the plasmid. (d) Give two advantages of bacterial insulin over insulin extracted from pig pancreas.
\tcblower
\textbf{(a)} A restriction endonuclease cuts DNA at a \textbf{specific recognition sequence} (e.g.\ EcoRI: 5'-GAATTC-3') and leaves characteristic single-stranded projections called \textbf{sticky ends}. Using the \emph{same} enzyme on both the human DNA (or cDNA) and the plasmid produces \textbf{complementary sticky ends} whose base sequences can pair by hydrogen bonding (A with T, G with C). This allows the insulin gene to fit precisely into the opened plasmid. If different enzymes were used, the ends would not match and ligation would fail or be inefficient.

\textbf{(b)} \textbf{DNA ligase} seals the gene into the plasmid by forming phosphodiester bonds in the sugar--phosphate backbone, producing stable recombinant DNA.

\textbf{(c)} The bacteria are spread on agar \textbf{containing ampicillin}. Only bacteria that have taken up the plasmid carry the \textbf{ampicillin-resistance marker gene} ($\mathrm{bla}$, encoding $\beta$-lactamase) and can survive and form colonies; bacteria without the plasmid are killed. The surviving colonies are therefore the \textbf{transformed} bacteria, which can then be screened further for the insulin gene itself (e.g.\ by colony PCR or restriction analysis).

\textbf{(d)} Advantages of bacterial (recombinant) human insulin:
\begin{itemize}
\item it is \textbf{identical to human insulin} (same amino acid sequence), so it causes fewer allergic reactions and no immune rejection, unlike pig insulin which differs by one amino acid (Thr$\rightarrow$Ala at B30);
\item it can be produced in \textbf{large, continuous quantities} in fermenters, independent of animal supply, and is free from the risk of transmitting animal pathogens (e.g.\ prions, viruses).
\end{itemize}
\textbf{Examiner's tip:} always name the \emph{tool} (restriction enzyme, ligase, plasmid, marker gene) and its \emph{role}; a step without its enzyme earns only half the marks.
\end{exoresolu}

\courssub{Skill 6: Reading restriction maps and gel electrophoresis}

\begin{methode}[How to interpret a restriction digest and a gel]
\begin{enumerate}
\item \textbf{Restriction digest:} count the number of \textbf{cut sites} for the enzyme; $n$ cuts in a linear DNA molecule give $n+1$ fragments; a circular plasmid gives $n$ fragments (if $n>0$).
\item \textbf{Fragment sizes} add up to the total DNA length; use this to check your answer.
\item \textbf{Gel electrophoresis:} DNA is negatively charged (phosphate groups), so it migrates to the \textbf{positive electrode (anode)}; \textbf{smaller fragments move faster and further} through the agarose gel matrix.
\item \textbf{Match bands to sizes}: the band nearest the bottom (furthest from the wells) is the smallest fragment. Use a \textbf{DNA ladder} (marker) of known sizes to estimate fragment lengths.
\item \textbf{Compare lanes} (e.g.\ suspect vs crime-scene DNA, or normal vs mutant allele) to draw the conclusion.
\end{enumerate}
\textbf{Reflex:} a mutation that \emph{destroys} a restriction site gives \emph{fewer, larger} fragments; a mutation that \emph{creates} a site gives \emph{more, smaller} fragments. This is the basis of RFLP analysis for sickle cell diagnosis.
\end{methode}

\begin{popfigure}
\begin{tikzpicture}[scale=0.9, every node/.style={font=\small}]
% Gel electrophoresis diagram
\draw[popdark, thick] (0,0) -- (8,0) -- (8,6) -- (0,6) -- cycle;
% Wells
\foreach \x in {1,2.5,4,5.5,7} {
    \draw[popdark, fill=popblue!20] (\x-0.3,5.5) rectangle (\x+0.3,5.8);
    \node[popdark, below] at (\x,5.3) {Well};
}
% Ladder lane (lane 1)
\foreach \y/\lbl in {5.0/10kb, 4.5/5kb, 4.0/3kb, 3.5/2kb, 3.0/1.5kb, 2.5/1kb, 2.0/0.5kb, 1.5/0.2kb} {
    \draw[popdark, thick] (0.7,\y) -- (1.3,\y);
    \node[popdark, right] at (1.4,\y) {\lbl};
}
\node[popdark, above] at (1,5.8) {Ladder};
% Lane 2: Father (HbA HbS)
\foreach \y in {4.8, 4.0} {
    \draw[popblue, thick, fill=popblue!30] (2.2,\y-0.05) rectangle (2.8,\y+0.05);
}
\node[popdark, above] at (2.5,5.8) {Father};
\node[popblue, right] at (2.9,4.8) {1.35 kb};
\node[popblue, right] at (2.9,4.0) {1.15 kb};
% Lane 3: Mother (HbA HbS)
\foreach \y in {4.8, 4.0, 1.5} {
    \draw[popgreen, thick, fill=popgreen!30] (3.7,\y-0.05) rectangle (4.3,\y+0.05);
}
\node[popdark, above] at (4,5.8) {Mother};
\node[popgreen, right] at (4.4,4.8) {1.35 kb};
\node[popgreen, right] at (4.4,4.0) {1.15 kb};
\node[popgreen, right] at (4.4,1.5) {0.20 kb};
% Lane 4: Child (HbS HbS)
\draw[poporange, thick, fill=poporange!30] (5.2,4.8-0.05) rectangle (5.8,4.8+0.05);
\node[popdark, above] at (5.5,5.8) {Child};
\node[poporange, right] at (5.9,4.8) {1.35 kb};
% Anode/cathode
\node[popdark] at (4,0.3) {Anode (+) $\leftarrow$ migration direction $\leftarrow$ Cathode (-)};
% Arrow
\draw[->, thick, popdark] (4,5.5) -- (4,1.0);
\end{tikzpicture}
\legende{RFLP analysis of $\beta$-globin gene with MstII. Normal allele (HbA) cut into 1.15 kb + 0.20 kb; sickle allele (HbS) uncut at 1.35 kb. Father and mother are carriers; child has sickle cell anaemia.}
\end{popfigure}

\begin{exoresolu}[Diagnosing sickle cell anaemia by RFLP]
The normal $\beta$-globin allele ($\mathrm{Hb}^A$) contains a restriction site for the enzyme \emph{Mst}II that is destroyed by the sickle cell mutation ($\mathrm{Hb}^S$). Digestion of the normal allele gives two fragments of 1.15 kb and 0.20 kb; the mutant allele gives a single fragment of 1.35 kb. DNA from a father (lane 1), mother (lane 2) and their child (lane 3) is digested and run on a gel. Lane 1 shows bands at 1.35 and 1.15 kb; lane 2 shows bands at 1.35, 1.15 and 0.20 kb; lane 3 shows a single band at 1.35 kb. (a) State the genotype of each person. (b) Explain why the 0.20 kb band may be faint or missing on a real gel. (c) What is the probability that a second child of this couple is also affected?
\tcblower
\textbf{(a)} Interpretation of the bands:
\begin{itemize}
\item \textbf{Lane 1 (father):} bands at 1.35 kb (mutant allele, uncut) and 1.15 kb (cut normal allele; the 0.20 kb fragment runs off or is too faint). He carries both alleles: genotype $\mathrm{Hb}^A\mathrm{Hb}^S$ -- a \textbf{carrier} (heterozygous), phenotypically normal.
\item \textbf{Lane 2 (mother):} bands at 1.35, 1.15 and 0.20 kb: again both alleles present, genotype $\mathrm{Hb}^A\mathrm{Hb}^S$ -- \textbf{carrier}.
\item \textbf{Lane 3 (child):} only the 1.35 kb band: both alleles are mutant, genotype $\mathrm{Hb}^S\mathrm{Hb}^S$ -- the child \textbf{has sickle cell anaemia}.
\end{itemize}
\textbf{(b)} The 0.20 kb fragment is very small: it binds little DNA stain (ethidium bromide or SYBR Safe), so it fluoresces weakly, and it migrates so fast that it may run off the end of the gel. Hence it is often faint or absent, and diagnosis relies on the 1.35/1.15 kb bands.

\textbf{(c)} Cross: $\mathrm{Hb}^A\mathrm{Hb}^S \times \mathrm{Hb}^A\mathrm{Hb}^S$.
\begin{center}
\begin{tabular}{|c|c|c|}
\hline
\rowcolor{popblueL} Gametes & $\mathrm{Hb}^A$ & $\mathrm{Hb}^S$ \\
\hline $\mathrm{Hb}^A$ & $\mathrm{Hb}^A\mathrm{Hb}^A$ & $\mathrm{Hb}^A\mathrm{Hb}^S$ \\
\hline $\mathrm{Hb}^S$ & $\mathrm{Hb}^A\mathrm{Hb}^S$ & $\mathrm{Hb}^S\mathrm{Hb}^S$ \\
\hline
\end{tabular}
\end{center}
P(affected child, $\mathrm{Hb}^S\mathrm{Hb}^S$) $= \dfrac{1}{4} = 0.25$, i.e.\ \textbf{25\%} for each pregnancy, independent of the first child's genotype.

\textbf{Examiner's tip:} this technique (RFLP + electrophoresis) is exactly how \textbf{prenatal diagnosis} of sickle cell disease is performed -- a very relevant application in Cameroon, where carrier frequency is high (up to 25\% in some zones). Do not confuse ``carrier'' (healthy heterozygote) with ``affected''.
\end{exoresolu}

\courssub{Skill 7: Evaluating applications of biotechnology (structured essay)}

\begin{methode}[How to write a balanced ``applications and implications'' essay]
\begin{enumerate}
\item \textbf{Define biotechnology} in one sentence: the use of living organisms, cells or their components (enzymes, DNA) to produce goods and services useful to humans.
\item \textbf{Organise by sector} -- health, agriculture, environment, industry -- with at least one precise example each:
\begin{itemize}
\item \textbf{Health:} recombinant insulin, growth hormone, vaccines (hepatitis B, HPV), monoclonal antibodies (cancer, autoimmune), gene therapy, prenatal diagnosis of sickle cell disease, PCR diagnosis of HIV/malaria/COVID-19.
\item \textbf{Agriculture:} GM crops (Bt cotton resistant to insects, Golden Rice enriched in $\beta$-carotene/provitamin A), disease-resistant plantain/banana varieties (e.g.\ against BXW, Fusarium wilt), biofertilisers (\emph{Rhizobium}, \emph{Azotobacter}), marker-assisted selection (MAS) for cassava mosaic resistance.
\item \textbf{Environment:} \textbf{bioremediation} (bacteria degrading oil spills in the Niger Delta, heavy metal biosorption), sewage treatment (activated sludge, anaerobic digesters for biogas), biofuels (bioethanol from sugarcane/molasses, biodiesel from \emph{Jatropha}), biosensors for pollution monitoring.
\item \textbf{Industry:} enzymes in detergents (proteases, lipases, amylases), brewing and food processing (rennin/chymosin for cheese, pectinases for fruit juice), bioplastics (PHA, PLA from microbial fermentation).
\end{itemize}
\item \textbf{Discuss risks/concerns}: gene flow to wild relatives (e.g.\ Bt gene to wild cotton), antibiotic-resistance markers (horizontal transfer risk), allergenicity of novel proteins, reduction of biodiversity (monocultures), cost and dependence of smallholder farmers on proprietary seeds, ethical issues (GMOs, germ-line editing, patenting life).
\item \textbf{Conclude} with a balanced judgement and the need for \textbf{regulation/biosafety frameworks} (Cartagena Protocol, national biosafety laws).
\end{enumerate}
\textbf{Reflex:} one well-explained example per sector scores more than a long list of names. Use Cameroonian context where natural.
\end{methode}

\begin{exoresolu}[GM crops for Cameroonian agriculture?]
``Genetically modified crops could improve food security in Cameroon, but their introduction must be carefully controlled.'' Discuss this statement, giving at least two potential benefits and two potential risks.
\tcblower
\textbf{Introduction.} Biotechnology offers tools, such as recombinant DNA technology, to introduce useful genes into crops. Whether Cameroon should adopt GM crops requires weighing benefits against risks within a national biosafety framework.

\textbf{Potential benefits.}
\begin{itemize}
\item \textbf{Pest resistance:} Bt crops carry a gene from the bacterium \emph{Bacillus thuringiensis} coding for a Cry protein toxic to certain insect larvae (Lepidoptera, Coleoptera). Bt cotton or Bt maize would reduce losses to stem borers (\emph{Busseola fusca}, \emph{Chilo partellus}) and cut \textbf{pesticide use}, lowering costs for smallholder farmers and reducing chemical pollution of watercourses (e.g.\ Wouri, Sanaga).
\item \textbf{Improved nutrition and stress tolerance:} biofortified crops (e.g.\ vitamin-A-enriched ``Golden Rice'', provitamin-A cassava developed by IITA) can combat \textbf{vitamin A deficiency}, a real public-health problem in the North and Far North; drought-tolerant maize (WEMA project) could secure yields in the Sudano-Sahelian zone where rainfall is unreliable.
\item \textbf{Disease resistance:} engineering plantains and cassava resistant to viral (e.g.\ banana bunchy top virus, cassava mosaic virus) and bacterial diseases (Xanthomonas wilt of banana) would protect staple foods of the forest and humid savanna zones.
\end{itemize}

\textbf{Potential risks.}
\begin{itemize}
\item \textbf{Environmental:} transgenes may spread by \textbf{pollen to wild relatives} (gene flow), possibly creating resistant weeds (e.g.\ wild cotton in the North); Bt toxin may affect \textbf{non-target insects} such as pollinators or beneficial predators, and pests may evolve \textbf{resistance} to the toxin (requiring refuge strategies).
\item \textbf{Health and socio-economic:} possible \textbf{allergenicity} of new proteins must be tested rigorously (Codex Alimentarius guidelines); seed patents could make farmers \textbf{dependent on companies} for seed each season, threatening local seed-saving traditions and agrobiodiversity (landraces of maize, sorghum, millet).
\end{itemize}

\textbf{Conclusion.} GM crops could indeed contribute to food security in Cameroon, but each variety must be assessed case by case under a \textbf{biosafety framework} (confined field trials, environmental risk assessment, food safety assessment, labelling, post-release monitoring), combined with protection of local varieties and support for agroecology. Biotechnology is a powerful tool, but not a substitute for good agronomy, infrastructure, and equitable policies.

\textbf{Examiner's tip:} a ``discuss'' question demands \emph{both} sides plus a personal, justified conclusion; a one-sided answer cannot score full marks. Mention the \textbf{Cartagena Protocol on Biosafety} (Cameroon is a party) and the national biosafety law.
\end{exoresolu}

\courssub{Skill 8: Calculations in biotechnology -- PCR and microbial growth kinetics}

\begin{methode}[How to calculate DNA amplification and microbial population growth]
\begin{enumerate}
\item \textbf{PCR:} each cycle \textbf{doubles} the number of DNA molecules (assuming 100\% efficiency). After $n$ cycles from $N_0$ starting molecules: $N = N_0 \times 2^{n}$.
\item \textbf{Bacterial growth (exponential phase):} if the generation (doubling) time is $g$ minutes, the number of generations in time $t$ is $n = \dfrac{t}{g}$, and the final population is $N = N_0 \times 2^{n}$.
\item \textbf{Solve for the unknown} using powers of 2 or logarithms: $n = \dfrac{\log(N/N_0)}{\log 2}$.
\item \textbf{Check units} (minutes vs hours) and state assumptions (exponential phase, unlimited nutrients, no inhibition, synchronous division).
\end{enumerate}
\textbf{Reflex:} PCR needs a \textbf{thermostable} DNA polymerase (\emph{Taq} polymerase from \emph{Thermus aquaticus}) because each cycle includes a $\sim 94\ ^\circ\mathrm{C}$ denaturation step; primers define the region amplified.
\end{methode}

\begin{popfigure}
\begin{tikzpicture}[scale=0.8]
\begin{axis}[
    width=12cm, height=7cm,
    xlabel={Time (hours)}, ylabel={Log$_{10}$ Population (CFU/mL)},
    xmin=0, xmax=10, ymin=2, ymax=10,
    xtick={0,2,4,6,8,10}, ytick={2,3,4,5,6,7,8,9,10},
    grid=major, grid style={popline},
    legend pos=south east,
    legend style={font=\small, fill=white, draw=popdark},
]
% Lag
\addplot[popblue, thick, mark=none, domain=0:1] {3};
% Exponential
\addplot[popblue, thick, mark=none, domain=1:5] {3 + (x-1)*1.2};
% Stationary
\addplot[popblue, thick, mark=none, domain=5:8] {7.8};
% Death
\addplot[popblue, thick, mark=none, domain=8:10] {7.8 - (x-8)*0.5};
% Phase labels
\node[popdark, font=\small] at (0.5, 3.5) {Lag};
\node[popdark, font=\small] at (3, 5.5) {Exponential (log)};
\node[popdark, font=\small] at (6.5, 8.2) {Stationary};
\node[popdark, font=\small] at (9, 7) {Death};
% Generation time indicator
\draw[<->, poporange, thick] (2, 4.2) -- (3, 5.4);
\node[poporange, right] at (3.2, 4.8) {Generation time $g$};
\end{axis}
\end{tikzpicture}
\legende{Bacterial growth curve in a batch fermenter. Exponential phase: constant generation time $g$. Stationary phase: nutrients exhausted, waste accumulated.}
\end{popfigure}

\begin{exoresolu}[PCR amplification and fermenter culture]
(a) Starting from a single DNA molecule, how many molecules are present after 30 PCR cycles? Give your answer in standard form to 3 significant figures. (b) A recombinant \emph{E.~coli} culture starts with $1.0 \times 10^{3}$ cells and divides every 20 minutes. Calculate the population after 4 hours, assuming exponential growth. (c) In practice the population stops growing exponentially. Give two reasons.
\tcblower
\textbf{(a)} Each PCR cycle doubles the DNA, so after $n$ cycles:
\[ N = N_0 \times 2^{n} = 1 \times 2^{30} = 1\,073\,741\,824 \approx 1.07 \times 10^{9}\ \text{molecules (3 s.f.)} \]
This enormous amplification from a single molecule is why PCR is used in diagnosis (e.g.\ detecting HIV, malaria parasite \emph{Plasmodium falciparum}, or SARS-CoV-2) and forensic analysis from tiny samples.

\textbf{(b)} Number of generations in 4 hours:
\[ n = \dfrac{t}{g} = \dfrac{4 \times 60\ \mathrm{min}}{20\ \mathrm{min}} = 12\ \text{generations} \]
Final population:
\[ N = N_0 \times 2^{n} = 1.0 \times 10^{3} \times 2^{12} = 10^{3} \times 4096 = 4.096 \times 10^{6} \approx 4.10 \times 10^{6}\ \text{cells (3 s.f.)} \]

\textbf{(c)} Exponential growth stops because:
\begin{itemize}
\item \textbf{nutrients (glucose, amino acids, phosphate) become exhausted} in the fermenter;
\item \textbf{toxic waste products accumulate} (e.g.\ organic acids lowering the pH, ethanol, ammonia), and oxygen may become limiting for aerobic cultures, so cells enter the stationary phase.
\end{itemize}
In industrial fermenters, sensors continuously monitor and adjust pH, temperature, dissolved oxygen and nutrient supply (fed-batch or continuous culture) to keep the culture productive for as long as possible.

\textbf{Examiner's tip:} the most frequent errors are forgetting to convert hours to minutes and writing $N = N_0 \times 2t$ instead of $N_0 \times 2^{t/g}$. Always compute the \emph{number of generations} first.
\end{exoresolu}

\courssub{Skill 9: Evaluation and remediation -- integrated review}

\begin{aretenir}[Key points for GCE Advanced Level Biotechnology]
\begin{itemize}
\item \textbf{Traditional biotechnology:} uses wild-type micro-organisms, no genetic modification. Examples: palm wine, \emph{fufu}, \emph{burukutu}, yoghurt, bread, cheese. Fermentation = anaerobic, 2 ATP/glucose.
\item \textbf{Modern biotechnology:} involves genetic manipulation (recombinant DNA, gene editing). Tools: restriction enzymes, ligase, vectors (plasmids, viruses), PCR, electrophoresis, CRISPR-Cas9.
\item \textbf{Recombinant DNA steps:} isolate gene $\rightarrow$ cut vector/gene with same RE $\rightarrow$ ligate $\rightarrow$ transform $\rightarrow$ select (marker) $\rightarrow$ screen $\rightarrow$ scale up in fermenter $\rightarrow$ purify product.
\item \textbf{Gene therapy:} somatic (non-heritable, e.g.\ lentiviral $\beta$-globin for sickle cell) vs germ-line (heritable, embryo editing, prohibited). CRISPR: gRNA guides Cas9 to target.
\item \textbf{RFLP diagnosis:} mutation alters restriction site $\rightarrow$ fragment size change $\rightarrow$ gel band pattern reveals genotype. Sickle cell: MstII site lost in HbS.
\item \textbf{Applications:} health (insulin, vaccines, mAbs, gene therapy), agriculture (Bt crops, biofortification, disease resistance), environment (bioremediation, biofuels), industry (enzymes, bioplastics).
\item \textbf{Risks/ethics:} gene flow, resistance evolution, allergenicity, farmer dependence, germ-line heritability, equity. Regulation: biosafety laws, Cartagena Protocol.
\item \textbf{Calculations:} PCR $N = N_0 2^n$; bacterial growth $N = N_0 2^{t/g}$. Always convert time to same units as generation time.
\item \textbf{Cameroon context:} sickle cell high prevalence (balanced polymorphism with malaria), palm wine/\emph{burukutu} traditional fermentation, need for disease-resistant plantain/cassava, vitamin A deficiency in North, oil spill bioremediation in coastal regions.
\end{itemize}
\end{aretenir}

\begin{exercice}[GCE-style practice questions]
\begin{enumerate}
\item \textbf{Traditional fermentation:} Describe the changes that occur when cassava tubers are soaked in water to produce \emph{fufu}. Name the micro-organisms involved, the type of respiration, the main products, and how they contribute to detoxification and texture change. (8 marks)
\item \textbf{Data analysis:} A student investigated the effect of pH on yeast fermentation. The volume of $\ce{CO2}$ produced in 30 min at different pH values was: pH 3: 5 cm$^3$; pH 5: 38 cm$^3$; pH 7: 22 cm$^3$; pH 9: 8 cm$^3$. (a) Plot a graph. (b) State the optimum pH. (c) Explain the low rate at pH 3 and pH 9. (d) Name two controlled variables. (7 marks)
\item \textbf{Recombinant DNA:} Outline the steps to produce a transgenic maize plant expressing the Bt toxin gene. Include the roles of restriction enzyme, DNA ligase, \emph{Agrobacterium tumefaciens}, and a selectable marker. (10 marks)
\item \textbf{RFLP:} The $\beta$-globin gene is 1.35 kb. MstII cuts the normal allele (HbA) into 1.15 kb and 0.20 kb fragments. The sickle allele (HbS) lacks the site. Draw the expected gel bands for: (i) HbA HbA, (ii) HbA HbS, (iii) HbS HbS. Explain why the 0.20 kb band is often not seen. (6 marks)
\item \textbf{Gene therapy:} Distinguish between somatic and germ-line gene therapy. Explain why somatic therapy for sickle cell disease targets haematopoietic stem cells. State two risks of using lentiviral vectors. (7 marks)
\item \textbf{Calculations:} (a) How many PCR cycles are needed to obtain at least $10^9$ molecules from 10 starting molecules? (b) A culture of \emph{Bacillus thuringiensis} doubles every 30 min. If the fermenter is inoculated with $5 \times 10^4$ spores/mL, what is the concentration after 5 hours? (5 marks)
\item \textbf{Essay:} ``Modern biotechnology offers solutions to Cameroon's agricultural challenges, but its adoption carries significant risks.'' Discuss this statement with reference to at least two specific crop examples. (15 marks)
\end{enumerate}
\end{exercice}

\begin{corrige}[Exercise 1 -- \emph{Fufu} fermentation]
\textbf{Micro-organisms:} Lactic acid bacteria (\emph{Lactobacillus}, \emph{Leuconostoc}) and sometimes yeasts. \textbf{Respiration:} Anaerobic (lactic fermentation). \textbf{Equation:} $\ce{C6H12O6 -> 2 CH3CHOHCOOH}$ + 2 ATP. \textbf{Detoxification:} Cassava contains cyanogenic glycosides (linamarin, lotaustralin). Soaking allows endogenous enzyme linamarase and microbial enzymes to hydrolyse them, releasing volatile HCN which diffuses out. Lactic acid lowers pH, inhibiting putrefactive bacteria. \textbf{Texture:} Pectinases and cellulases from microbes soften cell walls; starch granules swell and ferment, giving the characteristic doughy texture.
\end{corrige}

\begin{corrige}[Exercise 2 -- pH effect]
\textbf{(a)} Graph: pH on x-axis (3.5,7.9), $\ce{CO2}$ volume on y-axis. Peak at pH 5. \textbf{(b)} Optimum pH = 5. \textbf{(c)} pH 3: excess $\ce{H+}$ denatures enzymes (zymase), disrupts tertiary structure. pH 9: excess $\ce{OH-}$ similarly denatures enzymes; also yeast membrane transport disrupted. \textbf{(d)} Temperature, yeast mass, glucose concentration, volume, time, anaerobiosis.
\end{corrige}

\begin{corrige}[Exercise 3 -- Bt maize]
1. Isolate \emph{cry} gene from \emph{Bacillus thuringiensis}. 2. Cut gene and plant expression vector (Ti plasmid derivative) with same restriction enzyme $\rightarrow$ sticky ends. 3. Ligate with DNA ligase $\rightarrow$ recombinant Ti plasmid. 4. Transform into \emph{Agrobacterium tumefaciens}. 5. Co-cultivate \emph{Agrobacterium} with maize embryogenic callus $\rightarrow$ T-DNA transfer. 6. Select transformed cells on herbicide/antibiotic medium (marker gene, e.g.\ \emph{bar} for phosphinothricin resistance). 7. Regenerate whole plants via tissue culture. 8. Screen for \emph{cry} expression (ELISA/PCR). 9. Field trials, biosafety assessment.
\end{corrige}

\begin{corrige}[Exercise 4 -- RFLP gel]
\textbf{(i) HbA HbA:} two bands at 1.15 kb and 0.20 kb (0.20 kb may be faint/absent). \textbf{(ii) HbA HbS:} three bands at 1.35 kb, 1.15 kb, 0.20 kb (or two visible: 1.35, 1.15). \textbf{(iii) HbS HbS:} single band at 1.35 kb. \textbf{Reason:} 0.20 kb fragment binds little stain, runs fast, may run off gel.
\end{corrige}

\begin{corrige}[Exercise 5 -- Gene therapy]
\textbf{Somatic:} targets body cells, not inherited, treats patient only. \textbf{Germ-line:} targets gametes/embryo, heritable, affects offspring, ethically prohibited. \textbf{Stem cells:} self-renew and differentiate into all blood lineages $\rightarrow$ lifelong supply of corrected RBCs. \textbf{Lentiviral risks:} insertional mutagenesis (oncogene activation), immune response to vector, generation of replication-competent virus.
\end{corrige}

\begin{corrige}[Exercise 6 -- Calculations]
\textbf{(a)} $10 \times 2^n \ge 10^9 \Rightarrow 2^n \ge 10^8 \Rightarrow n \ge \log_2(10^8) = 8 \times \log_2(10) \approx 8 \times 3.322 = 26.57 \Rightarrow n = 27$ cycles. \textbf{(b)} $t = 5\ \mathrm{h} = 300\ \mathrm{min}$, $g = 30\ \mathrm{min}$, $n = 300/30 = 10$ generations. $N = 5 \times 10^4 \times 2^{10} = 5 \times 10^4 \times 1024 = 5.12 \times 10^7$ spores/mL.
\end{corrige}

\begin{corrige}[Exercise 7 -- Essay outline]
\textbf{Intro:} Define modern biotechnology, relevance to Cameroon (agriculture = 70\% employment, food security challenges). \textbf{Benefits:} (1) Bt cotton/maize -- reduce pesticide use, increase yield, farmer income; (2) Disease-resistant plantain/banana (BXW, Fusarium) -- staple food security; (3) Biofortified cassava (provitamin A) -- nutrition; (4) Drought-tolerant maize -- climate adaptation in North. \textbf{Risks:} (1) Gene flow to wild relatives (wild cotton, wild banana); (2) Pest resistance evolution (refuge compliance difficult for smallholders); (3) Seed dependency/patents vs farmer-saved seeds; (4) Biodiversity loss (displacement of landraces); (5) Socio-economic: cost, access, trade implications (EU GMO regulations). \textbf{Regulation:} Cameroon Biosafety Law (2003), Cartagena Protocol, case-by-case risk assessment, labelling, public participation. \textbf{Conclusion:} Potential is real but requires strong regulatory capacity, farmer involvement, coexistence measures, and integration with agroecology. Not a silver bullet.
\end{corrige}

\begin{savaistu}[Biotechnology in Cameroon -- did you know?]
\begin{itemize}
\item The \textbf{Institute of Agricultural Research for Development (IRAD)} has developed improved plantain varieties (e.g.\ \emph{FHIA} hybrids) using conventional breeding and tissue culture, and is researching GM banana for BXW resistance.
\item \textbf{Centre Pasteur du Cameroun} and \textbf{Centre de Recherche sur les Maladies Infectieuses (CREMI)} use PCR for diagnosis of HIV, tuberculosis, malaria, and COVID-19 -- modern biotechnology in public health.
\item \textbf{Sickle cell screening:} Neonatal screening using HPLC or electrophoresis is being piloted in Yaoundé and Douala; RFLP/PCR prenatal diagnosis is available in specialised centres.
\item \textbf{Bioremediation:} After oil spills in the Rio del Rey and Kribi areas, \emph{Pseudomonas} and \emph{Bacillus} strains are tested for hydrocarbon degradation.
\item \textbf{Traditional biotechnology economy:} Palm wine, \emph{bilibili} (millet beer), \emph{fufu}, \emph{eru} fermentation support thousands of rural livelihoods -- a cultural heritage and economic asset.
\end{itemize}
\end{savaistu}

\begin{attention}[Common GCE traps in Biotechnology]
\begin{itemize}
\item Confusing \textbf{traditional} (no genetic alteration) vs \textbf{modern} (recombinant DNA, gene editing) biotechnology.
\item Forgetting that fermentation yields \textbf{2 ATP} (not 36-38) and is \textbf{anaerobic}.
\item Writing ``yeast is killed by cold'' -- cold \emph{inactivates/slows} enzymes; heat \emph{denatures} them.
\item Omitting the \textbf{enzyme/tool} at each recombinant DNA step (restriction enzyme, ligase, reverse transcriptase, Taq polymerase).
\item Misinterpreting gel bands: smallest fragment runs \textbf{furthest} (not highest).
\item Confusing \textbf{carrier} (heterozygous, healthy) with \textbf{affected} (homozygous recessive, diseased).
\item Not distinguishing \textbf{somatic} (non-heritable) vs \textbf{germ-line} (heritable) gene therapy.
\item Calculation errors: not converting time units, using $2t$ instead of $2^{t/g}$.
\item One-sided essays: must give \textbf{balanced} arguments (pros + cons) + justified conclusion.
\end{itemize}
\end{attention}

\newpage

\coursec{Exercises}

\courssub{I check my knowledge}

\begin{exercice}[MCQ: Foundations of biotechnology]
For each question, choose the single correct answer.
\begin{enumerate}
\item The production of palm wine from raffia sap in the South-West Region of Cameroon is an example of:
\begin{enumerate}
\item recombinant DNA technology
\item traditional biotechnology
\item gene therapy
\item genome editing
\end{enumerate}
\item During alcoholic fermentation by \emph{Saccharomyces cerevisiae}, the role of NAD$^+$ is to:
\begin{enumerate}
\item accept electrons during glycolysis and be regenerated in the reduction of acetaldehyde
\item serve as the final electron acceptor of the respiratory chain
\item fix carbon dioxide into organic matter
\item hydrolyse sucrose into glucose and fructose
\end{enumerate}
\item A restriction endonuclease such as EcoRI:
\begin{enumerate}
\item joins two DNA fragments by forming phosphodiester bonds
\item cuts DNA at a specific recognition sequence
\item copies DNA into messenger RNA
\item amplifies a DNA fragment in vitro
\end{enumerate}
\item In gene therapy, a vector is:
\begin{enumerate}
\item a defective gene to be replaced
\item a carrier used to deliver a functional gene into target cells
\item an enzyme that cuts the chromosome
\item a protein that inhibits transcription
\end{enumerate}
\end{enumerate}
\end{exercice}

\begin{exercice}[True or false, with justification]
State whether each statement is true or false, and justify your answer in one or two sentences.
\begin{enumerate}
\item Fermentation is an aerobic process that completely oxidises glucose to carbon dioxide and water.
\item The bacterium \emph{Agrobacterium tumefaciens} can be used as a natural vector to transfer genes into plant cells.
\item CRISPR-Cas9 genome editing always modifies every cell of the organism with absolute specificity and no off-target effects.
\item Golden Rice has been engineered to produce $\beta$-carotene, a precursor of vitamin A, in its grains.
\item Recombinant human insulin is extracted directly from the pancreas of pigs.
\end{enumerate}
\end{exercice}

\begin{exercice}[Key definitions]
Define each of the following terms precisely, in two or three sentences each, and give one concrete example for each:
\begin{enumerate}
\item Traditional biotechnology
\item Recombinant DNA
\item Transgenic organism (GMO)
\item Somatic gene therapy
\item Bioremediation
\end{enumerate}
\end{exercice}

\begin{exercice}[Direct calculation: theoretical ethanol yield]
\emph{Saccharomyces cerevisiae} ferments glucose according to the overall equation:
\[ \ce{C6H12O6 -> 2 C2H5OH + 2 CO2} \]
A brewer ferments a wort containing $100\ \mathrm{g}$ of glucose. (Molar masses: glucose $= 180\ \mathrm{g\,mol^{-1}}$; ethanol $= 46\ \mathrm{g\,mol^{-1}}$.)
\begin{enumerate}
\item Calculate the number of moles of glucose fermented.
\item Deduce the theoretical mass of ethanol produced.
\item If the wort volume is $1\ \mathrm{L}$, express the theoretical ethanol yield in $\mathrm{g\,L^{-1}}$.
\item The brewer actually recovers $40\ \mathrm{g}$ of ethanol. Calculate the percentage yield and suggest one biological reason why the real yield is lower than the theoretical yield.
\end{enumerate}
\end{exercice}

\courssub{I practise}

\begin{exercice}[Restriction mapping and gel analysis]
A circular plasmid of $5.4\ \mathrm{kb}$ carries a $1.2\ \mathrm{kb}$ insert. The plasmid has one EcoRI site and one HindIII site, located $2.0\ \mathrm{kb}$ apart on the map, both outside the insert.
\begin{enumerate}
\item Predict the number and sizes of the fragments obtained after:
\begin{enumerate}
\item single digestion with EcoRI;
\item double digestion with EcoRI and HindIII.
\end{enumerate}
\item Draw the expected banding pattern on a $1\%$ agarose gel, indicating the position of each band relative to a molecular weight marker.
\item Explain how you would verify, from this gel, that the insert is present.
\end{enumerate}
\end{exercice}

\begin{exercice}[Steps of a recombinant DNA experiment]
A Cameroonian research team wants to produce human insulin in \emph{Escherichia coli}. Rearrange the following steps in the correct chronological order and justify the role of each step:
\begin{enumerate}
\item Transform competent bacteria and select recombinant colonies on a medium containing an antibiotic.
\item Isolate the human insulin gene (or synthesise its cDNA from pancreatic mRNA).
\item Extract the plasmids and verify the insert by restriction digestion.
\item Cut the plasmid vector and the insert with the same restriction enzyme.
\item Culture the transformed bacteria in a fermenter and purify the insulin produced.
\item Ligate the insert into the plasmid using DNA ligase.
\end{enumerate}
\end{exercice}

\begin{exercice}[Fermentation monitoring in a local brewery]
During the production of a traditional sorghum beer, a student measures the sugar concentration and the ethanol concentration of the must over five days:
\begin{center}
\begin{tabularx}{\linewidth}{|l|X|X|X|X|X|}
\hline
\rowcolor{popblueL} Day & 0 & 1 & 2 & 3 & 4 \\
\hline Sugar ($\mathrm{g\,L^{-1}}$) & 120 & 95 & 60 & 25 & 8 \\
\hline Ethanol ($\mathrm{g\,L^{-1}}$) & 0 & 12 & 30 & 48 & 55 \\
\hline
\end{tabularx}
\end{center}
\begin{enumerate}
\item Plot the two curves on the same graph (time on the x-axis).
\item Describe and explain the shape of each curve in relation to yeast metabolism.
\item Explain why ethanol production slows down after day 3, giving two possible reasons.
\item State two precautions a local producer should take to avoid contamination of the brew by unwanted microorganisms.
\end{enumerate}
\end{exercice}

\begin{exercice}[Biotechnology for support and locomotion]
Tissue engineering is increasingly used to repair damaged cartilage and bone.
\begin{enumerate}
\item Explain the principle of producing a cartilage implant from a patient's own stem cells cultured on a biodegradable scaffold.
\item State two advantages of using the patient's own cells rather than a donor's cells, in relation to the immune system.
\item Name one biotechnological product (other than cell therapy) used to support bone repair, and briefly describe how it is produced.
\end{enumerate}
\end{exercice}

\courssub{I prepare for the GCE}

\begin{exercice}[CRISPR-Cas9 therapy for sickle cell disease \hfill (20 marks)]
Sickle cell disease is frequent in Cameroon. A new treatment uses CRISPR-Cas9 to switch off the \emph{BCL11A} gene in the patient's haematopoietic stem cells, which reactivates the production of fetal haemoglobin (HbF).
\begin{enumerate}
\item Describe the molecular basis of sickle cell disease, indicating the gene concerned and the type of mutation involved. \textbf{(4 marks)}
\item Explain step by step how the CRISPR-Cas9 system can be guided to the \emph{BCL11A} gene and what happens after the cut. \textbf{(6 marks)}
\item Explain why reactivating HbF compensates for the defective adult haemoglobin. \textbf{(3 marks)}
\item Identify two risks associated with this approach, including off-target effects, and propose one molecular follow-up strategy to monitor treated patients. \textbf{(4 marks)}
\item Discuss one ethical or social issue raised by the availability of such an expensive therapy in Cameroon. \textbf{(3 marks)}
\end{enumerate}
\end{exercice}

\begin{exercice}[Golden Rice versus vitamin A supplementation \hfill (20 marks)]
Vitamin A deficiency causes blindness and increases child mortality in many tropical countries. Golden Rice is a transgenic rice engineered to accumulate $\beta$-carotene in the grain.
\begin{enumerate}
\item Outline the biotechnological steps used to obtain Golden Rice, naming the type of vector commonly used for monocot or dicot transformation. \textbf{(5 marks)}
\item Compare Golden Rice and periodic vitamin A capsule supplementation under the following headings: nutritional effectiveness, agronomic constraints, economic cost for families and the State, and ethical/environmental concerns. Present your comparison in a table. \textbf{(10 marks)}
\item As an adviser to the Cameroonian Ministry of Public Health, write a short argued paragraph (8--10 lines) recommending a strategy that combines both approaches for rural areas. \textbf{(5 marks)}
\end{enumerate}
\end{exercice}

\begin{exercice}[Risk assessment for a GM biosurfactant bacterium \hfill (20 marks)]
A company proposes to clean up an oil-polluted site in the Littoral Region of Cameroon by releasing a genetically modified bacterium that produces a biosurfactant degrading petroleum hydrocarbons.
\begin{enumerate}
\item Explain the principle of bioremediation and the role of the biosurfactant in making hydrocarbons accessible to the bacteria. \textbf{(4 marks)}
\item Draw up a risk assessment sheet identifying: (i) two risks for the local ecosystems (including horizontal gene transfer), (ii) two risks for human populations living near the site, and (iii) two containment or monitoring measures for each category of risk. Present your answer as a structured table. \textbf{(10 marks)}
\item Propose an experimental protocol, carried out before any field release, to test the survival and spread of the GM bacterium in soil microcosms. \textbf{(4 marks)}
\item State one argument in favour of, and one against, authorising this release, and conclude with a justified personal position. \textbf{(2 marks)}
\end{enumerate}
\end{exercice}

\newpage

\coursec{Solutions to the exercises}

\begin{corrige}[Exercise 1 — MCQ: Foundations of biotechnology]
\begin{enumerate}
\item \textbf{Correct answer: (b) traditional biotechnology.} Palm wine production exploits the natural fermentation of raffia sap sugars by wild yeasts, an ancestral practice that does not involve any manipulation of DNA. It therefore belongs to \cle{traditional biotechnology}, unlike recombinant DNA technology, gene therapy or genome editing, which all require direct modification of genetic material.
\item \textbf{Correct answer: (a).} During glycolysis, NAD$^+$ accepts electrons (and H$^+$) from glyceraldehyde-3-phosphate and is reduced to NADH. In alcoholic fermentation, NADH is re-oxidised to NAD$^+$ when acetaldehyde is reduced to ethanol, allowing glycolysis to continue in the absence of oxygen. Oxygen, not NAD$^+$, is the final electron acceptor of the respiratory chain.
\item \textbf{Correct answer: (b).} A restriction endonuclease such as EcoRI recognises a specific palindromic sequence (5'-GAATTC-3') and cuts the DNA backbone at that site. Joining fragments is the role of DNA \emph{ligase}; copying DNA into mRNA is transcription; in vitro amplification is PCR.
\item \textbf{Correct answer: (b).} In gene therapy, a \cle{vector} (modified virus, liposome, plasmid) is a carrier used to deliver a functional copy of a gene into the patient's target cells. It is neither the defective gene, nor a cutting enzyme, nor a transcription inhibitor.
\end{enumerate}
\end{corrige}

\begin{corrige}[Exercise 2 — True or false, with justification]
\begin{enumerate}
\item \textbf{False.} Fermentation is an \emph{anaerobic} process (it occurs without oxygen) and is an \emph{incomplete} oxidation of glucose: the end products (ethanol, lactate) still contain most of the chemical energy. Complete oxidation to \ce{CO2} and \ce{H2O} is aerobic respiration.
\item \textbf{True.} \emph{Agrobacterium tumefaciens} naturally transfers a fragment of its Ti plasmid (T-DNA) into the genome of plant cells, causing crown gall disease. Biotechnologists disarm the Ti plasmid and replace the tumour genes with a gene of interest, making it a natural vector for plant genetic engineering (mainly dicots).
\item \textbf{False.} CRISPR-Cas9 is highly specific but not absolute: \cle{off-target} cuts can occur at sequences similar to the target, and editing efficiency varies between cells (mosaicism), so not every cell is necessarily modified. This is why edited cells must be screened and sequenced.
\item \textbf{True.} Golden Rice was engineered by introducing genes (e.g.\ \emph{psy} from daffodil/maize and \emph{crtI} from a bacterium) so that the endosperm accumulates $\beta$-carotene (provitamin A), which the human body converts into vitamin A.
\item \textbf{False.} Recombinant human insulin is produced by bacteria (\emph{E.~coli}) or yeast carrying the \emph{human} insulin gene inserted into a plasmid; it is not extracted from pig pancreas. Animal (porcine/bovine) insulin was the older, pre-recombinant source.
\end{enumerate}
\end{corrige}

\begin{corrige}[Exercise 3 — Key definitions]
\begin{enumerate}
\item \textbf{Traditional biotechnology:} the use of living organisms or their enzymes, based on empirical ancestral know-how, to produce or transform food and other goods, without any direct manipulation of DNA. \emph{Example:} fermentation of raffia sap into palm wine, or of corn/cassava dough into ``nkui''-type products or sorghum beer in Cameroon.
\item \textbf{Recombinant DNA:} a DNA molecule artificially constructed \emph{in vitro} by joining DNA fragments from different sources (e.g.\ a human gene inserted into a bacterial plasmid) using restriction enzymes and DNA ligase. \emph{Example:} a plasmid carrying the human insulin gene used to transform \emph{E.~coli}.
\item \textbf{Transgenic organism (GMO):} an organism whose genome has been stably modified by the introduction of one or more genes from another species (a transgene), which is expressed and transmitted to offspring. \emph{Example:} Bt maize carrying a \emph{Bacillus thuringiensis} toxin gene, or Golden Rice.
\item \textbf{Somatic gene therapy:} the introduction of a functional gene into the \emph{somatic} (body, non-reproductive) cells of a patient to correct or compensate for a genetic defect; the modification is not transmitted to descendants. \emph{Example:} delivery of a functional CFTR gene to lung cells of a cystic fibrosis patient using a viral vector.
\item \textbf{Bioremediation:} the use of living organisms (bacteria, fungi, plants) to degrade, transform or immobilise pollutants in a contaminated environment (soil, water). \emph{Example:} hydrocarbon-degrading bacteria used to clean up oil-polluted soils in the Littoral Region.
\end{enumerate}
\end{corrige}

\begin{corrige}[Exercise 4 — Direct calculation: theoretical ethanol yield]
\begin{enumerate}
\item Number of moles of glucose:
\[ n(\text{glucose}) = \frac{m}{M} = \frac{100\ \mathrm{g}}{180\ \mathrm{g\,mol^{-1}}} \approx 0.556\ \mathrm{mol}. \]
\item From the equation $\ce{C6H12O6 -> 2 C2H5OH + 2 CO2}$, one mole of glucose gives two moles of ethanol:
\[ n(\text{ethanol}) = 2 \times 0.556 = 1.111\ \mathrm{mol}, \]
\[ m(\text{ethanol}) = n \times M = 1.111 \times 46 \approx 51.1\ \mathrm{g}. \]
\item For a wort volume of $1\ \mathrm{L}$, the theoretical yield is:
\[ \frac{51.1\ \mathrm{g}}{1\ \mathrm{L}} = 51.1\ \mathrm{g\,L^{-1}}. \]
\item Percentage yield:
\[ \text{yield} = \frac{40}{51.1} \times 100 \approx 78.3\ \%. \]
\textbf{Biological reason for the lower real yield (any one):} part of the glucose is used by the yeast for its own growth and respiration (biomass, ATP production) rather than for fermentation; some ethanol evaporates or remains dissolved in the lees; fermentation may stop before all sugar is consumed because ethanol becomes toxic to the yeast.
\end{enumerate}
\end{corrige}

\begin{corrige}[Exercise 5 — Restriction mapping and gel analysis]
\begin{enumerate}
\item \textbf{Predicted fragments:}
\begin{enumerate}
\item \emph{Single digestion with EcoRI:} a circular plasmid cut at \textbf{one} site is simply linearised, giving \textbf{one linear fragment of $5.4\ \mathrm{kb}$}.
\item \emph{Double digestion EcoRI + HindIII:} two cuts on a circular molecule give \textbf{two fragments}: one of $2.0\ \mathrm{kb}$ (the distance between the two sites) and one of $5.4 - 2.0 = 3.4\ \mathrm{kb}$.
\end{enumerate}
\item \textbf{Expected banding pattern:} on a $1\%$ agarose gel, small fragments migrate faster (further towards the anode, bottom of the gel).
\begin{popfigure}
\begin{tikzpicture}[scale=0.9]
\draw[fill=popblueL,draw=popdark] (0,0) rectangle (6,4.5);
\node[anchor=west,popdark] at (0,4.8) {\small Marker \quad EcoRI \quad EcoRI+HindIII};
\foreach \y/\s in {3.9/10, 3.3/5, 2.7/3.4, 2.1/2.0, 1.5/1.2, 0.9/0.5} {
  \draw[fill=popdark] (0.3,\y) rectangle (1.5,\y+0.18);
  \node[anchor=west,popmuted,scale=0.7] at (1.55,\y+0.09) {\s\ kb};
}
\draw[fill=poporange] (2.5,2.35) rectangle (3.7,2.53);
\node[anchor=west,poporange,scale=0.7] at (3.75,2.44) {5.4 kb};
\draw[fill=poppurple] (4.6,2.7) rectangle (5.8,2.88);
\draw[fill=poppurple] (4.6,2.1) rectangle (5.8,2.28);
\node[anchor=west,poppurple,scale=0.7] at (5.85,2.79) {3.4 kb};
\node[anchor=west,poppurple,scale=0.7] at (5.85,2.19) {2.0 kb};
\node[popmuted,scale=0.8] at (3,-0.4) {migration $\rightarrow$ anode (bottom)};
\end{tikzpicture}
\legende{Simulated agarose gel: marker (left), EcoRI single digest (middle), EcoRI+HindIII double digest (right).}
\end{popfigure}
\item \textbf{Verifying the insert is present:} the undigested/linearised recombinant plasmid is $5.4\ \mathrm{kb}$, whereas the empty vector would be only $5.4 - 1.2 = 4.2\ \mathrm{kb}$; the size shift of $1.2\ \mathrm{kb}$ relative to the empty-vector control shows the insert is present. To confirm directly, one can digest with an enzyme that cuts on each side of the insert to release a $1.2\ \mathrm{kb}$ band, or use the $1.2\ \mathrm{kb}$ insert as a probe in a Southern blot.
\end{enumerate}
\end{corrige}

\begin{corrige}[Exercise 6 — Steps of a recombinant DNA experiment]
\textbf{Correct chronological order:} (b) $\rightarrow$ (d) $\rightarrow$ (f) $\rightarrow$ (a) $\rightarrow$ (c) $\rightarrow$ (e).
\begin{enumerate}
\item \textbf{(b) Isolate the human insulin gene} (or synthesise its cDNA from pancreatic mRNA using reverse transcriptase). \emph{Role:} obtain the gene of interest; cDNA is preferred because bacteria cannot remove introns.
\item \textbf{(d) Cut the plasmid vector and the insert with the same restriction enzyme.} \emph{Role:} generate compatible (complementary sticky) ends so that the insert can be joined to the vector.
\item \textbf{(f) Ligate the insert into the plasmid using DNA ligase.} \emph{Role:} DNA ligase forms phosphodiester bonds, sealing the insert into the vector to create a \cle{recombinant plasmid}.
\item \textbf{(a) Transform competent bacteria and select recombinant colonies on an antibiotic medium.} \emph{Role:} introduce the plasmid into \emph{E.~coli}; the plasmid's antibiotic-resistance marker allows only transformed bacteria to grow.
\item \textbf{(c) Extract the plasmids and verify the insert by restriction digestion.} \emph{Role:} quality control — check that the selected colonies carry the correct insert (expected fragment sizes on a gel).
\item \textbf{(e) Culture the transformed bacteria in a fermenter and purify the insulin.} \emph{Role:} large-scale production, then extraction and purification of the recombinant human insulin for medical use.
\end{enumerate}
\end{corrige}

\begin{corrige}[Exercise 7 — Fermentation monitoring in a local brewery]
\begin{enumerate}
\item \textbf{Graph:}
\begin{popfigure}
\begin{tikzpicture}
\begin{axis}[width=11cm,height=6.5cm,xlabel={Time (days)},ylabel={Concentration ($\mathrm{g\,L^{-1}}$)},
xmin=0,xmax=4,ymin=0,ymax=130,xtick={0,1,2,3,4},legend pos=north east,grid=both,grid style={popline!40}]
\addplot[color=popblue,mark=*,thick] coordinates {(0,120)(1,95)(2,60)(3,25)(4,8)};
\addlegendentry{Sugar}
\addplot[color=poporange,mark=square*,thick] coordinates {(0,0)(1,12)(2,30)(3,48)(4,55)};
\addlegendentry{Ethanol}
\end{axis}
\end{tikzpicture}
\legende{Sugar consumption and ethanol production during alcoholic fermentation.}
\end{popfigure}
\item \textbf{Shape of the curves:} the sugar concentration \emph{decreases} continuously because the yeasts consume sugars as their respiratory/fermentative substrate; the ethanol concentration \emph{increases} because ethanol is the end product of alcoholic fermentation. The two curves are roughly mirror images: ethanol is produced as sugar is consumed, according to $\ce{C6H12O6 -> 2 C2H5OH + 2 CO2}$.
\item \textbf{Slowing of ethanol production after day 3 (two reasons):}
\begin{itemize}
\item the sugar substrate is becoming scarce ($25\ \mathrm{g\,L^{-1}}$ then $8\ \mathrm{g\,L^{-1}}$), so fermentation slows down by substrate limitation;
\item ethanol accumulates and becomes \emph{toxic} to the yeasts (it inhibits their enzymes and eventually kills the cells), reducing the fermenting population.
\end{itemize}
\item \textbf{Precautions against contamination (any two):} wash and, where possible, boil/sterilise the vessels and utensils; cover the brew to keep out dust, insects (e.g.\ vinegar flies that bring acetic acid bacteria) and airborne microbes; use clean water; inoculate with a healthy starter from a previous good batch rather than relying on chance contamination.
\end{enumerate}
\end{corrige}

\begin{corrige}[Exercise 8 — Biotechnology for support and locomotion]
\begin{enumerate}
\item \textbf{Principle of the cartilage implant:} stem cells (e.g.\ mesenchymal stem cells from the patient's bone marrow) are harvested, multiplied \emph{in vitro}, then seeded onto a \cle{biodegradable scaffold} (e.g.\ collagen or polylactic acid) that reproduces the shape of the damaged cartilage. In a bioreactor, with appropriate growth factors, the cells differentiate into \emph{chondrocytes} and secrete extracellular matrix. The scaffold progressively degrades while the new tissue forms, and the construct is then implanted at the lesion site.
\item \textbf{Advantages of autologous (own) cells:}
\begin{itemize}
\item \textbf{No immune rejection:} the cells carry the patient's own HLA (self) antigens, so the immune system does not recognise them as foreign;
\item \textbf{No need for immunosuppressive treatment} and no risk of transmitting donor-borne infections (e.g.\ viral contamination of the graft).
\end{itemize}
\item \textbf{Another biotechnological product for bone repair:} recombinant \emph{Bone Morphogenetic Protein} (rhBMP-2), a growth factor that stimulates osteoblast differentiation and bone formation. It is produced by recombinant DNA technology: the human BMP gene is inserted into host cells (e.g.\ mammalian cell cultures), which synthesise the protein, later purified and applied on a collagen sponge at the fracture site.
\end{enumerate}
\end{corrige}

\begin{corrige}[Exercise 9 — CRISPR-Cas9 therapy for sickle cell disease (20 marks)]
\begin{enumerate}
\item \textbf{Molecular basis (4 marks):} sickle cell disease is caused by a \emph{point mutation} (base substitution: A $\rightarrow$ T) in the gene coding for the $\beta$-chain of haemoglobin (\emph{HBB} gene, chromosome 11). The codon GAG (glutamic acid) becomes GTG (valine) at position 6 of the $\beta$-globin chain. The abnormal haemoglobin HbS polymerises when deoxygenated, deforming red blood cells into sickle shapes that block capillaries and are destroyed prematurely (chronic anaemia, painful crises). \emph{Examiner expects:} gene named, type of mutation, amino-acid substitution, one phenotypic consequence.
\item \textbf{CRISPR-Cas9 mechanism (6 marks):}
\begin{enumerate}
\item a \emph{guide RNA} (sgRNA) is synthesised, complementary to a specific sequence of the \emph{BCL11A} gene;
\item the sgRNA associates with the Cas9 nuclease and guides it, by complementary base pairing, to the target site (next to a PAM sequence);
\item Cas9 cuts both DNA strands at that precise site;
\item the cell repairs the break by \emph{non-homologous end joining} (NHEJ), which introduces small insertions/deletions that inactivate \emph{BCL11A};
\item the edited haematopoietic stem cells are reinfused into the patient, where they produce red blood cells expressing fetal haemoglobin.
\end{enumerate}
\item \textbf{Why HbF compensates (3 marks):} fetal haemoglobin (HbF, $\alpha_2\gamma_2$) does not contain $\beta$-globin, so it is unaffected by the sickle mutation; it transports oxygen normally and dilutes/inhibits the polymerisation of HbS, preventing sickling. Since \emph{BCL11A} normally represses the $\gamma$-globin gene after birth, switching it off reactivates HbF production.
\item \textbf{Risks and follow-up (4 marks):} risks — \emph{off-target cuts} elsewhere in the genome (possible mutations, cancer risk); incomplete editing (mosaicism) or genotoxicity of the conditioning chemotherapy; long-term effects of \emph{BCL11A} inactivation unknown. \emph{Follow-up:} regular molecular monitoring by PCR and DNA sequencing of the target locus and predicted off-target sites, plus measurement of HbF levels and blood counts.
\item \textbf{Ethical/social issue (3 marks):} the therapy is extremely expensive and requires highly specialised centres, so very few Cameroonian patients could access it, raising a problem of \emph{equity of access to care}; resources might alternatively fund prevention (premarital screening, genetic counselling). Any reasoned, balanced argument is credited.
\end{enumerate}
\end{corrige}

\begin{corrige}[Exercise 10 — Golden Rice versus vitamin A supplementation (20 marks)]
\begin{enumerate}
\item \textbf{Obtaining Golden Rice (5 marks):} (i) identify and isolate the genes of the $\beta$-carotene pathway: \emph{psy} (phytoene synthase, from daffodil or maize) and \emph{crtI} (bacterial carotene desaturase); (ii) place them under an endosperm-specific promoter so $\beta$-carotene accumulates in the grain; (iii) insert them into a suitable vector — the disarmed \emph{Ti plasmid of Agrobacterium tumefaciens} (standard for dicots, and adapted strains for monocots such as rice) or use biolistics (gene gun) for monocots; (iv) transform rice cells, regenerate whole plants by tissue culture on selective medium; (v) verify insertion and expression (PCR, carotenoid assays) and select stable lines.
\item \textbf{Comparison (10 marks):}
\begin{center}
\begin{tabularx}{\linewidth}{|l|X|X|}
\hline
\rowcolor{popblueL} Criterion & Golden Rice & Vitamin A capsules \\
\hline Nutritional effectiveness & Continuous daily supply of provitamin A with the staple food; dose depends on consumption and variety & Precise, high dose but only 1--2 times per year; risk of missed children \\
\hline Agronomic constraints & Requires adapted varieties, seed distribution, farmer adoption; yields must equal local varieties & None: no change to farming practices \\
\hline Economic cost & Low recurrent cost for families once seeds available; high initial research/seed-system cost for the State & Recurrent cost of purchase, storage and distribution campaigns, often donor-dependent \\
\hline Ethical/environmental concerns & GMO acceptance, possible gene flow to wild relatives, dependence on seed suppliers & Few ethical concerns; risk of overdose rare; no environmental impact \\
\hline
\end{tabularx}
\end{center}
\item \textbf{Argued recommendation (5 marks) — model paragraph:} I recommend a \emph{combined strategy}. In the short term, capsule supplementation campaigns must continue, because they reach children immediately and require no change in farming. In parallel, the State should support field trials of biofortified varieties (Golden Rice or, more realistically for Cameroon, biofortified orange-fleshed sweet potato and vitamin-A maize/cassava) adapted to local conditions, with free seed distribution to rural farmers. Biofortification then provides a sustainable, low-cost daily supply integrated into the diet, while capsules cover emergencies and populations not yet reached. Public information campaigns are essential to ensure acceptance of biofortified foods. This progressive combination guarantees both immediate coverage and long-term food autonomy.
\end{enumerate}
\end{corrige}

\begin{corrige}[Exercise 11 — Risk assessment for a GM biosurfactant bacterium (20 marks)]
\begin{enumerate}
\item \textbf{Principle (4 marks):} \cle{bioremediation} uses microorganisms to degrade pollutants into harmless products (\ce{CO2}, \ce{H2O}, biomass). Petroleum hydrocarbons are hydrophobic and poorly soluble in water, which limits their uptake by bacteria. The \emph{biosurfactant} is an amphiphilic molecule that emulsifies the hydrocarbons, increasing their dispersion and bioavailability, so the bacteria can absorb and oxidise them efficiently.
\item \textbf{Risk assessment sheet (10 marks):}
\begin{center}
\begin{tabularx}{\linewidth}{|l|X|X|}
\hline
\rowcolor{popblueL} Category & Risks & Containment / monitoring measures \\
\hline Ecosystems & (i) Horizontal gene transfer of the biosurfactant/antibiotic-resistance gene to native soil bacteria; (ii) disruption of indigenous microbial communities and soil food chains & (i) Use a strain without antibiotic-resistance markers, with a ``suicide gene'' limiting survival; (ii) confine treatment to the polluted plot with physical barriers and monitor soil microbiota by regular sampling \\
\hline Human populations & (i) Allergenicity/toxicity of the bacterium or its metabolites for workers and residents; (ii) spread of the GMO to water sources used for drinking or fishing & (i) Toxicity tests before release and protective equipment for workers; (ii) buffer zones around wells/streams and regular microbiological testing of water \\
\hline
\end{tabularx}
\end{center}
\item \textbf{Microcosm protocol (4 marks):} fill sealed soil columns (microcosms) with non-sterile local soil, some polluted with oil and some not; inoculate half with the GM bacterium (marked, e.g.\ with a fluorescent \emph{gfp} marker) and keep uninoculated controls; incubate under field-like temperature and humidity; sample at regular intervals (days 0, 7, 14, 30, 60) and quantify the GM strain by selective plating and PCR of its marker gene; measure hydrocarbon degradation and check for gene transfer to native bacteria; all manipulations in a contained laboratory (biosafety level 2) with autoclaving of all waste.
\item \textbf{Arguments and position (2 marks):} \emph{For:} the method is cheaper and more ecological than excavation/incineration and can restore farmland and fisheries. \emph{Against:} once released, a GMO cannot be recalled, and long-term ecological effects are uncertain. \emph{Personal position (example):} authorisation only after successful microcosm and greenhouse trials, with a marker-free, survival-limited strain and a strict monitoring plan — a precautionary but not prohibitive stance.
\end{enumerate}
\end{corrige}

\newpage

\coursec{Summary sheet}

\begin{popfigure}
\begin{center}\tcbox[colback=popink,colframe=popink,coltext=white,arc=9pt,boxsep=4pt,left=6pt,right=6pt]{\bfseries\small Biotechnology}\end{center}
\vspace{-2pt}
\begin{tcbraster}[raster columns=3,raster equal height=rows,raster column skip=6pt,raster row skip=6pt,enhanced,blankest]
\begin{tcolorbox}[enhanced,colback=popgreen,colframe=popgreen,coltext=white,boxrule=0.9pt,arc=3pt,left=4pt,right=4pt,top=3pt,bottom=3pt,boxsep=1pt,halign=left]\bfseries\scriptsize Traditional Biotechnology\end{tcolorbox}
\begin{tcolorbox}[enhanced,colback=poporange,colframe=poporange,coltext=white,boxrule=0.9pt,arc=3pt,left=4pt,right=4pt,top=3pt,bottom=3pt,boxsep=1pt,halign=left]\bfseries\scriptsize Homeostasis, Support, Locomotion\end{tcolorbox}
\begin{tcolorbox}[enhanced,colback=poppurple,colframe=poppurple,coltext=white,boxrule=0.9pt,arc=3pt,left=4pt,right=4pt,top=3pt,bottom=3pt,boxsep=1pt,halign=left]\bfseries\scriptsize Gene Therapy \& Genome Editing\end{tcolorbox}
\begin{tcolorbox}[enhanced,colback=poppink,colframe=poppink,coltext=white,boxrule=0.9pt,arc=3pt,left=4pt,right=4pt,top=3pt,bottom=3pt,boxsep=1pt,halign=left]\bfseries\scriptsize Recombinant DNA Technology\end{tcolorbox}
\begin{tcolorbox}[enhanced,colback=popgold,colframe=popgold,coltext=white,boxrule=0.9pt,arc=3pt,left=4pt,right=4pt,top=3pt,bottom=3pt,boxsep=1pt,halign=left]\bfseries\scriptsize Applications: Health, Agri, Environ, Industry\end{tcolorbox}
\end{tcbraster}

\legende{Mind map of Chapter BIO07: Biotechnology}
\end{popfigure}

\begin{bilan}
\textbf{Key Definitions}
\begin{itemize}
\item \cle{Biotechnology}: Use of living organisms or their components to make or modify products for specific use.
\item \cle{Traditional biotechnology}: Ancient processes like fermentation (e.g., \emph{kwacha} production, palm wine, cassava fermentation) relying on natural microbes.
\item \cle{Recombinant DNA technology}: Artificial combination of DNA from different sources to create transgenic organisms.
\item \cle{Gene therapy}: Introduction of functional genes into patient cells to treat genetic disorders (e.g., sickle cell disease trials).
\item \cle{Genome editing}: Precise alteration of DNA sequences using tools like CRISPR‑Cas9.
\item \cle{Homeostasis in biotech}: Use of biotechnological devices (biosensors, artificial organs) to maintain internal stability.
\end{itemize}

\textbf{Laws, Formulas \& Principles to Know}
\begin{itemize}
\item \cle{Chargaff’s rules} and \cle{Watson‑Crick base pairing} underpin DNA manipulation.
\item \cle{Restriction enzyme kinetics}: $v = \frac{V_{\max}[S]}{K_m + [S]}$ (Michaelis‑Menten) for predicting digestion efficiency.
\item \cle{PCR amplification}: $N = N_0 \times 2^n$ where $n$ = number of cycles.
\item \cle{Plasmid copy number} and \cle{selection marker} principles for cloning vectors.
\item \cle{Hardy‑Weinberg equilibrium} ($p^2+2pq+q^2=1$) for population genetics in GMO risk assessment.
\end{itemize}

\textbf{Methods (Step‑by‑Step)}
\begin{itemize}
\item \textbf{Fermentation monitoring}: (1) Inoculate sterile substrate; (2) Control pH, temperature, anaerobiosis; (3) Measure CO$_2$ evolution or ethanol yield; (4) Harvest product.
\item \textbf{Recombinant plasmid construction}: (1) Choose vector \& insert; (2) Digest both with compatible restriction enzymes; (3) Ligate with DNA ligase; (4) Transform competent \emph{E. coli}; (5) Screen colonies (blue-white, PCR, sequencing).
\item \textbf{CRISPR-Cas9 editing}: (1) Design sgRNA targeting gene; (2) Clone sgRNA into expression plasmid; (3) Deliver Cas9-sgRNA complex (electroporation, viral vector); (4) Verify edit by sequencing; (5) Assess off-target effects.
\item \textbf{Biosensor calibration}: (1) Immobilise enzyme/antibody; (2) Record signal vs known analyte concentrations; (3) Plot calibration curve; (4) Use for real-time homeostasis monitoring (e.g., glucose sensor).
\end{itemize}

\textbf{Common Mistakes (Traps)}
\begin{itemize}
\item Confusing \emph{traditional} (spontaneous, mixed cultures) with \emph{modern} (defined strains, controlled bioreactors) fermentation.
\item Forgetting that restriction enzymes cut at specific palindromic sites; using incompatible ends prevents ligation.
\item Assuming PCR doubles DNA every cycle indefinitely — plateau phase occurs due to reagent depletion.
\item Overlooking ethical \& regulatory frameworks (Cameroon’s biosafety law, Cartagena Protocol) when discussing GMO release.
\item Misapplying Hardy‑Weinberg: it assumes random mating, no selection, large population — rarely true for transgenic crops in field.
\end{itemize}

\textbf{GCE Advanced Level Tips}
\begin{itemize}
\item \textbf{Definitions}: Write concise, exam‑style definitions; mention at least one Cameroonian example (e.g., \emph{foléré} fermentation, sickle cell gene therapy trials in Yaoundé).
\item \textbf{Diagrams}: Draw labelled plasmid map (origin, selectable marker, MCS, insert) and CRISPR mechanism (Cas9, sgRNA, PAM, target DNA). Use TikZ‑style sketches in answers.
\item \textbf{Calculations}: Show all steps for PCR yield, restriction fragment sizes, allele frequencies. Units must be included.
\item \textbf{Essay questions}: Structure: Introduction → Traditional vs Modern → Recombinant DNA steps → Applications (health, agriculture, environment) → Ethical/safety concerns → Conclusion.
\item \textbf{Practicals}: Know how to interpret gel electrophoresis bands, fermentation graphs (lag, log, stationary phases), and biosensor calibration curves.
\end{itemize}
\end{bilan}

\findecours

\end{document}
